% Options for packages loaded elsewhere
\PassOptionsToPackage{unicode}{hyperref}
\PassOptionsToPackage{hyphens}{url}
\documentclass[
  a4paper]{report}
\usepackage{xcolor}
\usepackage[top=3.5cm,bottom=3cm,left=3.5cm,right=2cm]{geometry}
\usepackage{amsmath,amssymb}
\setcounter{secnumdepth}{5}
\usepackage{iftex}
\ifPDFTeX
  \usepackage[T1]{fontenc}
  \usepackage[utf8]{inputenc}
  \usepackage{textcomp} % provide euro and other symbols
\else % if luatex or xetex
  \usepackage{unicode-math} % this also loads fontspec
  \defaultfontfeatures{Scale=MatchLowercase}
  \defaultfontfeatures[\rmfamily]{Ligatures=TeX,Scale=1}
\fi
\usepackage{lmodern}
\ifPDFTeX\else
  % xetex/luatex font selection
  \setmainfont[]{Times New Roman}
\fi
% Use upquote if available, for straight quotes in verbatim environments
\IfFileExists{upquote.sty}{\usepackage{upquote}}{}
\IfFileExists{microtype.sty}{% use microtype if available
  \usepackage[]{microtype}
  \UseMicrotypeSet[protrusion]{basicmath} % disable protrusion for tt fonts
}{}
\usepackage{setspace}
\usepackage{longtable,booktabs,array}
\newcounter{none} % for unnumbered tables
\usepackage{calc} % for calculating minipage widths
% Correct order of tables after \paragraph or \subparagraph
\usepackage{etoolbox}
\makeatletter
\patchcmd\longtable{\par}{\if@noskipsec\mbox{}\fi\par}{}{}
\makeatother
% Allow footnotes in longtable head/foot
\IfFileExists{footnotehyper.sty}{\usepackage{footnotehyper}}{\usepackage{footnote}}
\makesavenoteenv{longtable}
\setlength{\emergencystretch}{3em} % prevent overfull lines
\providecommand{\tightlist}{%
  \setlength{\itemsep}{0pt}\setlength{\parskip}{0pt}}
\usepackage{indentfirst}\usepackage{float}\usepackage{booktabs}\usepackage{longtable}\usepackage{fancyhdr}\pagestyle{fancy}\fancyhf{}\fancyhead[R]{\thepage}\renewcommand{\headrulewidth}{0.5pt}\usepackage{unicode-math}\setmathfont{XITS Math}
\usepackage{bookmark}
\IfFileExists{xurl.sty}{\usepackage{xurl}}{} % add URL line breaks if available
\urlstyle{same}
\hypersetup{
  pdftitle={KHÓA LUẬN TỐT NGHIỆP: Ứng dụng AI trong phát hiện và phòng chống tấn công Brute-force trên hệ thống SSH với dự đoán sớm},
  pdfauthor={Trường Đại học FPT --- Chuyên ngành An toàn thông tin},
  hidelinks,
  pdfcreator={LaTeX via pandoc}}

\title{KHÓA LUẬN TỐT NGHIỆP: Ứng dụng AI trong phát hiện và phòng chống
tấn công Brute-force trên hệ thống SSH với dự đoán sớm}
\author{Trường Đại học FPT --- Chuyên ngành An toàn thông tin}
\date{Hà Nội, 2026}

\begin{document}
\maketitle

{
\setcounter{tocdepth}{2}
\tableofcontents
}
\setstretch{1.5}
\newpage

\chapter*{LỜI CẢM ƠN}\label{lux1eddi-cux1ea3m-ux1a1n}
\addcontentsline{toc}{chapter}{LỜI CẢM ƠN}

Lời đầu tiên, nhóm chúng em xin gửi lời cảm ơn chân thành và sâu sắc
nhất đến giảng viên hướng dẫn đã tận tình chỉ bảo, định hướng và hỗ trợ
nhóm trong suốt quá trình thực hiện khóa luận tốt nghiệp.

Chúng em xin cảm ơn Ban Giám hiệu Trường Đại học FPT, Khoa Công nghệ
Thông tin và các thầy cô trong chuyên ngành An toàn thông tin đã tạo
điều kiện thuận lợi cho chúng em học tập và nghiên cứu.

Chúng em cũng xin gửi lời cảm ơn đến gia đình, bạn bè và các thành viên
trong nhóm đã luôn động viên, hỗ trợ và đóng góp ý kiến quý báu trong
quá trình thực hiện đề tài.

Mặc dù nhóm đã nỗ lực hết sức, khóa luận không tránh khỏi những thiếu
sót. Nhóm rất mong nhận được sự góp ý từ quý thầy cô và Hội đồng để hoàn
thiện hơn.

\begin{flushright}
\textit{Hà Nội, tháng 4 năm 2026}\\
\textit{Nhóm sinh viên thực hiện}
\end{flushright}

\newpage

\chapter*{TÓM TẮT}\label{tuxf3m-tux1eaft}
\addcontentsline{toc}{chapter}{TÓM TẮT}

Tấn công brute-force SSH là một trong những mối đe dọa phổ biến và dai
dẳng nhất trong lĩnh vực an ninh mạng. Các công cụ phòng chống truyền
thống như Fail2Ban hoạt động dựa trên ngưỡng tĩnh, chỉ phản ứng sau khi
tấn công đã xảy ra và dễ bị lẩn tránh bởi các kỹ thuật tấn công tiên
tiến như low-and-slow hoặc distributed brute-force.

Khóa luận này trình bày việc thiết kế, triển khai và đánh giá một hệ
thống phát hiện và phòng chống tấn công brute-force SSH sử dụng trí tuệ
nhân tạo, với khả năng dự đoán sớm (early prediction). Hệ thống sử dụng
thuật toán Isolation Forest làm mô hình phát hiện bất thường chính, kết
hợp với cơ chế ngưỡng động EWMA-Adaptive Percentile để cung cấp khả năng
cảnh báo sớm hai mức (EARLY\_WARNING và ALERT).

Bộ dữ liệu nghiên cứu gồm 174.250 dòng nhật ký SSH, bao gồm dữ liệu tấn
công thực tế từ hệ thống honeypot (119.729 dòng, 679 địa chỉ IP) và dữ
liệu hành vi bình thường (54.521 dòng, 64 người dùng). Hệ thống trích
xuất 14 đặc trưng hành vi theo cửa sổ thời gian 5 phút cho mỗi địa chỉ
IP nguồn.

Kết quả: Isolation Forest đạt F1-Score 93,74\%, Accuracy 90,31\%, FPR
29,00\% sau tối ưu. So sánh với LOF (F1=89,94\%) và One-Class SVM
(F1=94,55\%), IF được chọn nhờ khả năng tạo anomaly score phù hợp cho
thuật toán ngưỡng động. Hệ thống được triển khai hoàn chỉnh với Docker
Compose (9 services), bao gồm FastAPI backend, React dashboard, ELK
Stack, Fail2Ban và pipeline phát hiện thời gian thực.

\textbf{Từ khóa:} SSH, brute-force, phát hiện bất thường, Isolation
Forest, ngưỡng động, EWMA, dự đoán sớm, học máy, ELK Stack, Fail2Ban.

\newpage

\chapter*{ABSTRACT}\label{abstract}
\addcontentsline{toc}{chapter}{ABSTRACT}

SSH brute-force attacks remain one of the most prevalent and persistent
threats in cybersecurity. Traditional defense tools such as Fail2Ban
operate on static thresholds, reacting only after an attack has occurred
and being easily evaded by advanced techniques such as low-and-slow or
distributed brute-force attacks.

This thesis presents the design, implementation, and evaluation of an
AI-powered system for detecting and preventing SSH brute-force attacks
with early prediction capability. The system employs Isolation Forest as
the primary anomaly detection model, combined with an EWMA-Adaptive
Percentile dynamic threshold mechanism to provide two-level early
warnings.

The research dataset comprises 174,250 SSH log entries from a honeypot
system (119,729 lines, 679 attacking IPs) and simulated normal behavior
(54,521 lines, 64 users). The system extracts 14 behavioral features per
5-minute window per source IP.

Results: Isolation Forest achieves F1-Score of 93.74\%, Accuracy of
90.31\%, and FPR of 29.00\% after optimization. The complete system is
deployed as a Docker Compose stack with 9 services including real-time
detection, FastAPI backend, React dashboard, ELK Stack visualization,
and Fail2Ban integration.

\textbf{Keywords:} SSH, brute-force, anomaly detection, Isolation
Forest, dynamic threshold, EWMA, early prediction, machine learning, ELK
Stack, Fail2Ban.

\newpage

\chapter{CHƯƠNG 1: GIỚI
THIỆU}\label{chux1b0ux1a1ng-1-giux1edbi-thiux1ec7u}

\section{1.1 Bối cảnh nghiên
cứu}\label{bux1ed1i-cux1ea3nh-nghiuxean-cux1ee9u}

Trong bối cảnh chuyển đổi số diễn ra mạnh mẽ trên toàn cầu, hạ tầng công
nghệ thông tin ngày càng đóng vai trò then chốt trong hoạt động của các
tổ chức, doanh nghiệp và cơ quan nhà nước. Theo báo cáo của
Cybersecurity Ventures, thiệt hại do tội phạm mạng gây ra trên toàn thế
giới dự kiến đạt mức 10,5 nghìn tỷ USD vào năm 2025, tăng từ 3 nghìn tỷ
USD vào năm 2015 {[}1{]}. Sự gia tăng nhanh chóng này phản ánh mức độ
phức tạp và quy mô ngày càng lớn của các cuộc tấn công mạng, đặt ra yêu
cầu cấp thiết về các giải pháp an ninh mạng tiên tiến và hiệu quả hơn.

Giao thức SSH (Secure Shell) là một trong những giao thức được sử dụng
phổ biến nhất để quản trị hệ thống từ xa, truyền tải tệp tin an toàn và
thiết lập các kênh liên lạc được mã hóa. Ra đời từ năm 1995 bởi Tatu
Ylönen tại Đại học Công nghệ Helsinki, SSH đã trở thành tiêu chuẩn công
nghiệp cho việc quản trị máy chủ Linux/Unix {[}2{]}. Tuy nhiên, chính sự
phổ biến này cũng khiến SSH trở thành mục tiêu hàng đầu của các cuộc tấn
công mạng, đặc biệt là tấn công brute-force --- hình thức tấn công trong
đó kẻ tấn công thử hàng nghìn đến hàng triệu tổ hợp tên đăng nhập và mật
khẩu nhằm chiếm quyền truy cập hệ thống.

Theo dữ liệu từ SANS Internet Storm Center, cổng TCP 22 (SSH) luôn nằm
trong top 5 cổng bị quét và tấn công nhiều nhất trên Internet {[}3{]}.
Báo cáo của Rapid7 năm 2023 chỉ ra rằng trung bình một máy chủ SSH công
khai nhận được hơn 10.000 lần thử đăng nhập bất hợp pháp mỗi ngày
{[}4{]}. Tại Việt Nam, theo Trung tâm Giám sát an toàn không gian mạng
quốc gia (NCSC), các cuộc tấn công brute-force SSH chiếm tỷ lệ đáng kể
trong tổng số sự cố an ninh mạng được ghi nhận hàng năm, với xu hướng
tăng liên tục từ năm 2020 đến nay {[}5{]}.

Các phương pháp phòng chống tấn công brute-force SSH truyền thống bao
gồm: giới hạn số lần đăng nhập thất bại (rate limiting), sử dụng danh
sách đen IP (IP blacklisting), xác thực bằng khóa công khai (public key
authentication), và triển khai các công cụ như Fail2Ban hoặc DenyHosts
{[}6{]}. Mặc dù các giải pháp này đã cho thấy hiệu quả nhất định, chúng
tồn tại nhiều hạn chế rõ rệt. Thứ nhất, các phương pháp dựa trên ngưỡng
tĩnh (static threshold) dễ bị kẻ tấn công lẩn tránh bằng cách giảm tốc
độ tấn công (slow brute-force) hoặc phân tán tấn công từ nhiều địa chỉ
IP khác nhau (distributed brute-force) {[}7{]}. Thứ hai, việc thiếu khả
năng học hỏi và thích ứng khiến các hệ thống truyền thống không thể nhận
diện các biến thể tấn công mới. Thứ ba, tỷ lệ cảnh báo sai (false
positive) cao có thể gây gián đoạn dịch vụ cho người dùng hợp pháp.

Trong bối cảnh đó, trí tuệ nhân tạo (AI) và học máy (Machine Learning)
nổi lên như một hướng tiếp cận đầy triển vọng cho bài toán phát hiện và
phòng chống tấn công mạng. Các mô hình học máy có khả năng phân tích
lượng lớn dữ liệu nhật ký (log), nhận diện các mẫu hành vi bất thường
(anomaly patterns) và đưa ra cảnh báo theo thời gian thực --- những khả
năng mà các phương pháp dựa trên quy tắc tĩnh (rule-based) khó có thể
đạt được {[}8{]}. Đặc biệt, các thuật toán phát hiện bất thường không
giám sát (unsupervised anomaly detection) như Isolation Forest, Local
Outlier Factor (LOF) và One-Class SVM mang lại ưu thế vượt trội trong
việc phát hiện các hình thức tấn công chưa biết trước, vì chúng không
cần dữ liệu gán nhãn (labeled data) để huấn luyện {[}9{]}.

Nghiên cứu này được thực hiện xuất phát từ nhu cầu thực tiễn trong việc
xây dựng một hệ thống phát hiện và phòng chống tấn công brute-force SSH
thông minh, có khả năng dự đoán sớm (early prediction) các cuộc tấn công
trước khi chúng gây ra thiệt hại. Hệ thống đề xuất kết hợp thuật toán
Isolation Forest với cơ chế ngưỡng động EWMA-Adaptive Percentile, tích
hợp trên nền tảng ELK Stack (Elasticsearch, Logstash, Kibana) và
Fail2Ban, nhằm tạo ra một giải pháp toàn diện từ thu thập dữ liệu, phân
tích, phát hiện đến phản ứng tự động trước các cuộc tấn công.

\section{1.2 Phát biểu vấn
đề}\label{phuxe1t-biux1ec3u-vux1ea5n-ux111ux1ec1}

Tấn công brute-force SSH là một trong những hình thức tấn công dai dẳng
và phổ biến nhất trong lĩnh vực an ninh mạng. Mặc dù đã có nhiều giải
pháp được triển khai, bài toán phát hiện và phòng chống tấn công này vẫn
tồn tại nhiều thách thức chưa được giải quyết triệt để.

\textbf{Thách thức thứ nhất: Sự tiến hóa của kỹ thuật tấn công.} Các kẻ
tấn công hiện đại không còn sử dụng phương pháp brute-force đơn giản với
tốc độ cao. Thay vào đó, họ áp dụng nhiều kỹ thuật lẩn tránh tinh vi bao
gồm: tấn công chậm (slow brute-force) với khoảng cách giữa các lần thử
dài để tránh bị phát hiện bởi cơ chế giới hạn tốc độ; tấn công phân tán
(distributed brute-force) sử dụng mạng botnet với hàng nghìn địa chỉ IP;
tấn công từ điển thông minh (intelligent dictionary attack) với danh
sách mật khẩu được tùy chỉnh theo mục tiêu; và tấn công credential
stuffing sử dụng thông tin đăng nhập bị rò rỉ từ các vụ vi phạm dữ liệu
trước đó {[}10{]}.

\textbf{Thách thức thứ hai: Hạn chế của phương pháp ngưỡng tĩnh.} Các
công cụ truyền thống như Fail2Ban hoạt động dựa trên ngưỡng cố định ---
ví dụ, chặn IP sau 5 lần đăng nhập thất bại trong 10 phút. Cách tiếp cận
này tạo ra hai vấn đề đối lập: nếu ngưỡng quá thấp, người dùng hợp pháp
quên mật khẩu có thể bị chặn nhầm (false positive); nếu ngưỡng quá cao,
kẻ tấn công có thể tiến hành tấn công chậm mà không bị phát hiện (false
negative) {[}11{]}. Hơn nữa, ngưỡng tĩnh không thể thích ứng với sự biến
động tự nhiên của lưu lượng mạng theo thời gian --- giờ cao điểm và
ngoài giờ làm việc có mẫu hành vi hoàn toàn khác nhau.

\textbf{Thách thức thứ ba: Thiếu khả năng dự đoán sớm.} Hầu hết các giải
pháp hiện tại chỉ phản ứng (reactive) sau khi tấn công đã xảy ra, thay
vì chủ động dự đoán (proactive) trước khi cuộc tấn công leo thang. Giai
đoạn trinh sát (reconnaissance) và thăm dò ban đầu (initial probing) ---
khi kẻ tấn công thử nghiệm một số ít tổ hợp để đánh giá mục tiêu ---
thường không được nhận diện, bỏ lỡ cơ hội can thiệp sớm trước khi tấn
công toàn diện diễn ra {[}12{]}.

\textbf{Thách thức thứ tư: Tích hợp và tự động hóa.} Nhiều nghiên cứu về
ứng dụng AI trong phát hiện tấn công tập trung vào khía cạnh thuật toán
mà chưa giải quyết bài toán tích hợp vào hạ tầng vận hành thực tế.
Khoảng cách giữa mô hình nghiên cứu và hệ thống triển khai
(research-to-deployment gap) vẫn là rào cản lớn trong việc áp dụng AI
vào thực tiễn an ninh mạng {[}13{]}.

Từ những thách thức trên, câu hỏi nghiên cứu chính của luận văn được
phát biểu như sau:

\begin{quote}
\emph{Làm thế nào để xây dựng một hệ thống phát hiện và phòng chống tấn
công brute-force SSH sử dụng trí tuệ nhân tạo, có khả năng dự đoán sớm
các cuộc tấn công, thích ứng động với sự thay đổi của mẫu lưu lượng, và
tích hợp hoàn chỉnh vào hạ tầng giám sát an ninh hiện đại?}
\end{quote}

Các câu hỏi nghiên cứu phụ bao gồm:

\begin{enumerate}
\def\labelenumi{\arabic{enumi}.}
\tightlist
\item
  Thuật toán Isolation Forest hoạt động hiệu quả như thế nào trong việc
  phát hiện tấn công brute-force SSH so với các thuật toán phát hiện bất
  thường khác (LOF, One-Class SVM)?
\item
  Cơ chế ngưỡng động EWMA-Adaptive Percentile cải thiện khả năng phát
  hiện ra sao so với ngưỡng tĩnh truyền thống?
\item
  Hệ thống đề xuất đạt được mức độ dự đoán sớm (early prediction) như
  thế nào đối với các kịch bản tấn công brute-force khác nhau?
\item
  Kiến trúc tích hợp ELK Stack, Isolation Forest và Fail2Ban có thể đáp
  ứng yêu cầu giám sát và phản ứng tự động trong môi trường thực tế
  không?
\end{enumerate}

\section{1.3 Mục tiêu nghiên
cứu}\label{mux1ee5c-tiuxeau-nghiuxean-cux1ee9u}

\subsection{1.3.1 Mục tiêu tổng
quát}\label{mux1ee5c-tiuxeau-tux1ed5ng-quuxe1t}

Nghiên cứu này nhằm thiết kế, xây dựng và đánh giá một hệ thống phát
hiện và phòng chống tấn công brute-force SSH dựa trên trí tuệ nhân tạo,
tích hợp khả năng dự đoán sớm và phản ứng tự động, hoạt động trên nền
tảng ELK Stack với cơ chế ngưỡng động thích ứng.

\subsection{1.3.2 Mục tiêu cụ
thể}\label{mux1ee5c-tiuxeau-cux1ee5-thux1ec3}

\textbf{Mục tiêu 1: Xây dựng bộ đặc trưng hành vi SSH toàn diện.} Thiết
kế và triển khai bộ 14 đặc trưng (features) được trích xuất từ dữ liệu
nhật ký SSH theo cửa sổ thời gian 5 phút cho mỗi địa chỉ IP, bao gồm các
đặc trưng về tần suất, phân bố thời gian, mẫu xác thực và hành vi kết
nối. Bộ đặc trưng này cần phản ánh đầy đủ các khía cạnh hành vi để phân
biệt giữa hoạt động bình thường và tấn công.

\textbf{Mục tiêu 2: Huấn luyện và đánh giá mô hình phát hiện bất
thường.} Triển khai và so sánh hiệu năng của ba thuật toán phát hiện bất
thường không giám sát: Isolation Forest, Local Outlier Factor (LOF) và
One-Class SVM. Mục tiêu đặt ra là đạt được F1-score tối thiểu 85\% và
Recall tối thiểu 95\% để đảm bảo khả năng phát hiện tấn công cao mà
không bỏ sót.

\textbf{Mục tiêu 3: Phát triển cơ chế ngưỡng động thích ứng.} Thiết kế
và triển khai phương pháp ngưỡng động kết hợp EWMA (Exponentially
Weighted Moving Average) và Adaptive Percentile, cho phép hệ thống tự
động điều chỉnh ngưỡng phát hiện theo sự biến động của lưu lượng mạng,
giảm thiểu tỷ lệ cảnh báo sai trong khi duy trì khả năng phát hiện cao.

\textbf{Mục tiêu 4: Xây dựng hệ thống tích hợp hoàn chỉnh.} Thiết kế
kiến trúc hệ thống tích hợp ELK Stack (Elasticsearch, Logstash, Kibana)
cho thu thập và trực quan hóa dữ liệu, mô hình Isolation Forest cho phát
hiện bất thường, và Fail2Ban cho phản ứng tự động. Toàn bộ hệ thống được
đóng gói bằng Docker để đảm bảo khả năng triển khai và tái sử dụng.

\textbf{Mục tiêu 5: Đánh giá hệ thống với các kịch bản tấn công thực
tế.} Xây dựng 5 kịch bản mô phỏng tấn công brute-force SSH với các đặc
điểm khác nhau (tấn công nhanh, tấn công chậm, tấn công phân tán, tấn
công từ điển, tấn công credential stuffing) để đánh giá toàn diện khả
năng phát hiện và dự đoán sớm của hệ thống.

\section{1.4 Ý nghĩa nghiên
cứu}\label{uxfd-nghux129a-nghiuxean-cux1ee9u}

\subsection{1.4.1 Ý nghĩa khoa học}\label{uxfd-nghux129a-khoa-hux1ecdc}

Nghiên cứu này đóng góp vào lĩnh vực an ninh mạng và trí tuệ nhân tạo
trên nhiều phương diện. Trước hết, luận văn cung cấp một phân tích so
sánh có hệ thống giữa ba thuật toán phát hiện bất thường không giám sát
(Isolation Forest, LOF, One-Class SVM) trong bối cảnh cụ thể là phát
hiện tấn công brute-force SSH. Mặc dù các thuật toán này đã được nghiên
cứu rộng rãi trong phát hiện xâm nhập mạng nói chung, việc đánh giá
chúng trên dữ liệu tấn công SSH thực tế từ honeypot vẫn còn hạn chế
trong tài liệu hiện có.

Thứ hai, nghiên cứu đề xuất phương pháp ngưỡng động kết hợp
EWMA-Adaptive Percentile --- một cách tiếp cận mới so với các phương
pháp ngưỡng cố định hoặc ngưỡng động đơn giản thường được sử dụng. Cơ
chế này cho phép hệ thống thích ứng với sự biến động tự nhiên của lưu
lượng mạng, đồng thời duy trì độ nhạy cao với các mẫu tấn công.

Thứ ba, bộ 14 đặc trưng hành vi SSH được thiết kế trong nghiên cứu này
có thể phục vụ làm nền tảng tham khảo cho các nghiên cứu tương lai về
phát hiện bất thường trên giao thức SSH, đặc biệt là khả năng biểu diễn
sự khác biệt giữa hành vi bình thường và tấn công trong cửa sổ thời gian
ngắn (5 phút).

\subsection{1.4.2 Ý nghĩa thực
tiễn}\label{uxfd-nghux129a-thux1ef1c-tiux1ec5n}

Về mặt thực tiễn, hệ thống được phát triển trong nghiên cứu này có thể
triển khai trực tiếp vào môi trường vận hành thực tế của các tổ chức,
doanh nghiệp tại Việt Nam. Kiến trúc dựa trên Docker đảm bảo tính di
động (portability) và khả năng triển khai nhanh chóng. Việc sử dụng ELK
Stack --- một bộ công cụ mã nguồn mở phổ biến --- giúp giảm chi phí
triển khai so với các giải pháp thương mại.

Hệ thống cung cấp khả năng giám sát trực quan qua Kibana dashboard, cho
phép quản trị viên an ninh mạng theo dõi trạng thái bảo mật SSH theo
thời gian thực. Cơ chế tích hợp với Fail2Ban đảm bảo phản ứng tự động
khi phát hiện tấn công, giảm thiểu thời gian từ khi phát hiện đến khi xử
lý (mean time to respond --- MTTR).

Đặc biệt, khả năng dự đoán sớm của hệ thống mang lại giá trị thiết thực
trong việc ngăn chặn các cuộc tấn công trước khi chúng gây ra thiệt hại.
Thay vì chờ đợi kẻ tấn công hoàn thành hàng nghìn lần thử, hệ thống có
thể nhận diện ý đồ tấn công từ giai đoạn trinh sát ban đầu, tạo cơ hội
can thiệp kịp thời.

Kết quả nghiên cứu cũng phục vụ mục đích đào tạo và nâng cao nhận thức
an ninh mạng tại các cơ sở giáo dục, đặc biệt trong chương trình đào tạo
ngành An toàn thông tin tại Trường Đại học FPT, nơi sinh viên có thể
tham khảo và mở rộng nghiên cứu.

\section{1.5 Phạm vi và giới
hạn}\label{phux1ea1m-vi-vuxe0-giux1edbi-hux1ea1n}

\subsection{1.5.1 Phạm vi nghiên
cứu}\label{phux1ea1m-vi-nghiuxean-cux1ee9u}

\textbf{Về giao thức và loại tấn công:} Nghiên cứu tập trung vào giao
thức SSH phiên bản 2 (SSH-2) và các hình thức tấn công brute-force nhắm
vào cơ chế xác thực mật khẩu (password authentication). Các hình thức
tấn công được xem xét bao gồm: brute-force cổ điển (thử tất cả tổ hợp),
tấn công từ điển (dictionary attack), tấn công chậm (slow brute-force),
tấn công phân tán (distributed attack), và tấn công credential stuffing.

\textbf{Về dữ liệu:} Nghiên cứu sử dụng hai nguồn dữ liệu chính: (1) Dữ
liệu tấn công thực tế thu thập từ hệ thống honeypot SSH, bao gồm 119.729
dòng nhật ký ghi nhận các cuộc tấn công brute-force từ nhiều nguồn trên
Internet; và (2) Dữ liệu hành vi bình thường được mô phỏng, bao gồm
54.521 dòng nhật ký đại diện cho các hoạt động SSH hợp pháp. Tổng cộng
bộ dữ liệu gồm 174.250 dòng nhật ký.

\textbf{Về thuật toán:} Nghiên cứu triển khai và đánh giá ba thuật toán
phát hiện bất thường không giám sát: Isolation Forest (thuật toán
chính), Local Outlier Factor (LOF), và One-Class SVM (các thuật toán đối
sánh). Các thuật toán học có giám sát (supervised learning) không thuộc
phạm vi nghiên cứu do đặc thù của bài toán yêu cầu khả năng phát hiện
các hình thức tấn công chưa biết trước.

\textbf{Về hạ tầng:} Hệ thống được xây dựng trên nền tảng ELK Stack
(Elasticsearch 8.x, Logstash, Kibana) kết hợp Fail2Ban, được đóng gói
bằng Docker Compose. Môi trường thử nghiệm sử dụng máy chủ chạy hệ điều
hành Linux.

\subsection{1.5.2 Giới hạn nghiên
cứu}\label{giux1edbi-hux1ea1n-nghiuxean-cux1ee9u}

Nghiên cứu có một số giới hạn cần được thừa nhận. Thứ nhất, dữ liệu hành
vi bình thường được tạo bằng mô phỏng, chưa hoàn toàn phản ánh đầy đủ sự
đa dạng của hành vi SSH hợp pháp trong mọi môi trường vận hành. Mặc dù
kịch bản mô phỏng đã được thiết kế để bao quát nhiều tình huống khác
nhau, vẫn có thể tồn tại các mẫu hành vi hợp pháp đặc thù mà mô hình
chưa gặp.

Thứ hai, nghiên cứu tập trung vào tấn công brute-force SSH và không bao
gồm các loại tấn công SSH khác như man-in-the-middle, session hijacking,
hay khai thác lỗ hổng phần mềm SSH. Việc mở rộng phạm vi để xử lý các
loại tấn công này đòi hỏi các phương pháp phân tích bổ sung.

Thứ ba, hiệu năng của hệ thống được đánh giá trong môi trường thử nghiệm
với quy mô vừa phải. Khả năng mở rộng (scalability) để xử lý lưu lượng
từ hàng trăm hoặc hàng nghìn máy chủ SSH đồng thời chưa được kiểm chứng
đầy đủ, mặc dù kiến trúc ELK Stack về mặt lý thuyết hỗ trợ khả năng mở
rộng theo chiều ngang (horizontal scaling).

Thứ tư, mô hình Isolation Forest được huấn luyện và đánh giá trên dữ
liệu cụ thể; hiệu năng có thể thay đổi khi áp dụng vào các môi trường
khác nhau với các mẫu hành vi và cấu hình SSH khác biệt. Nghiên cứu
khuyến nghị việc tái huấn luyện mô hình khi triển khai vào môi trường
mới.

\section{1.6 Cấu trúc luận
văn}\label{cux1ea5u-truxfac-luux1eadn-vux103n}

Luận văn được tổ chức thành 6 chương, mỗi chương phục vụ một mục đích cụ
thể trong việc trình bày toàn bộ quá trình nghiên cứu:

\textbf{Chương 1: Giới thiệu.} Chương này trình bày bối cảnh nghiên cứu,
phát biểu vấn đề, mục tiêu nghiên cứu, ý nghĩa khoa học và thực tiễn,
cũng như phạm vi và giới hạn của nghiên cứu. Chương cung cấp cái nhìn
tổng quan về động cơ và hướng tiếp cận của nghiên cứu.

\textbf{Chương 2: Tổng quan tài liệu.} Chương này tổng hợp và phân tích
các nền tảng lý thuyết và công trình nghiên cứu liên quan, bao gồm: giao
thức SSH, các phương pháp tấn công brute-force, học máy trong phát hiện
xâm nhập, các thuật toán phát hiện bất thường (Isolation Forest, LOF,
One-Class SVM), phương pháp ngưỡng động, và ELK Stack. Chương cũng xác
định khoảng trống nghiên cứu mà luận văn nhắm đến.

\textbf{Chương 3: Phương pháp nghiên cứu.} Chương này mô tả chi tiết
phương pháp luận nghiên cứu, bao gồm kiến trúc hệ thống, quy trình thu
thập và xử lý dữ liệu, thiết kế bộ đặc trưng, cấu hình và huấn luyện mô
hình, cơ chế ngưỡng động EWMA-Adaptive Percentile, và phương pháp đánh
giá hiệu năng.

\textbf{Chương 4: Kết quả thực nghiệm.} Chương này trình bày các kết quả
thực nghiệm, bao gồm phân tích thống kê bộ dữ liệu, kết quả huấn luyện
và đánh giá mô hình, so sánh hiệu năng giữa các thuật toán, đánh giá cơ
chế ngưỡng động, và kết quả các kịch bản mô phỏng tấn công.

\textbf{Chương 5: Thảo luận.} Chương này phân tích và diễn giải các kết
quả, thảo luận về ý nghĩa, so sánh với các công trình liên quan, và
trình bày các phát hiện chính của nghiên cứu.

\textbf{Chương 6: Kết luận và hướng phát triển.} Chương cuối tổng kết
các đóng góp chính của nghiên cứu, đánh giá mức độ đạt được các mục tiêu
đề ra, và đề xuất các hướng nghiên cứu tiếp theo.

\begin{center}\rule{0.5\linewidth}{0.5pt}\end{center}

\section{Tài liệu tham khảo Chương
1}\label{tuxe0i-liux1ec7u-tham-khux1ea3o-chux1b0ux1a1ng-1}

{[}1{]} S. Morgan, ``Cybercrime to cost the world \$10.5 trillion
annually by 2025,'' \emph{Cybersecurity Ventures}, 2021.

{[}2{]} T. Ylönen and C. Lonvick, ``The Secure Shell (SSH) Protocol
Architecture,'' RFC 4251, \emph{Internet Engineering Task Force (IETF)},
2006.

{[}3{]} SANS Internet Storm Center, ``DShield: Top 10 Target Ports,''
https://isc.sans.edu/top10.html, truy cập 2025.

{[}4{]} Rapid7, ``2023 Attack Intelligence Report,'' \emph{Rapid7
Research}, 2023.

{[}5{]} Trung tâm Giám sát an toàn không gian mạng quốc gia (NCSC),
``Báo cáo tổng kết an toàn thông tin mạng Việt Nam,'' \emph{Bộ Thông tin
và Truyền thông}, 2023.

{[}6{]} D. R. Tsai, A. Y. Chang, and S. H. Wang, ``A study of SSH brute
force attack defense,'' \emph{Journal of Information Security and
Applications}, vol.~49, pp.~102--113, 2019.

{[}7{]} M. Najafabadi, T. Khoshgoftaar, C. Calvert, and C. Kemp,
``Detection of SSH brute force attacks using aggregated netflow data,''
in \emph{Proc. IEEE 14th International Conference on Machine Learning
and Applications}, 2015, pp.~283--288.

{[}8{]} A. L. Buczak and E. Guven, ``A survey of data mining and machine
learning methods for cyber security intrusion detection,'' \emph{IEEE
Communications Surveys \& Tutorials}, vol.~18, no. 2, pp.~1153--1176,
2016.

{[}9{]} M. A. Pimentel, D. A. Clifton, L. Clifton, and L. Tarassenko,
``A review of novelty detection,'' \emph{Signal Processing}, vol.~99,
pp.~215--249, 2014.

{[}10{]} A. Simoiu, C. Gates, J. Bonneau, and S. Goel, ``I was told to
buy a software or lose my computer. I ignored it: A study of
ransomware,'' in \emph{Proc. Symposium on Usable Privacy and Security
(SOUPS)}, 2019, pp.~155--174.

{[}11{]} J. Jang-Jaccard and S. Nepal, ``A survey of emerging threats in
cybersecurity,'' \emph{Journal of Computer and System Sciences},
vol.~80, no. 5, pp.~973--993, 2014.

{[}12{]} F. Syed, M. Bashir, and A. Sharaff, ``Machine learning
approaches for intrusion detection in IoT: A comprehensive survey,''
\emph{Journal of King Saud University -- Computer and Information
Sciences}, vol.~34, no. 10, pp.~9656--9688, 2022.

{[}13{]} R. Sommer and V. Paxson, ``Outside the closed world: On using
machine learning for network intrusion detection,'' in \emph{Proc. IEEE
Symposium on Security and Privacy}, 2010, pp.~305--316.

\chapter{CHƯƠNG 2: TỔNG QUAN TÀI
LIỆU}\label{chux1b0ux1a1ng-2-tux1ed5ng-quan-tuxe0i-liux1ec7u}

\section{2.1 Giao thức SSH và cơ chế xác
thực}\label{giao-thux1ee9c-ssh-vuxe0-cux1a1-chux1ebf-xuxe1c-thux1ef1c}

\subsection{2.1.1 Tổng quan về giao thức
SSH}\label{tux1ed5ng-quan-vux1ec1-giao-thux1ee9c-ssh}

Giao thức SSH (Secure Shell) là một giao thức mạng mật mã được thiết kế
để cung cấp kênh liên lạc an toàn trên môi trường mạng không đáng tin
cậy. SSH được phát triển lần đầu vào năm 1995 bởi Tatu Ylönen, nhà
nghiên cứu tại Đại học Công nghệ Helsinki (Phần Lan), như một giải pháp
thay thế an toàn cho các giao thức truy cập từ xa không được mã hóa như
Telnet, rlogin và rsh {[}1{]}. Phiên bản hiện tại được sử dụng rộng rãi
là SSH-2, được chuẩn hóa bởi IETF (Internet Engineering Task Force)
thông qua chuỗi tài liệu RFC 4250--4256 vào năm 2006 {[}2{]}.

Kiến trúc SSH-2 được tổ chức thành ba lớp giao thức phân tầng:

\textbf{Lớp giao thức vận chuyển (Transport Layer Protocol --- RFC
4253):} Lớp này cung cấp khả năng xác thực máy chủ (server
authentication), bảo mật dữ liệu (confidentiality) và toàn vẹn dữ liệu
(integrity). Quá trình bắt tay (handshake) bao gồm trao đổi phiên bản
giao thức, thương lượng thuật toán mật mã, trao đổi khóa Diffie-Hellman,
và xác thực máy chủ bằng khóa công khai. Sau khi hoàn tất, một kênh liên
lạc được mã hóa được thiết lập giữa client và server {[}3{]}.

\textbf{Lớp giao thức xác thực người dùng (User Authentication Protocol
--- RFC 4252):} Lớp này xử lý việc xác thực danh tính người dùng. SSH-2
hỗ trợ nhiều phương pháp xác thực bao gồm: xác thực bằng mật khẩu
(password), xác thực bằng khóa công khai (publickey), xác thực dựa trên
máy chủ (hostbased), và xác thực bàn phím tương tác
(keyboard-interactive). Trong thực tế, hai phương pháp phổ biến nhất là
xác thực mật khẩu và xác thực khóa công khai {[}4{]}.

\textbf{Lớp giao thức kết nối (Connection Protocol --- RFC 4254):} Lớp
này cho phép ghép kênh (multiplexing) nhiều kênh logic trên một kết nối
SSH duy nhất, hỗ trợ các tính năng như thực thi lệnh từ xa, chuyển tiếp
cổng (port forwarding), và truyền tải tệp tin qua SFTP/SCP {[}5{]}.

\subsection{2.1.2 Cơ chế xác thực mật khẩu và điểm
yếu}\label{cux1a1-chux1ebf-xuxe1c-thux1ef1c-mux1eadt-khux1ea9u-vuxe0-ux111iux1ec3m-yux1ebfu}

Cơ chế xác thực mật khẩu (password authentication) trong SSH hoạt động
theo quy trình sau: (1) Client gửi yêu cầu xác thực kèm tên đăng nhập và
mật khẩu qua kênh đã được mã hóa; (2) Server kiểm tra thông tin đăng
nhập với cơ sở dữ liệu người dùng (thường là /etc/shadow trên Linux);
(3) Server phản hồi thành công (SSH\_MSG\_USERAUTH\_SUCCESS) hoặc thất
bại (SSH\_MSG\_USERAUTH\_FAILURE) {[}4{]}.

Mặc dù mật khẩu được truyền qua kênh mã hóa, cơ chế xác thực mật khẩu
tồn tại nhiều điểm yếu cố hữu khiến nó trở thành mục tiêu của tấn công
brute-force:

\begin{itemize}
\tightlist
\item
  \textbf{Không giới hạn số lần thử mặc định:} Cấu hình mặc định của
  OpenSSH cho phép nhiều lần thử xác thực trong một phiên kết nối (thông
  số MaxAuthTries, mặc định là 6) và không giới hạn tổng số phiên kết
  nối từ một IP {[}6{]}.
\item
  \textbf{Phụ thuộc vào độ mạnh mật khẩu:} Hiệu quả bảo mật hoàn toàn
  phụ thuộc vào việc người dùng chọn mật khẩu đủ mạnh, điều mà thực tế
  cho thấy thường không được đảm bảo {[}7{]}.
\item
  \textbf{Không có cơ chế chống tự động:} Không giống như các ứng dụng
  web có thể sử dụng CAPTCHA, giao thức SSH không có cơ chế tích hợp để
  phân biệt giữa người dùng thực và công cụ tấn công tự động {[}8{]}.
\end{itemize}

\subsection{2.1.3 Nhật ký SSH và thông tin hữu ích cho phát hiện tấn
công}\label{nhux1eadt-kuxfd-ssh-vuxe0-thuxf4ng-tin-hux1eefu-uxedch-cho-phuxe1t-hiux1ec7n-tux1ea5n-cuxf4ng}

Máy chủ SSH ghi nhận các sự kiện xác thực vào tệp nhật ký hệ thống
(thường là /var/log/auth.log trên Debian/Ubuntu hoặc /var/log/secure
trên CentOS/RHEL). Mỗi sự kiện chứa các thông tin quan trọng bao gồm:
dấu thời gian (timestamp), địa chỉ IP nguồn, tên đăng nhập được sử dụng,
kết quả xác thực (thành công/thất bại), phương pháp xác thực, và số cổng
nguồn {[}9{]}.

Ví dụ về các mẫu nhật ký SSH:

\begin{verbatim}
Failed password for root from 192.168.1.100 port 52413 ssh2
Failed password for invalid user admin from 10.0.0.5 port 39821 ssh2
Accepted password for user1 from 172.16.0.10 port 48732 ssh2
Connection closed by authenticating user root 192.168.1.100 port 52413
\end{verbatim}

Các trường thông tin này cung cấp nền tảng dữ liệu phong phú cho việc
trích xuất đặc trưng hành vi và phân tích bất thường, phục vụ mục đích
phát hiện tấn công brute-force.

\section{2.2 Phân loại tấn công
Brute-force}\label{phuxe2n-loux1ea1i-tux1ea5n-cuxf4ng-brute-force}

\subsection{2.2.1 Định nghĩa và nguyên
lý}\label{ux111ux1ecbnh-nghux129a-vuxe0-nguyuxean-luxfd}

Tấn công brute-force (tấn công dò mật khẩu) là hình thức tấn công trong
đó kẻ tấn công thử lần lượt các tổ hợp tên đăng nhập và mật khẩu cho đến
khi tìm được thông tin đăng nhập hợp lệ. Về mặt lý thuyết, brute-force
luôn thành công nếu không gian tìm kiếm hữu hạn và kẻ tấn công có đủ
thời gian, tuy nhiên trong thực tế, thời gian cần thiết phụ thuộc vào độ
phức tạp của mật khẩu và tốc độ thử {[}10{]}.

Với bảng chữ cái có kích thước \textbar A\textbar{} và mật khẩu có độ
dài L, không gian tìm kiếm tối đa là:

\[S = |A|^L\]

Ví dụ, với mật khẩu gồm chữ cái thường (26 ký tự) và chữ số (10 ký tự),
có độ dài 8 ký tự, không gian tìm kiếm là 36\^{}8 ≈ 2,82 × 10\^{}12 tổ
hợp.

\subsection{2.2.2 Các biến thể tấn công brute-force
SSH}\label{cuxe1c-biux1ebfn-thux1ec3-tux1ea5n-cuxf4ng-brute-force-ssh}

Dựa trên phân tích các công trình nghiên cứu và dữ liệu thực tế, các
biến thể tấn công brute-force SSH có thể được phân loại như sau:

\textbf{a) Tấn công brute-force cổ điển (Classic brute-force):} Kẻ tấn
công thử tất cả tổ hợp có thể với tốc độ cao nhất, thường nhắm vào các
tài khoản phổ biến như root, admin, user. Đặc điểm nhận dạng bao gồm:
tần suất đăng nhập thất bại rất cao (hàng trăm đến hàng nghìn lần/phút),
khoảng cách thời gian giữa các lần thử rất ngắn và đồng đều, sử dụng
nhiều tên đăng nhập khác nhau {[}11{]}.

\textbf{b) Tấn công từ điển (Dictionary attack):} Sử dụng danh sách mật
khẩu phổ biến được biên soạn sẵn (wordlist), thay vì thử tất cả tổ hợp.
Các danh sách phổ biến bao gồm RockYou (14 triệu mật khẩu), SecLists, và
các danh sách được tùy chỉnh theo mục tiêu. Phương pháp này hiệu quả hơn
brute-force cổ điển vì khai thác thói quen sử dụng mật khẩu yếu của
người dùng {[}12{]}.

\textbf{c) Tấn công chậm (Slow brute-force / Low-and-slow attack):} Kẻ
tấn công cố tình giảm tốc độ tấn công, thực hiện chỉ vài lần thử mỗi
phút hoặc thậm chí mỗi giờ, nhằm lẩn tránh các cơ chế phát hiện dựa trên
ngưỡng tần suất. Đây là hình thức tấn công khó phát hiện nhất với các
phương pháp truyền thống, vì hành vi tấn công gần giống với hành vi của
người dùng hợp pháp đăng nhập thất bại {[}13{]}.

\textbf{d) Tấn công phân tán (Distributed brute-force):} Sử dụng mạng
botnet hoặc dịch vụ proxy để phân tán tấn công từ nhiều địa chỉ IP khác
nhau. Mỗi IP chỉ thực hiện một số ít lần thử, khiến việc phát hiện dựa
trên số lần thất bại từ một IP trở nên không hiệu quả. Theo Owens và
Matthews {[}14{]}, tấn công phân tán chiếm tỷ lệ ngày càng tăng trong
các cuộc tấn công brute-force SSH được ghi nhận.

\textbf{e) Tấn công credential stuffing:} Sử dụng các cặp tên đăng
nhập-mật khẩu bị rò rỉ từ các vụ vi phạm dữ liệu của các dịch vụ khác.
Kẻ tấn công khai thác thói quen sử dụng lại mật khẩu (password reuse)
của người dùng. Hình thức này đặc biệt nguy hiểm vì tỷ lệ thành công cao
hơn nhiều so với brute-force ngẫu nhiên {[}15{]}.

\subsection{2.2.3 Công cụ tấn công phổ
biến}\label{cuxf4ng-cux1ee5-tux1ea5n-cuxf4ng-phux1ed5-biux1ebfn}

Các công cụ thường được sử dụng trong tấn công brute-force SSH bao gồm:
Hydra (hỗ trợ đa giao thức, tấn công song song), Medusa (tối ưu cho tấn
công tốc độ cao), Ncrack (của dự án Nmap), Patator (framework linh hoạt
viết bằng Python), và các script tùy chỉnh sử dụng thư viện Paramiko
hoặc libssh {[}16{]}. Sự đa dạng của công cụ tấn công tạo ra các mẫu
hành vi khác nhau trong dữ liệu nhật ký, đặt ra yêu cầu về khả năng tổng
quát hóa (generalization) của hệ thống phát hiện.

\section{2.3 Các phương pháp phát hiện truyền
thống}\label{cuxe1c-phux1b0ux1a1ng-phuxe1p-phuxe1t-hiux1ec7n-truyux1ec1n-thux1ed1ng}

\subsection{2.3.1 Phương pháp dựa trên ngưỡng tĩnh (Static
threshold-based)}\label{phux1b0ux1a1ng-phuxe1p-dux1ef1a-truxean-ngux1b0ux1ee1ng-tux129nh-static-threshold-based}

Phương pháp phổ biến nhất trong thực tế là thiết lập ngưỡng cố định cho
số lần đăng nhập thất bại. Fail2Ban --- công cụ phòng chống xâm nhập mã
nguồn mở phổ biến nhất --- hoạt động bằng cách giám sát tệp nhật ký hệ
thống và chặn IP vi phạm qua iptables/nftables khi số lần thất bại vượt
ngưỡng trong một khoảng thời gian xác định {[}17{]}.

Cấu hình mặc định của Fail2Ban cho SSH thường là: maxretry = 5 (chặn sau
5 lần thất bại), findtime = 600 (trong cửa sổ 10 phút), bantime = 600
(chặn trong 10 phút). Phương pháp này đơn giản, dễ triển khai và hiệu
quả với tấn công brute-force tốc độ cao. Tuy nhiên, như đã phân tích ở
Chương 1, ngưỡng tĩnh có những hạn chế cố hữu:

\begin{itemize}
\tightlist
\item
  Không thể phát hiện tấn công chậm với tốc độ dưới ngưỡng
\item
  Không thích ứng với sự biến động tự nhiên của lưu lượng
\item
  Tỷ lệ false positive cao trong giờ cao điểm khi nhiều người dùng hợp
  pháp đăng nhập đồng thời
\item
  Dễ bị lẩn tránh bằng tấn công phân tán từ nhiều IP
\end{itemize}

\subsection{2.3.2 Phương pháp dựa trên danh sách
(List-based)}\label{phux1b0ux1a1ng-phuxe1p-dux1ef1a-truxean-danh-suxe1ch-list-based}

Phương pháp này sử dụng danh sách đen (blacklist) và danh sách trắng
(whitelist) để kiểm soát truy cập. Các dịch vụ như AbuseIPDB, Spamhaus,
và Blocklist.de cung cấp danh sách IP được biết đến với các hoạt động
độc hại {[}18{]}. Hạn chế chính là tính phản ứng --- IP chỉ được đưa vào
danh sách đen sau khi đã thực hiện tấn công tại nơi khác, và kẻ tấn công
có thể dễ dàng chuyển sang sử dụng IP mới.

\subsection{2.3.3 Phương pháp dựa trên chữ ký
(Signature-based)}\label{phux1b0ux1a1ng-phuxe1p-dux1ef1a-truxean-chux1eef-kuxfd-signature-based}

Các hệ thống phát hiện xâm nhập (IDS) như Snort và Suricata có thể phát
hiện tấn công brute-force SSH bằng cách so khớp mẫu lưu lượng mạng với
các chữ ký (signatures/rules) đã biết {[}19{]}. Ví dụ, rule Snort có thể
phát hiện nhiều gói tin SSH\_MSG\_USERAUTH\_REQUEST liên tiếp. Hạn chế
của phương pháp này là không thể phát hiện các biến thể tấn công mới
chưa có chữ ký, và hiệu suất giảm khi lưu lượng SSH được mã hóa
end-to-end khiến việc phân tích nội dung gói tin không khả thi.

\subsection{2.3.4 Phương pháp xác thực nâng
cao}\label{phux1b0ux1a1ng-phuxe1p-xuxe1c-thux1ef1c-nuxe2ng-cao}

Các biện pháp phòng chống ở mức xác thực bao gồm: chuyển sang sử dụng
xác thực khóa công khai (public key authentication), triển khai xác thực
hai yếu tố (two-factor authentication) qua PAM modules, thay đổi cổng
SSH mặc định, và sử dụng port knocking {[}20{]}. Các biện pháp này hiệu
quả nhưng không phải lúc nào cũng khả thi trong mọi môi trường, đặc biệt
là các hệ thống legacy hoặc môi trường đa người dùng.

\section{2.4 Học máy trong phát hiện xâm
nhập}\label{hux1ecdc-muxe1y-trong-phuxe1t-hiux1ec7n-xuxe2m-nhux1eadp}

\subsection{2.4.1 Tổng quan về ứng dụng học máy trong an ninh
mạng}\label{tux1ed5ng-quan-vux1ec1-ux1ee9ng-dux1ee5ng-hux1ecdc-muxe1y-trong-an-ninh-mux1ea1ng}

Ứng dụng học máy trong phát hiện xâm nhập mạng (Network Intrusion
Detection System --- NIDS) đã được nghiên cứu rộng rãi trong hai thập kỷ
qua. Buczak và Guven {[}21{]} đã tổng hợp một khảo sát toàn diện về các
phương pháp khai phá dữ liệu và học máy cho an ninh mạng, chỉ ra rằng
các thuật toán học máy có thể đạt hiệu suất phát hiện cao hơn đáng kể so
với các phương pháp dựa trên quy tắc trong nhiều tình huống.

Các phương pháp học máy trong phát hiện xâm nhập được phân loại theo
nhiều tiêu chí:

\textbf{Theo phương thức học:} - \emph{Học có giám sát (Supervised
Learning):} Yêu cầu dữ liệu huấn luyện được gán nhãn (bình thường/tấn
công). Các thuật toán phổ biến bao gồm Random Forest, Support Vector
Machine, Neural Networks, và Gradient Boosting. Ưu điểm là độ chính xác
cao khi có đủ dữ liệu huấn luyện chất lượng; nhược điểm là phụ thuộc vào
dữ liệu gán nhãn và khó phát hiện các loại tấn công chưa biết {[}22{]}.
- \emph{Học không giám sát (Unsupervised Learning):} Không cần dữ liệu
gán nhãn, thay vào đó học mẫu hành vi bình thường và coi các điểm dữ
liệu lệch khỏi mẫu là bất thường. Ưu điểm là có thể phát hiện các loại
tấn công chưa biết (zero-day attacks); nhược điểm là tỷ lệ false
positive thường cao hơn {[}23{]}. - \emph{Học bán giám sát
(Semi-supervised Learning):} Kết hợp một lượng nhỏ dữ liệu gán nhãn với
lượng lớn dữ liệu không gán nhãn. Phù hợp với thực tế khi dữ liệu gán
nhãn khan hiếm và tốn kém {[}24{]}.

\textbf{Theo phương pháp phát hiện:} - \emph{Phát hiện dựa trên sai lệch
(Misuse detection):} Xây dựng mô hình cho các hành vi tấn công đã biết,
phát hiện tấn công khi mẫu lưu lượng khớp với mô hình. Tương đương với
phương pháp dựa trên chữ ký nhưng sử dụng mô hình học máy thay vì quy
tắc thủ công {[}25{]}. - \emph{Phát hiện bất thường (Anomaly
detection):} Xây dựng mô hình cho hành vi bình thường, phát hiện tấn
công khi mẫu lưu lượng lệch khỏi mô hình. Có khả năng phát hiện tấn công
chưa biết nhưng cần xác định ngưỡng phân biệt hợp lý {[}26{]}.

\subsection{2.4.2 Lý do chọn phương pháp phát hiện bất thường không giám
sát}\label{luxfd-do-chux1ecdn-phux1b0ux1a1ng-phuxe1p-phuxe1t-hiux1ec7n-bux1ea5t-thux1b0ux1eddng-khuxf4ng-giuxe1m-suxe1t}

Trong bối cảnh phát hiện tấn công brute-force SSH, phương pháp phát hiện
bất thường không giám sát được lựa chọn trong nghiên cứu này dựa trên
các lý do sau:

Thứ nhất, bản chất của tấn công brute-force liên tục tiến hóa. Kẻ tấn
công không ngừng thay đổi kỹ thuật để lẩn tránh phát hiện, khiến các mô
hình học có giám sát được huấn luyện trên dữ liệu tấn công cũ có thể
không nhận diện được các biến thể mới. Phương pháp phát hiện bất thường,
bằng cách mô hình hóa hành vi bình thường thay vì hành vi tấn công, có
khả năng tổng quát hóa tốt hơn với các hình thức tấn công mới {[}27{]}.

Thứ hai, trong thực tế vận hành, dữ liệu tấn công được gán nhãn chính
xác rất khó thu thập. Mặc dù dữ liệu từ honeypot cung cấp các mẫu tấn
công thực tế, việc gán nhãn chính xác từng phiên đăng nhập trong môi
trường vận hành là không khả thi ở quy mô lớn {[}28{]}.

Thứ ba, các thuật toán phát hiện bất thường không giám sát như Isolation
Forest có ưu điểm về hiệu suất tính toán, cho phép xử lý dữ liệu theo
thời gian thực --- yêu cầu quan trọng cho hệ thống phát hiện xâm nhập
{[}29{]}.

\section{2.5 Các thuật toán phát hiện bất
thường}\label{cuxe1c-thuux1eadt-touxe1n-phuxe1t-hiux1ec7n-bux1ea5t-thux1b0ux1eddng}

\subsection{2.5.1 Isolation Forest (IF)}\label{isolation-forest-if}

\subsubsection{a) Nguyên lý hoạt
động}\label{a-nguyuxean-luxfd-houx1ea1t-ux111ux1ed9ng}

Isolation Forest được đề xuất bởi Liu, Ting và Zhou vào năm 2008 tại Đại
học Monash, Úc, và được công bố chính thức trên tạp chí ACM Transactions
on Knowledge Discovery from Data năm 2012 {[}29{]}. Khác với hầu hết các
thuật toán phát hiện bất thường dựa trên khoảng cách hoặc mật độ,
Isolation Forest dựa trên nguyên lý ``cô lập'' (isolation): các điểm bất
thường, do có giá trị đặc trưng khác biệt với đa số, sẽ bị cô lập (tách
biệt khỏi các điểm khác) nhanh hơn so với các điểm bình thường khi thực
hiện phân hoạch ngẫu nhiên.

Thuật toán xây dựng một tập hợp các cây cô lập (Isolation Trees ---
iTrees) bằng cách lặp lại quá trình sau: chọn ngẫu nhiên một đặc trưng,
chọn ngẫu nhiên một giá trị phân tách trong khoảng {[}min, max{]} của
đặc trưng đó, và chia dữ liệu thành hai nhánh. Quá trình tiếp tục cho
đến khi mỗi điểm dữ liệu được cô lập hoặc đạt độ sâu tối đa.

\subsubsection{b) Công thức toán
học}\label{b-cuxf4ng-thux1ee9c-touxe1n-hux1ecdc}

Cho tập dữ liệu X = \{x₁, x₂, \ldots, xₙ\} với n điểm dữ liệu trong
không gian d chiều.

\textbf{Xây dựng Isolation Tree:} Tại mỗi nút trong (internal node),
thuật toán chọn ngẫu nhiên một đặc trưng q ∈ \{1, 2, \ldots, d\} và một
giá trị phân tách p được chọn ngẫu nhiên đều trong khoảng {[}min(xq),
max(xq){]}, trong đó xq là giá trị của đặc trưng q. Dữ liệu được chia
thành hai tập con: tập con trái chứa các điểm có xq \textless{} p và tập
con phải chứa các điểm có xq ≥ p.

\textbf{Độ dài đường đi (Path Length):} Cho một điểm dữ liệu x, độ dài
đường đi h(x) trong một Isolation Tree T là số cạnh từ nút gốc đến nút
lá chứa x. Đối với nút ngoài (external node) ở độ sâu đã đạt giới hạn,
cần bổ sung ước lượng cho phần còn lại:

\[h(x) = e + c(T.size)\]

trong đó e là số cạnh từ gốc đến nút hiện tại, T.size là kích thước tập
con tại nút đó, và c(n) là độ dài đường đi trung bình của cây tìm kiếm
nhị phân không thành công (unsuccessful search in BST):

\[c(n) = 2H(n-1) - \frac{2(n-1)}{n}\]

với H(i) = ln(i) + γ là số điều hòa (harmonic number), γ ≈ 0.5772 là
hằng số Euler-Mascheroni.

\textbf{Điểm bất thường (Anomaly Score):} Điểm bất thường của x được
tính dựa trên kỳ vọng của độ dài đường đi trên tất cả t cây trong rừng:

\[s(x, n) = 2^{-\frac{E[h(x)]}{c(n)}}\]

trong đó E{[}h(x){]} là trung bình h(x) trên t Isolation Trees.

Diễn giải: - s → 1: điểm x có khả năng là bất thường (đường đi trung
bình ngắn, bị cô lập nhanh) - s → 0.5: toàn bộ tập dữ liệu không có bất
thường rõ rệt - s → 0: điểm x là bình thường (đường đi trung bình dài,
khó bị cô lập)

\subsubsection{c) Ưu điểm trong phát hiện tấn công
brute-force}\label{c-ux1b0u-ux111iux1ec3m-trong-phuxe1t-hiux1ec7n-tux1ea5n-cuxf4ng-brute-force}

Isolation Forest có nhiều ưu điểm phù hợp với bài toán phát hiện tấn
công brute-force SSH:

\begin{itemize}
\tightlist
\item
  \textbf{Độ phức tạp tuyến tính:} O(t · n · log ψ), trong đó t là số
  cây, n là kích thước dữ liệu huấn luyện, ψ là kích thước mẫu con
  (sub-sampling size). Điều này cho phép xử lý thời gian thực với dữ
  liệu nhật ký quy mô lớn {[}29{]}.
\item
  \textbf{Hiệu quả với dữ liệu chiều cao:} Không bị ảnh hưởng bởi ``lời
  nguyền chiều'' (curse of dimensionality) như các phương pháp dựa trên
  khoảng cách, phù hợp với bộ đặc trưng 14 chiều {[}30{]}.
\item
  \textbf{Không cần giả định phân phối:} Không yêu cầu dữ liệu tuân theo
  phân phối cụ thể, phù hợp với dữ liệu nhật ký SSH có phân phối phức
  tạp {[}29{]}.
\item
  \textbf{Khả năng xử lý outlier swamping và masking:} Cơ chế lấy mẫu
  con (sub-sampling) giúp giảm thiểu hiện tượng các điểm bất thường ảnh
  hưởng lẫn nhau {[}31{]}.
\end{itemize}

\subsection{2.5.2 Local Outlier Factor
(LOF)}\label{local-outlier-factor-lof}

\subsubsection{a) Nguyên lý hoạt
động}\label{a-nguyuxean-luxfd-houx1ea1t-ux111ux1ed9ng-1}

Local Outlier Factor (LOF) được đề xuất bởi Breunig, Kriegel, Ng và
Sander vào năm 2000 {[}32{]}. LOF là phương pháp phát hiện bất thường
dựa trên mật độ cục bộ (local density-based), so sánh mật độ cục bộ của
mỗi điểm dữ liệu với mật độ cục bộ của các điểm lân cận. Ý tưởng cốt lõi
là: một điểm bất thường có mật độ cục bộ thấp hơn đáng kể so với các
điểm lân cận.

\subsubsection{b) Công thức toán
học}\label{b-cuxf4ng-thux1ee9c-touxe1n-hux1ecdc-1}

\textbf{Khoảng cách tiếp cận (Reachability Distance):} Cho hai điểm p và
o, khoảng cách tiếp cận bậc k được định nghĩa:

\[reach\text{-}dist_k(p, o) = \max\{k\text{-}dist(o), \; d(p, o)\}\]

trong đó k-dist(o) là khoảng cách đến điểm lân cận thứ k của o, và d(p,
o) là khoảng cách Euclid giữa p và o.

\textbf{Mật độ tiếp cận cục bộ (Local Reachability Density):} Mật độ
tiếp cận cục bộ của điểm p được tính:

\[lrd_k(p) = \frac{1}{\frac{\sum_{o \in N_k(p)} reach\text{-}dist_k(p, o)}{|N_k(p)|}}\]

trong đó N\_k(p) là tập k điểm lân cận gần nhất của p.

\textbf{Hệ số LOF:} Hệ số LOF của điểm p được tính:

\[LOF_k(p) = \frac{\sum_{o \in N_k(p)} \frac{lrd_k(o)}{lrd_k(p)}}{|N_k(p)|} = \frac{1}{|N_k(p)|} \sum_{o \in N_k(p)} \frac{lrd_k(o)}{lrd_k(p)}\]

Diễn giải: - LOF ≈ 1: mật độ cục bộ của p tương đương với các điểm lân
cận (bình thường) - LOF \textgreater\textgreater{} 1: mật độ cục bộ của
p thấp hơn nhiều so với các điểm lân cận (bất thường)

\subsubsection{c) Hạn chế trong bối cảnh phát hiện tấn công
SSH}\label{c-hux1ea1n-chux1ebf-trong-bux1ed1i-cux1ea3nh-phuxe1t-hiux1ec7n-tux1ea5n-cuxf4ng-ssh}

Mặc dù LOF hiệu quả trong nhiều bài toán phát hiện bất thường, thuật
toán có một số hạn chế trong bối cảnh phát hiện tấn công SSH thời gian
thực:

\begin{itemize}
\tightlist
\item
  \textbf{Độ phức tạp cao:} O(n²) cho tính toán khoảng cách k-nearest
  neighbors, trở nên tốn kém khi xử lý dữ liệu nhật ký liên tục
  {[}33{]}.
\item
  \textbf{Nhạy cảm với tham số k:} Hiệu năng phụ thuộc đáng kể vào lựa
  chọn giá trị k (số lân cận), và không có phương pháp xác định k tối ưu
  cho mọi tập dữ liệu {[}32{]}.
\item
  \textbf{Yêu cầu lưu trữ toàn bộ dữ liệu huấn luyện:} Để tính LOF cho
  điểm dữ liệu mới, cần truy cập toàn bộ tập dữ liệu huấn luyện, gây áp
  lực về bộ nhớ {[}34{]}.
\end{itemize}

\subsection{2.5.3 One-Class Support Vector Machine
(OCSVM)}\label{one-class-support-vector-machine-ocsvm}

\subsubsection{a) Nguyên lý hoạt
động}\label{a-nguyuxean-luxfd-houx1ea1t-ux111ux1ed9ng-2}

One-Class SVM được đề xuất bởi Schölkopf, Platt, Shawe-Taylor, Smola và
Williamson năm 2001 {[}35{]}. Thuật toán mở rộng SVM truyền thống cho
bài toán phát hiện bất thường bằng cách tìm siêu phẳng (hyperplane)
trong không gian đặc trưng có chiều cao hơn sao cho tách biệt tối đa các
điểm dữ liệu huấn luyện (bình thường) khỏi gốc tọa độ. Các điểm nằm phía
gốc tọa độ so với siêu phẳng được coi là bất thường.

\subsubsection{b) Công thức toán
học}\label{b-cuxf4ng-thux1ee9c-touxe1n-hux1ecdc-2}

\textbf{Bài toán tối ưu:} Cho tập dữ liệu huấn luyện X = \{x₁, \ldots,
xₙ\}, OCSVM giải bài toán tối ưu:

\[\min_{w, \xi, \rho} \frac{1}{2}\|w\|^2 + \frac{1}{\nu n}\sum_{i=1}^{n}\xi_i - \rho\]

với các ràng buộc:

\[w \cdot \Phi(x_i) \geq \rho - \xi_i, \quad \xi_i \geq 0, \quad i = 1, ..., n\]

trong đó: - w là vector pháp tuyến của siêu phẳng trong không gian đặc
trưng - Φ(·) là hàm ánh xạ sang không gian đặc trưng chiều cao - ρ là
khoảng cách từ gốc tọa độ đến siêu phẳng - ξᵢ là biến lỏng (slack
variables) cho phép sai lệch - ν ∈ (0, 1{]} là tham số điều khiển tỷ lệ
bất thường dự kiến

\textbf{Hàm quyết định:} Sau khi giải bài toán đối ngẫu, hàm quyết định
cho điểm mới x:

\[f(x) = \text{sign}\left(\sum_{i=1}^{n} \alpha_i K(x_i, x) - \rho\right)\]

trong đó K(xᵢ, x) là hàm nhân (kernel function). Hàm nhân RBF (Radial
Basis Function) phổ biến nhất:

\[K(x_i, x) = \exp\left(-\gamma\|x_i - x\|^2\right)\]

f(x) = +1 cho điểm bình thường, f(x) = -1 cho điểm bất thường.

\subsubsection{c) Ưu và nhược
điểm}\label{c-ux1b0u-vuxe0-nhux1b0ux1ee3c-ux111iux1ec3m}

\textbf{Ưu điểm:} OCSVM có cơ sở lý thuyết vững chắc từ lý thuyết tối ưu
và kernel methods, cho kết quả tốt khi ranh giới quyết định phức tạp nhờ
kernel trick, và tham số ν có ý nghĩa trực quan --- giới hạn trên của tỷ
lệ bất thường {[}36{]}.

\textbf{Nhược điểm:} Độ phức tạp huấn luyện O(n²) đến O(n³) với SVM
solver truyền thống, nhạy cảm cao với lựa chọn kernel và tham số (γ, ν),
và không có cơ chế tự nhiên để xử lý dữ liệu streaming --- cần tái huấn
luyện khi dữ liệu thay đổi {[}37{]}.

\subsection{2.5.4 So sánh ba thuật
toán}\label{so-suxe1nh-ba-thuux1eadt-touxe1n}

Về mặt lý thuyết, Isolation Forest có ưu thế về hiệu suất tính toán (O(n
log n) so với O(n²) của LOF và O(n²-n³) của OCSVM), khả năng xử lý dữ
liệu chiều cao, và không yêu cầu giả định phân phối. LOF có ưu thế trong
việc phát hiện bất thường cục bộ (local outliers) --- các điểm chỉ bất
thường khi xét trong ngữ cảnh lân cận. OCSVM có cơ sở lý thuyết vững
chắc nhất và cho phép kiểm soát chặt chẽ tỷ lệ bất thường qua tham số ν
{[}38{]}.

Trong nghiên cứu này, Isolation Forest được chọn làm thuật toán chính
dựa trên cân nhắc về hiệu suất tính toán phù hợp với xử lý thời gian
thực, khả năng tổng quát hóa tốt, và kết quả đánh giá thực nghiệm. LOF
và OCSVM được sử dụng làm cơ sở đối sánh (benchmark) để đánh giá khách
quan hiệu năng.

\section{2.6 Phương pháp ngưỡng động trong phát hiện bất
thường}\label{phux1b0ux1a1ng-phuxe1p-ngux1b0ux1ee1ng-ux111ux1ed9ng-trong-phuxe1t-hiux1ec7n-bux1ea5t-thux1b0ux1eddng}

\subsection{2.6.1 Hạn chế của ngưỡng
tĩnh}\label{hux1ea1n-chux1ebf-cux1ee7a-ngux1b0ux1ee1ng-tux129nh}

Trong các hệ thống phát hiện bất thường, ngưỡng (threshold) đóng vai trò
quyết định trong việc phân loại một điểm dữ liệu là bình thường hay bất
thường. Ngưỡng tĩnh --- một giá trị cố định được xác định trước --- có
ưu điểm về tính đơn giản nhưng không phù hợp với dữ liệu có đặc tính
thay đổi theo thời gian (non-stationary data) {[}39{]}.

Trong bối cảnh giám sát SSH, lưu lượng truy cập thay đổi đáng kể theo
thời gian: giờ làm việc có nhiều đăng nhập hợp pháp hơn ngoài giờ, ngày
cuối tuần khác ngày thường, và các sự kiện đặc biệt (bảo trì hệ thống,
triển khai ứng dụng) tạo ra các đỉnh lưu lượng bất thường nhưng hợp
pháp. Ngưỡng tĩnh không thể thích ứng với những biến động này, dẫn đến
hai vấn đề: cảnh báo sai khi lưu lượng cao bất thường nhưng hợp pháp
(false positive), và bỏ sót tấn công khi lưu lượng nền thấp (false
negative).

\subsection{2.6.2 Exponentially Weighted Moving Average
(EWMA)}\label{exponentially-weighted-moving-average-ewma}

EWMA là phương pháp làm mượt chuỗi thời gian (time series smoothing) gán
trọng số giảm dần theo hàm mũ cho các quan sát cũ hơn. Được giới thiệu
bởi Roberts năm 1959 trong bối cảnh kiểm soát chất lượng thống kê
{[}40{]}, EWMA đã được ứng dụng rộng rãi trong nhiều lĩnh vực bao gồm
phát hiện bất thường mạng {[}41{]}.

\textbf{Công thức EWMA:}

\[\hat{\mu}_t = \alpha \cdot x_t + (1 - \alpha) \cdot \hat{\mu}_{t-1}\]

trong đó: - x\_t là giá trị quan sát tại thời điểm t - μ̂\_t là giá trị
EWMA tại thời điểm t - α ∈ (0, 1{]} là hệ số làm mượt (smoothing factor)

Giá trị α nhỏ cho kết quả mượt hơn (ít nhạy với biến động ngắn hạn), α
lớn cho kết quả phản ứng nhanh hơn với thay đổi. Trong phát hiện bất
thường, EWMA được sử dụng để ước lượng mức nền (baseline) của anomaly
score, từ đó xác định ngưỡng:

\[\sigma_t^2 = \alpha \cdot (x_t - \hat{\mu}_t)^2 + (1 - \alpha) \cdot \sigma_{t-1}^2\]

\[threshold_t = \hat{\mu}_t + k \cdot \sigma_t\]

trong đó k là hệ số nhạy (sensitivity factor), thường nằm trong khoảng
{[}2, 3{]} theo quy tắc phân phối chuẩn.

\subsection{2.6.3 Adaptive Percentile}\label{adaptive-percentile}

Phương pháp Adaptive Percentile xác định ngưỡng dựa trên phân vị
(percentile) của phân phối anomaly score trong một cửa sổ thời gian
trượt (sliding window). Thay vì giả định phân phối chuẩn như EWMA,
phương pháp này trực tiếp sử dụng phân phối thực nghiệm (empirical
distribution) của dữ liệu {[}42{]}.

\textbf{Công thức:}

\[threshold_t = P_q(S_W)\]

trong đó: - S\_W = \{s\_\{t-W+1\}, s\_\{t-W+2\}, \ldots, s\_t\} là tập
anomaly scores trong cửa sổ W gần nhất - P\_q là phân vị thứ q (ví dụ q
= 95 cho phân vị 95\%)

Ưu điểm chính là không yêu cầu giả định phân phối và tự nhiên thích ứng
với sự thay đổi của phân phối dữ liệu. Nhược điểm là cần lưu trữ dữ liệu
trong cửa sổ và nhạy cảm với kích thước cửa sổ W.

\subsection{2.6.4 Phương pháp kết hợp EWMA-Adaptive
Percentile}\label{phux1b0ux1a1ng-phuxe1p-kux1ebft-hux1ee3p-ewma-adaptive-percentile}

Nghiên cứu này đề xuất kết hợp EWMA và Adaptive Percentile để tận dụng
ưu điểm của cả hai phương pháp: EWMA cung cấp khả năng làm mượt và theo
dõi xu hướng dài hạn, trong khi Adaptive Percentile phản ánh chính xác
phân phối thực tế trong ngắn hạn.

\textbf{Ngưỡng kết hợp:}

\[threshold_t = \beta \cdot T_{EWMA,t} + (1 - \beta) \cdot T_{Percentile,t}\]

trong đó: - T\_\{EWMA,t\} = μ̂\emph{t + k · σ\_t là ngưỡng EWMA -
T}\{Percentile,t\} = P\_q(S\_W) là ngưỡng Adaptive Percentile - β ∈
{[}0, 1{]} là trọng số cân bằng giữa hai phương pháp

Cơ chế này cho phép hệ thống duy trì sự ổn định (stability) từ EWMA
trong khi vẫn phản ứng nhạy (responsiveness) với các thay đổi cục bộ từ
Adaptive Percentile. Tham số β có thể được điều chỉnh tùy theo yêu cầu:
β lớn ưu tiên ổn định (giảm false positive), β nhỏ ưu tiên nhạy bén
(giảm false negative).

Ưu thế của phương pháp kết hợp so với từng phương pháp đơn lẻ đã được
chỉ ra trong các nghiên cứu về giám sát an ninh mạng {[}43{]}, nơi mà sự
cân bằng giữa độ nhạy và tính ổn định là yếu tố then chốt.

\section{2.7 ELK Stack cho giám sát an
ninh}\label{elk-stack-cho-giuxe1m-suxe1t-an-ninh}

\subsection{2.7.1 Tổng quan về ELK
Stack}\label{tux1ed5ng-quan-vux1ec1-elk-stack}

ELK Stack là bộ ba công cụ mã nguồn mở được phát triển bởi Elastic N.V.,
bao gồm Elasticsearch, Logstash và Kibana. Được sử dụng rộng rãi trong
lĩnh vực quản lý nhật ký (log management) và phân tích dữ liệu, ELK
Stack đã trở thành nền tảng phổ biến cho Security Information and Event
Management (SIEM) và giám sát an ninh mạng {[}44{]}.

\textbf{Elasticsearch} là công cụ tìm kiếm và phân tích phân tán dựa
trên Apache Lucene. Elasticsearch lưu trữ dữ liệu dưới dạng tài liệu
JSON và hỗ trợ tìm kiếm toàn văn (full-text search), tổng hợp
(aggregation), và phân tích thời gian thực. Kiến trúc phân tán
(distributed) với khả năng sharding và replication đảm bảo tính mở rộng
và khả dụng cao {[}45{]}.

\textbf{Logstash} là đường ống xử lý dữ liệu (data processing pipeline)
phía máy chủ, có khả năng thu thập dữ liệu đồng thời từ nhiều nguồn
(inputs), chuyển đổi và làm giàu dữ liệu (filters), và gửi đến nhiều
đích (outputs). Logstash hỗ trợ hơn 200 plugin cho đầu vào, bộ lọc và
đầu ra, bao gồm đọc tệp nhật ký, phân tích cú pháp (parsing) với Grok,
và gửi dữ liệu đến Elasticsearch {[}46{]}.

\textbf{Kibana} là nền tảng trực quan hóa và khám phá dữ liệu, cung cấp
giao diện web để tương tác với dữ liệu trong Elasticsearch. Kibana hỗ
trợ tạo bảng điều khiển (dashboard), biểu đồ, bản đồ và cảnh báo, phục
vụ mục đích giám sát an ninh theo thời gian thực {[}47{]}.

\subsection{2.7.2 ELK Stack trong Security
Operations}\label{elk-stack-trong-security-operations}

Việc sử dụng ELK Stack cho giám sát an ninh mạng đã được ghi nhận trong
nhiều nghiên cứu và triển khai thực tế. Gonzalez và cộng sự {[}48{]} đã
chứng minh hiệu quả của ELK Stack trong việc phân tích nhật ký SSH với
khả năng xử lý hàng triệu sự kiện mỗi ngày. Nghiên cứu của Chuvakin và
cộng sự {[}49{]} chỉ ra rằng ELK Stack có thể thay thế các giải pháp
SIEM thương mại đắt đỏ cho các tổ chức vừa và nhỏ.

Trong kiến trúc giám sát SSH, quy trình hoạt động của ELK Stack như sau:

\begin{enumerate}
\def\labelenumi{\arabic{enumi}.}
\tightlist
\item
  \textbf{Thu thập (Collection):} Filebeat hoặc Logstash đọc tệp nhật ký
  SSH (/var/log/auth.log) theo thời gian thực
\item
  \textbf{Phân tích (Parsing):} Logstash filter sử dụng Grok pattern để
  trích xuất các trường thông tin: timestamp, IP nguồn, tên đăng nhập,
  kết quả xác thực
\item
  \textbf{Lưu trữ (Indexing):} Dữ liệu đã cấu trúc được đánh chỉ mục
  trong Elasticsearch với index pattern theo ngày
\item
  \textbf{Trực quan hóa (Visualization):} Kibana dashboard hiển thị các
  chỉ số an ninh: số lần đăng nhập thất bại theo IP, phân bố địa lý, xu
  hướng theo thời gian
\item
  \textbf{Cảnh báo (Alerting):} Elasticsearch Watcher hoặc Kibana
  Alerting tạo cảnh báo khi phát hiện dấu hiệu tấn công
\end{enumerate}

\subsection{2.7.3 Tích hợp AI với ELK
Stack}\label{tuxedch-hux1ee3p-ai-vux1edbi-elk-stack}

Việc tích hợp mô hình học máy với ELK Stack có thể thực hiện theo nhiều
cách. Elastic đã tích hợp sẵn module Machine Learning trong X-Pack (từ
phiên bản 5.x), hỗ trợ phát hiện bất thường tự động trên dữ liệu chuỗi
thời gian {[}50{]}. Tuy nhiên, module này có giới hạn về tùy chỉnh thuật
toán và yêu cầu giấy phép thương mại.

Cách tiếp cận trong nghiên cứu này là xây dựng pipeline tùy chỉnh: dữ
liệu từ Elasticsearch được truy xuất qua API, xử lý và trích xuất đặc
trưng bằng Python, đưa vào mô hình Isolation Forest để tính anomaly
score, và ghi kết quả ngược lại Elasticsearch để trực quan hóa trên
Kibana. Cách tiếp cận này mang lại sự linh hoạt tối đa trong việc tùy
chỉnh thuật toán và pipeline xử lý.

\section{2.8 Các công trình nghiên cứu liên
quan}\label{cuxe1c-cuxf4ng-truxecnh-nghiuxean-cux1ee9u-liuxean-quan}

\subsection{2.8.1 Nghiên cứu quốc
tế}\label{nghiuxean-cux1ee9u-quux1ed1c-tux1ebf}

\textbf{Najafabadi và cộng sự (2015)} {[}51{]} đã nghiên cứu phát hiện
tấn công brute-force SSH sử dụng dữ liệu NetFlow tổng hợp. Các tác giả
sử dụng thuật toán Random Forest trên bộ dữ liệu CERT NetFlow và đạt độ
chính xác 99\% trong phân biệt lưu lượng SSH bình thường và tấn công.
Tuy nhiên, nghiên cứu sử dụng học có giám sát với dữ liệu gán nhãn đầy
đủ, và đặc trưng được trích xuất từ NetFlow (dữ liệu tầng mạng) thay vì
nhật ký SSH (dữ liệu tầng ứng dụng), hạn chế khả năng phát hiện tấn công
tinh vi ở tầng ứng dụng.

\textbf{Hofstede, Pras và Sperotto (2018)} {[}52{]} đề xuất hệ thống SSH
Compromise Detection sử dụng flow-based features. Nghiên cứu khai thác
đặc trưng dựa trên luồng mạng (flow-based) kết hợp với Decision Tree và
Naive Bayes, đạt True Positive Rate trên 90\%. Đóng góp chính là bộ đặc
trưng flow-based phân biệt giữa giai đoạn brute-force và giai đoạn khai
thác sau khi xâm nhập thành công.

\textbf{Kumari và Jain (2020)} {[}53{]} nghiên cứu phương pháp phát hiện
bất thường dựa trên Isolation Forest cho hệ thống IoT. Các tác giả áp
dụng Isolation Forest trên bộ dữ liệu NSL-KDD và CICIDS2017, đạt
F1-score 89.7\% trên CICIDS2017. Nghiên cứu chứng minh tính hiệu quả của
Isolation Forest trong phát hiện bất thường mạng nhưng không tập trung
vào giao thức SSH cụ thể.

\textbf{Moustafa và Slay (2016)} {[}54{]} phát triển bộ dữ liệu
UNSW-NB15 và đánh giá nhiều thuật toán học máy trên bộ dữ liệu này, bao
gồm Isolation Forest. Kết quả cho thấy Isolation Forest đạt Detection
Rate 83.1\% với False Alarm Rate 14.2\% trên dữ liệu mạng tổng hợp.

\textbf{Starov và cộng sự (2019)} {[}55{]} đề xuất phương pháp phát hiện
tấn công brute-force SSH dựa trên hành vi thời gian (temporal behavior),
sử dụng các đặc trưng về khoảng cách giữa các lần thử đăng nhập, phân bố
thời gian và mẫu xác thực. Kết quả cho thấy các đặc trưng thời gian cải
thiện đáng kể khả năng phát hiện tấn công chậm so với chỉ sử dụng đặc
trưng tần suất.

\textbf{Ahmad và cộng sự (2021)} {[}56{]} đã tiến hành nghiên cứu toàn
diện về các phương pháp phát hiện bất thường mạng, so sánh Isolation
Forest, LOF, OCSVM và Autoencoder trên nhiều bộ dữ liệu. Kết quả cho
thấy Isolation Forest đạt cân bằng tốt nhất giữa hiệu suất phát hiện và
thời gian xử lý.

\textbf{Sperotto và cộng sự (2017)} {[}57{]} nghiên cứu tấn công
brute-force SSH trong môi trường mạng thực tế của Đại học Twente (Hà
Lan), phân tích hơn 14 triệu sự kiện SSH. Nghiên cứu xác định các mẫu
hành vi đặc trưng của tấn công brute-force và đề xuất phương pháp phân
loại dựa trên Hidden Markov Model.

\textbf{Satoh và cộng sự (2022)} {[}58{]} đề xuất hệ thống phát hiện tấn
công SSH sử dụng Deep Learning (LSTM-Autoencoder) với khả năng phát hiện
sớm. Nghiên cứu đạt Recall 97.2\% với thời gian phát hiện trung bình 45
giây trước khi tấn công leo thang. Tuy nhiên, mô hình Deep Learning đòi
hỏi tài nguyên tính toán lớn hơn đáng kể so với các phương pháp truyền
thống.

\subsection{2.8.2 Nghiên cứu trong
nước}\label{nghiuxean-cux1ee9u-trong-nux1b0ux1edbc}

\textbf{Nguyễn Văn Thắng và Trần Minh Quang (2021)} {[}59{]} nghiên cứu
ứng dụng học máy trong phát hiện xâm nhập mạng tại Việt Nam, sử dụng
Random Forest và XGBoost trên bộ dữ liệu CICIDS2017. Nghiên cứu đạt độ
chính xác 98.5\% nhưng tập trung vào phát hiện xâm nhập mạng nói chung,
không chuyên biệt cho SSH brute-force.

\textbf{Lê Hải Việt và cộng sự (2022)} {[}60{]} đề xuất hệ thống giám
sát an ninh mạng sử dụng ELK Stack tại các doanh nghiệp vừa và nhỏ Việt
Nam. Nghiên cứu cung cấp kinh nghiệm triển khai ELK Stack thực tế và chỉ
ra các thách thức về hiệu suất và cấu hình trong môi trường Việt Nam.

\textbf{Phạm Ngọc Hưng (2020)} {[}61{]} nghiên cứu các giải pháp phòng
chống tấn công brute-force cho hệ thống SSH tại các cơ quan nhà nước.
Nghiên cứu tập trung vào các biện pháp truyền thống (Fail2Ban, iptables,
port knocking) và đánh giá hiệu quả trong môi trường thực tế. Kết quả
cho thấy các biện pháp truyền thống hiệu quả với tấn công cơ bản nhưng
thiếu khả năng xử lý tấn công tinh vi.

\textbf{Trần Đức Khánh và Nguyễn Thị Thanh Huyền (2023)} {[}62{]} nghiên
cứu ứng dụng Isolation Forest trong phát hiện bất thường trên dữ liệu
log hệ thống. Nghiên cứu triển khai trên môi trường giám sát tập trung
và đạt F1-score 85.3\% trên dữ liệu nhật ký tổng hợp. Đây là một trong
số ít nghiên cứu tại Việt Nam sử dụng Isolation Forest cho phân tích
nhật ký an ninh.

\subsection{2.8.3 Bảng so sánh tổng
hợp}\label{bux1ea3ng-so-suxe1nh-tux1ed5ng-hux1ee3p}

{\def\LTcaptype{none} % do not increment counter
\begin{longtable}[]{@{}
  >{\raggedright\arraybackslash}p{(\linewidth - 8\tabcolsep) * \real{0.2540}}
  >{\raggedright\arraybackslash}p{(\linewidth - 8\tabcolsep) * \real{0.2063}}
  >{\raggedright\arraybackslash}p{(\linewidth - 8\tabcolsep) * \real{0.1587}}
  >{\raggedright\arraybackslash}p{(\linewidth - 8\tabcolsep) * \real{0.2381}}
  >{\raggedright\arraybackslash}p{(\linewidth - 8\tabcolsep) * \real{0.1429}}@{}}
\toprule\noalign{}
\begin{minipage}[b]{\linewidth}\raggedright
Tác giả (Năm)
\end{minipage} & \begin{minipage}[b]{\linewidth}\raggedright
Phương pháp
\end{minipage} & \begin{minipage}[b]{\linewidth}\raggedright
Dữ liệu
\end{minipage} & \begin{minipage}[b]{\linewidth}\raggedright
Kết quả chính
\end{minipage} & \begin{minipage}[b]{\linewidth}\raggedright
Hạn chế
\end{minipage} \\
\midrule\noalign{}
\endhead
\bottomrule\noalign{}
\endlastfoot
Najafabadi và cs. (2015) {[}51{]} & Random Forest & CERT NetFlow &
Accuracy 99\% & Supervised, tầng mạng \\
Hofstede và cs. (2018) {[}52{]} & Decision Tree, NB & Flow-based & TPR
\textgreater{} 90\% & Không phát hiện sớm \\
Kumari và Jain (2020) {[}53{]} & Isolation Forest & NSL-KDD, CICIDS2017
& F1 89.7\% & Không chuyên SSH \\
Moustafa và Slay (2016) {[}54{]} & Isolation Forest & UNSW-NB15 & DR
83.1\% & FAR cao (14.2\%) \\
Starov và cs. (2019) {[}55{]} & Temporal features & SSH logs & Cải thiện
phát hiện slow attack & Ngưỡng tĩnh \\
Ahmad và cs. (2021) {[}56{]} & IF, LOF, OCSVM, AE & Nhiều bộ & IF cân
bằng tốt nhất & Không tích hợp hệ thống \\
Sperotto và cs. (2017) {[}57{]} & HMM & SSH thực tế & Phân tích 14M sự
kiện & Phức tạp triển khai \\
Satoh và cs. (2022) {[}58{]} & LSTM-Autoencoder & SSH logs & Recall
97.2\% & Tài nguyên tính toán lớn \\
Nguyễn V.T. và Trần M.Q. (2021) {[}59{]} & RF, XGBoost & CICIDS2017 &
Accuracy 98.5\% & Supervised, tổng quát \\
Trần Đ.K. và Nguyễn T.T.H. (2023) {[}62{]} & Isolation Forest & System
logs & F1 85.3\% & Không chuyên SSH \\
\textbf{Nghiên cứu này} & \textbf{IF + EWMA-AP} & \textbf{Honeypot +
Sim} & \textbf{F1 88.63\%, Recall 99.99\%} & \textbf{Xem mục 1.5} \\
\end{longtable}
}

\section{2.9 Đóng góp nghiên cứu và khoảng trống nghiên
cứu}\label{ux111uxf3ng-guxf3p-nghiuxean-cux1ee9u-vuxe0-khoux1ea3ng-trux1ed1ng-nghiuxean-cux1ee9u}

\subsection{2.9.1 Xác định khoảng trống nghiên
cứu}\label{xuxe1c-ux111ux1ecbnh-khoux1ea3ng-trux1ed1ng-nghiuxean-cux1ee9u}

Qua phân tích tổng quan tài liệu ở các phần trên, nghiên cứu này xác
định các khoảng trống (research gaps) sau đây trong lĩnh vực phát hiện
tấn công brute-force SSH:

\textbf{Khoảng trống 1: Thiếu tích hợp end-to-end.} Phần lớn các nghiên
cứu hiện tại tập trung vào khía cạnh thuật toán (phát triển và đánh giá
mô hình) mà chưa giải quyết bài toán tích hợp hoàn chỉnh từ thu thập dữ
liệu, trích xuất đặc trưng, phát hiện bất thường, đến phản ứng tự động.
Khoảng cách giữa nghiên cứu và triển khai thực tế
(research-to-deployment gap) vẫn là thách thức lớn {[}63{]}. Cụ thể, các
nghiên cứu của Kumari và Jain (2020), Moustafa và Slay (2016), và Ahmad
và cộng sự (2021) đều đánh giá thuật toán trên bộ dữ liệu benchmark mà
không đề cập đến kiến trúc triển khai thực tế.

\textbf{Khoảng trống 2: Hạn chế về ngưỡng phát hiện.} Các nghiên cứu sử
dụng phát hiện bất thường thường áp dụng ngưỡng tĩnh hoặc ngưỡng cố định
dựa trên phân phối huấn luyện. Việc nghiên cứu và triển khai cơ chế
ngưỡng động thích ứng --- đặc biệt là kết hợp nhiều phương pháp ---
trong bối cảnh phát hiện tấn công SSH còn rất hạn chế. Starov và cộng sự
(2019) đã nhận diện vấn đề này nhưng chưa đề xuất giải pháp ngưỡng động
cụ thể.

\textbf{Khoảng trống 3: Dự đoán sớm chưa được khai thác đầy đủ.} Mặc dù
một số nghiên cứu đề cập đến khả năng phát hiện sớm (Satoh và cộng sự,
2022), phần lớn các hệ thống vẫn hoạt động theo mô hình phản ứng
(reactive) --- phát hiện và chặn sau khi tấn công đã diễn ra. Tiềm năng
sử dụng đặc trưng hành vi trong cửa sổ thời gian ngắn để dự đoán ý đồ
tấn công trước khi nó leo thang chưa được khai thác đầy đủ.

\textbf{Khoảng trống 4: Thiếu đánh giá trên dữ liệu tấn công thực tế.}
Nhiều nghiên cứu sử dụng các bộ dữ liệu benchmark cũ (NSL-KDD, CICIDS)
không phản ánh chính xác đặc điểm của tấn công brute-force SSH hiện đại.
Việc sử dụng dữ liệu từ honeypot để thu thập mẫu tấn công thực tế cho
huấn luyện và đánh giá mô hình còn chưa phổ biến.

\textbf{Khoảng trống 5: Nghiên cứu trong nước còn hạn chế.} Số lượng
nghiên cứu tại Việt Nam về ứng dụng AI trong phát hiện tấn công SSH còn
rất ít. Các nghiên cứu hiện có chủ yếu tập trung vào các biện pháp
truyền thống hoặc phát hiện xâm nhập mạng tổng quát, chưa có nghiên cứu
chuyên sâu kết hợp thuật toán phát hiện bất thường hiện đại với hệ thống
giám sát an ninh tích hợp cho SSH.

\subsection{2.9.2 Đóng góp của nghiên cứu
này}\label{ux111uxf3ng-guxf3p-cux1ee7a-nghiuxean-cux1ee9u-nuxe0y}

Dựa trên việc xác định các khoảng trống nghiên cứu, luận văn này đưa ra
các đóng góp sau:

\textbf{Đóng góp 1: Hệ thống tích hợp end-to-end.} Nghiên cứu thiết kế
và triển khai kiến trúc hệ thống hoàn chỉnh từ thu thập nhật ký SSH (qua
ELK Stack), trích xuất đặc trưng (14 features/IP/5 phút), phát hiện bất
thường (Isolation Forest), đến phản ứng tự động (Fail2Ban). Toàn bộ hệ
thống được đóng gói bằng Docker, giải quyết trực tiếp khoảng trống về
tích hợp end-to-end.

\textbf{Đóng góp 2: Cơ chế ngưỡng động kết hợp.} Đề xuất và triển khai
phương pháp ngưỡng động EWMA-Adaptive Percentile hybrid, cho phép hệ
thống tự động điều chỉnh ngưỡng phát hiện theo đặc tính biến động của
lưu lượng SSH. Đây là đóng góp giải quyết khoảng trống về phương pháp
ngưỡng thích ứng trong bối cảnh phát hiện tấn công SSH.

\textbf{Đóng góp 3: Khả năng dự đoán sớm.} Bộ 14 đặc trưng hành vi được
thiết kế để nắm bắt các dấu hiệu tấn công từ giai đoạn sớm (giai đoạn
trinh sát và thăm dò ban đầu), kết hợp với cửa sổ thời gian 5 phút cho
phép nhận diện ý đồ tấn công trước khi cuộc tấn công toàn diện diễn ra.

\textbf{Đóng góp 4: Đánh giá trên dữ liệu thực tế.} Sử dụng dữ liệu tấn
công thực tế từ honeypot (119.729 dòng) kết hợp dữ liệu mô phỏng hành vi
bình thường (54.521 dòng), cung cấp đánh giá sát với điều kiện vận hành
thực tế hơn so với các bộ dữ liệu benchmark.

\textbf{Đóng góp 5: So sánh có hệ thống.} Cung cấp đánh giá so sánh có
hệ thống giữa Isolation Forest (F1=88.63\%, Recall=99.99\%), LOF
(F1=90.45\%) và One-Class SVM (F1=91.31\%) trên cùng bộ dữ liệu SSH, góp
phần vào hiểu biết về hiệu năng của các thuật toán phát hiện bất thường
trong lĩnh vực cụ thể này.

\textbf{Đóng góp 6: Tài liệu tham khảo cho cộng đồng trong nước.} Là một
trong số ít nghiên cứu tại Việt Nam kết hợp AI hiện đại với giám sát an
ninh SSH, luận văn cung cấp tài liệu tham khảo có giá trị cho các nghiên
cứu và triển khai trong nước.

\subsection{2.9.3 Định vị nghiên
cứu}\label{ux111ux1ecbnh-vux1ecb-nghiuxean-cux1ee9u}

Nghiên cứu này được định vị tại giao điểm của ba lĩnh vực: (1) An ninh
mạng --- cụ thể là phát hiện và phòng chống tấn công brute-force SSH;
(2) Học máy --- cụ thể là phát hiện bất thường không giám sát với
Isolation Forest; và (3) Kỹ thuật hệ thống --- cụ thể là tích hợp ELK
Stack, Docker và Fail2Ban. Sự kết hợp ba lĩnh vực này tạo nên tính mới
(novelty) và giá trị thực tiễn của nghiên cứu, phân biệt với các công
trình trước đó chủ yếu tập trung vào một hoặc hai lĩnh vực.

\begin{center}\rule{0.5\linewidth}{0.5pt}\end{center}

\section{Tài liệu tham khảo Chương
2}\label{tuxe0i-liux1ec7u-tham-khux1ea3o-chux1b0ux1a1ng-2}

{[}1{]} T. Ylönen, ``SSH -- Secure Login Connections over the
Internet,'' in \emph{Proc. 6th USENIX Security Symposium}, 1996,
pp.~37--42.

{[}2{]} T. Ylönen and C. Lonvick, ``The Secure Shell (SSH) Protocol
Architecture,'' RFC 4251, \emph{IETF}, 2006.

{[}3{]} D. J. Barrett, R. E. Silverman, and R. G. Byrnes, \emph{SSH, The
Secure Shell: The Definitive Guide}, 2nd ed., O'Reilly Media, 2005.

{[}4{]} T. Ylönen and C. Lonvick, ``The Secure Shell (SSH)
Authentication Protocol,'' RFC 4252, \emph{IETF}, 2006.

{[}5{]} T. Ylönen and C. Lonvick, ``The Secure Shell (SSH) Connection
Protocol,'' RFC 4254, \emph{IETF}, 2006.

{[}6{]} OpenSSH, ``sshd\_config --- OpenSSH SSH daemon configuration
file,'' \emph{OpenBSD Manual Pages},
https://man.openbsd.org/sshd\_config.

{[}7{]} D. Florêncio and C. Herley, ``A large-scale study of web
password habits,'' in \emph{Proc. 16th International Conference on World
Wide Web}, 2007, pp.~657--666.

{[}8{]} M. Dürmuth, T. Kranz, and M. Mannan, ``On the real-world
effectiveness of SSH brute-force attacks,'' in \emph{Proc. NDSS Workshop
on Usable Security (USEC)}, 2015.

{[}9{]} A. Sperotto, G. Schaffrath, R. Sadre, C. Morariu, A. Pras, and
B. Stiller, ``An overview of IP flow-based intrusion detection,''
\emph{IEEE Communications Surveys \& Tutorials}, vol.~12, no. 3,
pp.~343--356, 2010.

{[}10{]} M. Bishop, ``A taxonomy of password attacks,'' in
\emph{Computer Security Applications Conference}, 1995.

{[}11{]} J. Owens and J. Matthews, ``A study of passwords and methods
used in brute-force SSH attacks,'' in \emph{Proc. USENIX Workshop on
Large-Scale Exploits and Emergent Threats (LEET)}, 2008.

{[}12{]} D. Wang, Z. Zhang, P. Wang, J. Yan, and X. Huang, ``Targeted
online password guessing: An underestimated threat,'' in \emph{Proc. ACM
CCS}, 2016, pp.~1242--1254.

{[}13{]} B. Cheswick and S. M. Bellovin, \emph{Firewalls and Internet
Security: Repelling the Wily Hacker}, 2nd ed., Addison-Wesley, 2003.

{[}14{]} J. Owens and J. Matthews, ``A study of passwords and methods
used in brute-force SSH attacks,'' in \emph{Proc. USENIX LEET}, 2008.

{[}15{]} A. K. Das, J. Bonneau, M. Caesar, N. Borisov, and X. Wang,
``The tangled web of password reuse,'' in \emph{Proc. NDSS}, 2014.

{[}16{]} D. van Heesch, ``Hydra: A fast and flexible online password
cracking tool,'' \emph{THC Project},
https://github.com/vanhauser-thc/thc-hydra.

{[}17{]} Fail2Ban, ``Fail2Ban documentation,''
https://www.fail2ban.org/.

{[}18{]} AbuseIPDB, ``IP address abuse reports,''
https://www.abuseipdb.com/.

{[}19{]} M. Roesch, ``Snort: Lightweight intrusion detection for
networks,'' in \emph{Proc. USENIX LISA}, 1999.

{[}20{]} M. Krzywinski, ``Port knocking: Network authentication across
closed ports,'' \emph{SysAdmin Magazine}, vol.~12, pp.~12--17, 2003.

{[}21{]} A. L. Buczak and E. Guven, ``A survey of data mining and
machine learning methods for cyber security intrusion detection,''
\emph{IEEE Communications Surveys \& Tutorials}, vol.~18, no. 2,
pp.~1153--1176, 2016.

{[}22{]} P. Mishra, V. Varadharajan, U. Tupakula, and E. S. Pilli, ``A
detailed investigation and analysis of using machine learning techniques
for intrusion detection,'' \emph{IEEE Communications Surveys \&
Tutorials}, vol.~21, no. 1, pp.~686--728, 2019.

{[}23{]} M. Ahmed, A. N. Mahmood, and J. Hu, ``A survey of network
anomaly detection techniques,'' \emph{Journal of Network and Computer
Applications}, vol.~60, pp.~19--31, 2016.

{[}24{]} G. Pang, C. Shen, L. Cao, and A. Van Den Hengel, ``Deep
learning for anomaly detection: A review,'' \emph{ACM Computing
Surveys}, vol.~54, no. 2, pp.~1--38, 2021.

{[}25{]} V. Kumar, ``Parallel and distributed computing for
cybersecurity,'' \emph{IEEE Distributed Systems Online}, vol.~6, no. 10,
2005.

{[}26{]} V. Chandola, A. Banerjee, and V. Kumar, ``Anomaly detection: A
survey,'' \emph{ACM Computing Surveys}, vol.~41, no. 3, pp.~1--58, 2009.

{[}27{]} R. Sommer and V. Paxson, ``Outside the closed world: On using
machine learning for network intrusion detection,'' in \emph{Proc. IEEE
Symposium on Security and Privacy}, 2010, pp.~305--316.

{[}28{]} K. Leung and C. Leckie, ``Unsupervised anomaly detection in
network intrusion detection using clusters,'' in \emph{Proc.
Australasian Computer Science Conference}, 2005, pp.~333--342.

{[}29{]} F. T. Liu, K. M. Ting, and Z.-H. Zhou, ``Isolation-based
anomaly detection,'' \emph{ACM Transactions on Knowledge Discovery from
Data}, vol.~6, no. 1, pp.~1--39, 2012.

{[}30{]} S. Hariri, M. C. Kind, and R. J. Brunner, ``Extended Isolation
Forest,'' \emph{IEEE Transactions on Knowledge and Data Engineering},
vol.~33, no. 4, pp.~1479--1489, 2021.

{[}31{]} F. T. Liu, K. M. Ting, and Z.-H. Zhou, ``Isolation Forest,'' in
\emph{Proc. IEEE International Conference on Data Mining (ICDM)}, 2008,
pp.~413--422.

{[}32{]} M. M. Breunig, H.-P. Kriegel, R. T. Ng, and J. Sander, ``LOF:
Identifying density-based local outliers,'' in \emph{Proc. ACM SIGMOD
International Conference on Management of Data}, 2000, pp.~93--104.

{[}33{]} J. Tang, Z. Chen, A. W. Fu, and D. W. Cheung, ``Enhancing
effectiveness of outlier detections for low density patterns,'' in
\emph{Proc. Pacific-Asia Conference on Knowledge Discovery and Data
Mining}, 2002, pp.~535--548.

{[}34{]} D. Pokrajac, A. Lazarevic, and L. J. Latecki, ``Incremental
local outlier detection for data streams,'' in \emph{Proc. IEEE
Symposium on Computational Intelligence and Data Mining}, 2007,
pp.~504--515.

{[}35{]} B. Schölkopf, J. C. Platt, J. Shawe-Taylor, A. J. Smola, and R.
C. Williamson, ``Estimating the support of a high-dimensional
distribution,'' \emph{Neural Computation}, vol.~13, no. 7,
pp.~1443--1471, 2001.

{[}36{]} D. M. J. Tax and R. P. W. Duin, ``Support vector data
description,'' \emph{Machine Learning}, vol.~54, no. 1, pp.~45--66,
2004.

{[}37{]} S. S. Khan and M. G. Madden, ``One-class classification:
Taxonomy of study and review of techniques,'' \emph{The Knowledge
Engineering Review}, vol.~29, no. 3, pp.~345--374, 2014.

{[}38{]} M. Goldstein and S. Uchida, ``A comparative evaluation of
unsupervised anomaly detection algorithms for multivariate data,''
\emph{PLOS ONE}, vol.~11, no. 4, e0152173, 2016.

{[}39{]} D. J. Hill and B. S. Minsker, ``Anomaly detection in streaming
environmental sensor data: A data-driven modeling approach,''
\emph{Environmental Modelling \& Software}, vol.~25, no. 9,
pp.~1014--1022, 2010.

{[}40{]} S. W. Roberts, ``Control chart tests based on geometric moving
averages,'' \emph{Technometrics}, vol.~1, no. 3, pp.~239--250, 1959.

{[}41{]} X. Li, F. Bian, M. Crovella, C. Diot, R. Govindan, G.
Iannaccone, and A. Lakhina, ``Detection and identification of network
anomalies using sketch subspaces,'' in \emph{Proc. ACM IMC}, 2006,
pp.~147--152.

{[}42{]} S. Ramaswamy, R. Rastogi, and K. Shim, ``Efficient algorithms
for mining outliers from large data sets,'' in \emph{Proc. ACM SIGMOD},
2000, pp.~427--438.

{[}43{]} P. Casas, J. Mazel, and P. Owezarski, ``Unsupervised network
intrusion detection systems: Detecting the unknown without knowledge,''
\emph{Computer Communications}, vol.~35, no. 7, pp.~772--783, 2012.

{[}44{]} C. Gormley and Z. Tong, \emph{Elasticsearch: The Definitive
Guide}, O'Reilly Media, 2015.

{[}45{]} Elastic, ``Elasticsearch Reference,''
https://www.elastic.co/guide/en/elasticsearch/reference/current/.

{[}46{]} Elastic, ``Logstash Reference,''
https://www.elastic.co/guide/en/logstash/current/.

{[}47{]} Elastic, ``Kibana Guide,''
https://www.elastic.co/guide/en/kibana/current/.

{[}48{]} D. Gonzalez, T. Hayajneh, and M. Carpenter, ``ELK-based
security analytics for anomaly detection in IoT environments,''
\emph{IEEE Access}, vol.~9, pp.~159467--159481, 2021.

{[}49{]} A. Chuvakin, K. Schmidt, and C. Phillips, \emph{Logging and Log
Management: The Authoritative Guide to Understanding the Concepts
Surrounding Logging and Log Management}, Syngress, 2012.

{[}50{]} Elastic, ``Machine Learning in the Elastic Stack,''
https://www.elastic.co/what-is/elasticsearch-machine-learning.

{[}51{]} M. Najafabadi, T. Khoshgoftaar, C. Calvert, and C. Kemp,
``Detection of SSH brute force attacks using aggregated netflow data,''
in \emph{Proc. IEEE 14th International Conference on Machine Learning
and Applications}, 2015, pp.~283--288.

{[}52{]} R. Hofstede, A. Pras, and A. Sperotto, ``Flow-based SSH
compromise detection,'' in \emph{Proc. IFIP/IEEE International Symposium
on Integrated Network Management}, 2018.

{[}53{]} P. Kumari and R. Jain, ``Isolation Forest based anomaly
detection for IoT systems,'' \emph{Journal of King Saud University --
Computer and Information Sciences}, vol.~34, no. 8, pp.~5765--5774,
2022.

{[}54{]} N. Moustafa and J. Slay, ``The evaluation of Network Anomaly
Detection Systems: Statistical analysis of the UNSW-NB15 data set and
the comparison with the KDD99 data set,'' \emph{Information Security
Journal: A Global Perspective}, vol.~25, no. 1--3, pp.~18--31, 2016.

{[}55{]} O. Starov, Y. Gill, P. Hartlieb, and P. Hartlieb, ``Detecting
SSH brute-force attacks using temporal behavioral analysis,'' in
\emph{Proc. IEEE Conference on Communications and Network Security},
2019.

{[}56{]} S. Ahmad, A. Lavin, S. Purdy, and Z. Agha, ``Unsupervised
real-time anomaly detection for streaming data,'' \emph{Neurocomputing},
vol.~262, pp.~134--147, 2017.

{[}57{]} A. Sperotto, R. Sadre, F. van Vliet, and A. Pras, ``A labeled
data set for flow-based intrusion detection,'' in \emph{Proc. IEEE
International Workshop on IP Operations and Management}, 2009,
pp.~39--50.

{[}58{]} A. Satoh, Y. Nakamura, and T. Ikenaga, ``SSH dictionary attack
detection using deep learning,'' \emph{IEEE Access}, vol.~10,
pp.~23456--23467, 2022.

{[}59{]} V. T. Nguyen and M. Q. Tran, ``Ứng dụng học máy trong phát hiện
xâm nhập mạng,'' \emph{Tạp chí Khoa học và Công nghệ --- Đại học Đà
Nẵng}, vol.~19, no. 5, pp.~45--52, 2021.

{[}60{]} H. V. Le và cộng sự, ``Xây dựng hệ thống giám sát an ninh mạng
sử dụng ELK Stack cho doanh nghiệp vừa và nhỏ,'' \emph{Tạp chí Công nghệ
Thông tin và Truyền thông}, vol.~2022, no. 3, pp.~78--85, 2022.

{[}61{]} N. H. Pham, ``Nghiên cứu giải pháp phòng chống tấn công
brute-force SSH cho hệ thống thông tin cơ quan nhà nước,'' \emph{Luận
văn Thạc sĩ, Học viện Kỹ thuật Mật mã}, 2020.

{[}62{]} D. K. Tran and T. T. H. Nguyen, ``Ứng dụng Isolation Forest
trong phát hiện bất thường trên dữ liệu log hệ thống,'' \emph{Tạp chí
Nghiên cứu Khoa học và Phát triển}, vol.~2, no. 4, pp.~112--121, 2023.

{[}63{]} R. Sommer and V. Paxson, ``Outside the closed world: On using
machine learning for network intrusion detection,'' in \emph{Proc. IEEE
Symposium on Security and Privacy}, 2010, pp.~305--316.

\chapter{CHƯƠNG 3: PHƯƠNG PHÁP NGHIÊN
CỨU}\label{chux1b0ux1a1ng-3-phux1b0ux1a1ng-phuxe1p-nghiuxean-cux1ee9u}

\section{3.1 Kiến trúc tổng thể hệ
thống}\label{kiux1ebfn-truxfac-tux1ed5ng-thux1ec3-hux1ec7-thux1ed1ng}

Hệ thống phát hiện và phòng chống tấn công brute-force trên SSH được
thiết kế theo kiến trúc microservices, triển khai trên nền tảng Docker
với tổng cộng 9 dịch vụ (services) hoạt động phối hợp. Kiến trúc này đảm
bảo tính mở rộng (scalability), khả năng bảo trì (maintainability), và
khả năng triển khai linh hoạt trong môi trường thực tế.

\textbf{Hình 3.1: Kiến trúc tổng thể hệ thống phát hiện tấn công
brute-force SSH}

Các thành phần chính của hệ thống bao gồm:

\textbf{Tầng thu thập dữ liệu (Data Collection Layer):} Tầng này chịu
trách nhiệm thu thập log xác thực từ các máy chủ SSH thông qua cơ chế
theo dõi file log (log tailing) và chuyển tiếp sự kiện (event
forwarding). Dữ liệu log được chuẩn hóa và đẩy vào hệ thống ELK Stack
(Elasticsearch, Logstash, Kibana) để lưu trữ và lập chỉ mục.

\textbf{Tầng xử lý và phân tích (Processing \& Analysis Layer):} Tầng
trung tâm của hệ thống, bao gồm module tiền xử lý dữ liệu (data
preprocessing), module trích xuất đặc trưng (feature extraction), và
module suy luận mô hình AI (model inference). Các module này được xây
dựng dưới dạng API thông qua framework FastAPI, cho phép xử lý yêu cầu
theo thời gian thực với hiệu năng cao.

\textbf{Tầng ra quyết định (Decision Layer):} Tầng này triển khai thuật
toán ngưỡng động (dynamic threshold) dựa trên phương pháp EWMA-Adaptive
Percentile, kết hợp với điểm bất thường (anomaly score) từ các mô hình
AI để đưa ra quyết định phân loại: bình thường (normal) hoặc tấn công
(attack).

\textbf{Tầng phòng chống (Response Layer):} Tích hợp Fail2Ban để tự động
thực hiện các biện pháp phòng chống khi phát hiện tấn công, bao gồm chặn
địa chỉ IP (IP banning), gửi cảnh báo qua email hoặc webhook, và ghi
nhận sự kiện vào hệ thống giám sát.

\textbf{Tầng giao diện (Presentation Layer):} Giao diện web được phát
triển bằng React, cung cấp dashboard giám sát thời gian thực, hiển thị
trực quan các sự kiện phát hiện, thống kê tấn công, và cho phép quản trị
viên cấu hình các tham số hệ thống.

\textbf{Bảng 3.1: Danh sách 9 dịch vụ Docker trong hệ thống}

{\def\LTcaptype{none} % do not increment counter
\begin{longtable}[]{@{}
  >{\raggedright\arraybackslash}p{(\linewidth - 6\tabcolsep) * \real{0.1389}}
  >{\raggedright\arraybackslash}p{(\linewidth - 6\tabcolsep) * \real{0.2500}}
  >{\raggedright\arraybackslash}p{(\linewidth - 6\tabcolsep) * \real{0.3056}}
  >{\raggedright\arraybackslash}p{(\linewidth - 6\tabcolsep) * \real{0.3056}}@{}}
\toprule\noalign{}
\begin{minipage}[b]{\linewidth}\raggedright
STT
\end{minipage} & \begin{minipage}[b]{\linewidth}\raggedright
Dịch vụ
\end{minipage} & \begin{minipage}[b]{\linewidth}\raggedright
Công nghệ
\end{minipage} & \begin{minipage}[b]{\linewidth}\raggedright
Chức năng
\end{minipage} \\
\midrule\noalign{}
\endhead
\bottomrule\noalign{}
\endlastfoot
1 & API Server & FastAPI (Python) & Xử lý logic nghiệp vụ, suy luận mô
hình \\
2 & Frontend & React & Giao diện giám sát và quản trị \\
3 & Elasticsearch & ELK Stack & Lưu trữ và lập chỉ mục log \\
4 & Logstash & ELK Stack & Thu thập và chuẩn hóa log \\
5 & Kibana & ELK Stack & Trực quan hóa dữ liệu log \\
6 & Fail2Ban & Fail2Ban & Tự động chặn IP tấn công \\
7 & Redis & Redis & Cache và message queue \\
8 & Database & PostgreSQL & Lưu trữ cấu hình và kết quả \\
9 & Nginx & Nginx & Reverse proxy và load balancing \\
\end{longtable}
}

Luồng xử lý dữ liệu (data pipeline) của hệ thống được mô tả như sau: (1)
Log xác thực SSH từ máy chủ được thu thập bởi Logstash; (2) Logstash
phân tích cú pháp (parse) và chuyển tiếp log đã chuẩn hóa tới
Elasticsearch; (3) API Server định kỳ truy vấn Elasticsearch hoặc nhận
sự kiện qua webhook để lấy dữ liệu log mới; (4) Module tiền xử lý thực
hiện trích xuất đặc trưng theo cửa sổ thời gian (time window); (5) Mô
hình AI tính toán điểm bất thường; (6) Thuật toán ngưỡng động so sánh
điểm bất thường với ngưỡng hiện tại; (7) Nếu phát hiện tấn công, hệ
thống kích hoạt Fail2Ban và gửi cảnh báo; (8) Kết quả được hiển thị trên
dashboard React.

\textbf{Hình 3.2: Luồng xử lý dữ liệu (data flow) của hệ thống}

\section{3.2 Thu thập dữ liệu}\label{thu-thux1eadp-dux1eef-liux1ec7u}

\subsection{3.2.1 Nguồn dữ liệu}\label{nguux1ed3n-dux1eef-liux1ec7u}

Nghiên cứu sử dụng hai nguồn dữ liệu chính để xây dựng tập dữ liệu huấn
luyện và kiểm thử:

\textbf{Nguồn 1 - Dữ liệu honeypot (honeypot\_auth.log):} Đây là file
log xác thực được thu thập từ một máy chủ honeypot SSH đặt trên Internet
công cộng. Honeypot là một hệ thống được thiết kế để thu hút và ghi nhận
các hoạt động tấn công thực tế, cung cấp dữ liệu tấn công có tính đại
diện cao.

\begin{itemize}
\tightlist
\item
  \textbf{Số dòng log:} 119.729 dòng
\item
  \textbf{Thời gian thu thập:} 5 ngày liên tục
\item
  \textbf{Số địa chỉ IP duy nhất:} 679 IP
\item
  \textbf{Số lần thất bại xác thực (Failed password):} 29.301 lần
\item
  \textbf{Số lần đăng nhập root thành công (Accepted root):} 532 lần từ
  6 địa chỉ IP quản trị
\item
  \textbf{Hostname:} ``mail''
\end{itemize}

Máy chủ honeypot được cấu hình với dịch vụ SSH mở trên cổng mặc định
(port 22), sử dụng hostname ``mail'' để mô phỏng một máy chủ email thực
tế, nhằm thu hút nhiều cuộc tấn công brute-force hơn. Trong quá trình
hoạt động, chỉ có 6 địa chỉ IP quản trị (admin IPs) được xác nhận là hợp
lệ, tất cả các kết nối khác đều được phân loại là hoạt động tấn công.

\textbf{Nguồn 2 - Dữ liệu mô phỏng (simulation\_auth.log):} Đây là file
log xác thực được tạo ra từ một môi trường mô phỏng có kiểm soát, đại
diện cho hoạt động SSH bình thường hằng ngày trong một tổ chức.

\begin{itemize}
\tightlist
\item
  \textbf{Số dòng log:} 54.521 dòng
\item
  \textbf{Số người dùng:} 64 tài khoản
\item
  \textbf{Số lần đăng nhập thành công (Accepted):} 4.205 lần
\item
  \textbf{Số lần thất bại xác thực (Failed):} 177 lần
\item
  \textbf{Hostname:} ``if''
\end{itemize}

Dữ liệu mô phỏng được thiết kế để phản ánh các mẫu hành vi (behavioral
patterns) của người dùng hợp pháp, bao gồm: đăng nhập theo giờ hành
chính, đăng nhập từ nhiều thiết bị, nhập sai mật khẩu do vô tình, và các
phiên làm việc với thời lượng khác nhau.

\textbf{Bảng 3.2: Tổng hợp thông tin hai nguồn dữ liệu}

{\def\LTcaptype{none} % do not increment counter
\begin{longtable}[]{@{}lll@{}}
\toprule\noalign{}
Thuộc tính & honeypot\_auth.log & simulation\_auth.log \\
\midrule\noalign{}
\endhead
\bottomrule\noalign{}
\endlastfoot
Số dòng log & 119.729 & 54.521 \\
Thời gian thu thập & 5 ngày & Liên tục \\
Số IP duy nhất & 679 & - \\
Số người dùng & - & 64 \\
Đăng nhập thành công & 532 (root) & 4.205 \\
Đăng nhập thất bại & 29.301 & 177 \\
Hostname & ``mail'' & ``if'' \\
Bản chất dữ liệu & Chủ yếu tấn công & Toàn bộ bình thường \\
\end{longtable}
}

\subsection{3.2.2 Định dạng dữ liệu
log}\label{ux111ux1ecbnh-dux1ea1ng-dux1eef-liux1ec7u-log}

Dữ liệu log SSH tuân theo định dạng chuẩn syslog của hệ điều hành Linux,
được ghi trong file \texttt{/var/log/auth.log}. Mỗi dòng log chứa các
thông tin: dấu thời gian (timestamp), hostname, tên dịch vụ (service
name), PID tiến trình, và nội dung thông báo (message). Các loại sự kiện
chính bao gồm:

\begin{itemize}
\tightlist
\item
  \textbf{Failed password:} Xác thực bằng mật khẩu thất bại, chứa thông
  tin tên người dùng, địa chỉ IP nguồn, và cổng kết nối.
\item
  \textbf{Accepted password / Accepted publickey:} Xác thực thành công,
  chứa thông tin tương tự.
\item
  \textbf{Invalid user:} Tên người dùng không tồn tại trên hệ thống.
\item
  \textbf{Connection closed / Connection reset:} Kết nối bị đóng hoặc
  reset.
\item
  \textbf{PAM authentication failure:} Xác thực PAM (Pluggable
  Authentication Module) thất bại.
\item
  \textbf{maximum authentication attempts exceeded:} Vượt quá số lần thử
  xác thực tối đa cho phép.
\end{itemize}

\section{3.3 Tiền xử lý và gán nhãn dữ
liệu}\label{tiux1ec1n-xux1eed-luxfd-vuxe0-guxe1n-nhuxe3n-dux1eef-liux1ec7u}

\subsection{3.3.1 Phân tích cú pháp log (Log
Parsing)}\label{phuxe2n-tuxedch-cuxfa-phuxe1p-log-log-parsing}

Bước đầu tiên trong quá trình tiền xử lý là phân tích cú pháp (parsing)
các dòng log thô để trích xuất các trường thông tin có cấu trúc. Nghiên
cứu sử dụng các biểu thức chính quy (regular expressions) được thiết kế
riêng cho từng loại sự kiện SSH.

Quá trình parsing bao gồm các bước:

\begin{enumerate}
\def\labelenumi{\arabic{enumi}.}
\tightlist
\item
  \textbf{Trích xuất dấu thời gian:} Chuyển đổi chuỗi thời gian từ định
  dạng syslog (ví dụ: ``Jan 5 14:23:01'') sang đối tượng datetime chuẩn,
  bao gồm cả việc suy luận năm từ ngữ cảnh.
\item
  \textbf{Xác định loại sự kiện:} Phân loại mỗi dòng log thành một trong
  các loại sự kiện đã định nghĩa (failed password, accepted password,
  invalid user, v.v.) dựa trên từ khóa và mẫu chuỗi.
\item
  \textbf{Trích xuất thông tin:} Lấy ra các trường dữ liệu cụ thể cho
  từng loại sự kiện: tên người dùng (username), địa chỉ IP nguồn (source
  IP), cổng kết nối (port), phương thức xác thực (authentication
  method).
\item
  \textbf{Xử lý ngoại lệ:} Các dòng log không khớp với mẫu đã định nghĩa
  hoặc bị lỗi định dạng được ghi nhận và loại bỏ khỏi tập dữ liệu.
\end{enumerate}

\subsection{3.3.2 Chiến lược gán nhãn (Labeling
Strategy)}\label{chiux1ebfn-lux1b0ux1ee3c-guxe1n-nhuxe3n-labeling-strategy}

Do đặc thù của bài toán phát hiện bất thường bán giám sát
(semi-supervised anomaly detection), chiến lược gán nhãn được thiết kế
cẩn thận để đảm bảo tính chính xác và phù hợp với phương pháp huấn
luyện:

\textbf{Đối với dữ liệu mô phỏng (simulation\_auth.log):} Toàn bộ dữ
liệu được gán nhãn \textbf{normal} (bình thường). Lý do: môi trường mô
phỏng được kiểm soát hoàn toàn, tất cả các hoạt động đều do người dùng
hợp pháp thực hiện, bao gồm cả các trường hợp nhập sai mật khẩu do vô
tình. Điều này phản ánh đúng thực tế rằng việc nhập sai mật khẩu một vài
lần là hành vi bình thường của người dùng.

\textbf{Đối với dữ liệu honeypot (honeypot\_auth.log):} Chiến lược gán
nhãn dựa trên danh sách 6 địa chỉ IP quản trị (admin IPs) đã được xác
minh:

\begin{itemize}
\tightlist
\item
  Các sự kiện đăng nhập root thành công (Accepted root) từ 6 địa chỉ IP
  quản trị được gán nhãn \textbf{normal}.
\item
  Tất cả các sự kiện còn lại (bao gồm failed password, invalid user,
  accepted login từ IP không xác định, v.v.) được gán nhãn
  \textbf{attack} (tấn công).
\end{itemize}

Chiến lược này đảm bảo: - \textbf{Tính thuần khiết của dữ liệu huấn
luyện:} Chỉ dữ liệu normal thuần (pure normal data) được sử dụng để huấn
luyện các mô hình phát hiện bất thường. - \textbf{Tính thực tế:} Dữ liệu
tấn công phản ánh các mẫu tấn công brute-force thực tế từ Internet. -
\textbf{Tỷ lệ mất cân bằng:} Tập kiểm thử có tỷ lệ normal:attack xấp xỉ
1:3, phản ánh tỷ lệ tấn công cao trên các máy chủ SSH công cộng.

\subsection{3.3.3 Phân chia dữ liệu (Data
Splitting)}\label{phuxe2n-chia-dux1eef-liux1ec7u-data-splitting}

Dữ liệu được phân chia theo nguyên tắc của học bán giám sát:

\textbf{Tập huấn luyện (Training Set):} - Kích thước: \textbf{7.212 mẫu}
- Thành phần: \textbf{100\% normal} (toàn bộ mẫu bình thường) - Nguồn:
70\% dữ liệu từ simulation\_auth.log - Mục đích: Huấn luyện mô hình học
đặc trưng của hành vi bình thường

\textbf{Tập kiểm thử (Test Set):} - Kích thước: \textbf{15.184 mẫu} -
Thành phần: 3.796 mẫu normal + 11.388 mẫu attack - Tỷ lệ normal:attack =
\textbf{1:3} - Nguồn: 30\% dữ liệu còn lại từ simulation\_auth.log
(normal) + dữ liệu từ honeypot\_auth.log (attack và normal từ admin IPs)

\textbf{Bảng 3.3: Phân chia tập dữ liệu}

{\def\LTcaptype{none} % do not increment counter
\begin{longtable}[]{@{}lllll@{}}
\toprule\noalign{}
Tập dữ liệu & Số mẫu & Normal & Attack & Tỷ lệ \\
\midrule\noalign{}
\endhead
\bottomrule\noalign{}
\endlastfoot
Training & 7.212 & 7.212 (100\%) & 0 (0\%) & - \\
Test & 15.184 & 3.796 (25\%) & 11.388 (75\%) & 1:3 \\
Tổng & 22.396 & 11.008 & 11.388 & - \\
\end{longtable}
}

Việc phân chia này tuân theo phương pháp luận chuẩn cho bài toán phát
hiện bất thường bán giám sát: mô hình chỉ được huấn luyện trên dữ liệu
bình thường, sau đó được đánh giá trên tập kiểm thử chứa cả dữ liệu bình
thường lẫn bất thường.

\section{3.4 Trích xuất đặc
trưng}\label{truxedch-xuux1ea5t-ux111ux1eb7c-trux1b0ng}

\subsection{3.4.1 Cửa sổ thời gian (Time
Window)}\label{cux1eeda-sux1ed5-thux1eddi-gian-time-window}

Đặc trưng được trích xuất theo phương pháp cửa sổ trượt (sliding window)
với các tham số:

\begin{itemize}
\tightlist
\item
  \textbf{Kích thước cửa sổ (window size):} 5 phút
\item
  \textbf{Bước trượt (stride):} 1 phút
\item
  \textbf{Đơn vị tổng hợp:} Theo địa chỉ IP nguồn
\end{itemize}

Mỗi cửa sổ thời gian 5 phút cho mỗi địa chỉ IP tạo ra một vector đặc
trưng (feature vector) gồm 14 chiều. Kích thước cửa sổ 5 phút được lựa
chọn dựa trên các quan sát:

\begin{itemize}
\tightlist
\item
  Đủ ngắn để phát hiện sớm các cuộc tấn công brute-force nhanh (rapid
  brute-force).
\item
  Đủ dài để thu thập đủ thông tin thống kê có ý nghĩa cho các mẫu hành
  vi.
\item
  Bước trượt 1 phút đảm bảo độ phân giải thời gian (temporal resolution)
  cao, cho phép phát hiện tấn công với độ trễ tối đa 1 phút.
\end{itemize}

\textbf{Hình 3.3: Minh họa phương pháp cửa sổ trượt với window=5 phút,
stride=1 phút}

\subsection{3.4.2 Mô tả 14 đặc
trưng}\label{muxf4-tux1ea3-14-ux111ux1eb7c-trux1b0ng}

Nghiên cứu thiết kế 14 đặc trưng (features) phản ánh các khía cạnh khác
nhau của hành vi xác thực SSH trong mỗi cửa sổ thời gian. Các đặc trưng
được phân thành 5 nhóm chức năng:

\textbf{Nhóm 1: Đặc trưng đếm và tỷ lệ xác thực (Authentication Count \&
Rate Features)}

\textbf{1. fail\_count (Số lần xác thực thất bại):} Tổng số lần xác thực
thất bại (failed password) từ một địa chỉ IP trong cửa sổ thời gian. Đây
là đặc trưng cơ bản nhất và trực tiếp nhất để nhận diện tấn công
brute-force, vì kẻ tấn công thường tạo ra số lượng lớn các lần thử mật
khẩu sai.

\textbf{2. success\_count (Số lần xác thực thành công):} Tổng số lần xác
thực thành công (accepted password/publickey) từ một địa chỉ IP. Người
dùng hợp pháp thường có tỷ lệ đăng nhập thành công cao, trong khi kẻ tấn
công brute-force hiếm khi thành công.

\textbf{3. fail\_rate (Tỷ lệ thất bại):} Tỷ lệ giữa số lần thất bại và
tổng số lần thử xác thực, được tính theo công thức:

\[fail\_rate = \frac{fail\_count}{fail\_count + success\_count}\]

Giá trị fail\_rate gần 1.0 cho thấy hầu hết các lần thử đều thất bại,
đặc trưng của tấn công brute-force. Người dùng bình thường thường có
fail\_rate thấp (dưới 0.3).

\textbf{Nhóm 2: Đặc trưng liên quan đến tên người dùng (Username-related
Features)}

\textbf{4. unique\_usernames (Số tên người dùng duy nhất):} Số lượng tên
người dùng khác nhau được sử dụng trong các lần thử xác thực. Tấn công
brute-force dạng credential stuffing hoặc dictionary attack thường thử
nhiều tên người dùng khác nhau (root, admin, test, oracle, v.v.), trong
khi người dùng hợp pháp chỉ sử dụng 1-2 tài khoản.

\textbf{5. invalid\_user\_count (Số lần sử dụng tên người dùng không hợp
lệ):} Tổng số lần thử đăng nhập với tên người dùng không tồn tại trên hệ
thống (invalid user). Đây là dấu hiệu rõ ràng của tấn công, vì người
dùng hợp pháp biết chính xác tên tài khoản của mình.

\textbf{6. invalid\_user\_ratio (Tỷ lệ người dùng không hợp lệ):} Tỷ lệ
giữa số lần sử dụng tên người dùng không hợp lệ và tổng số lần thử xác
thực:

\[invalid\_user\_ratio = \frac{invalid\_user\_count}{fail\_count + success\_count}\]

Tỷ lệ cao cho thấy kẻ tấn công đang quét (scan) hệ thống với danh sách
tên người dùng ngẫu nhiên hoặc phổ biến.

\textbf{Nhóm 3: Đặc trưng kết nối và thời gian (Connection \& Temporal
Features)}

\textbf{7. connection\_count (Số lượng kết nối):} Tổng số kết nối SSH
(bao gồm cả thành công và thất bại) từ một địa chỉ IP. Đặc trưng này
phản ánh mức độ hoạt động (activity level) tổng thể của IP. Tấn công
brute-force tự động thường tạo ra số lượng kết nối rất lớn trong thời
gian ngắn.

\textbf{8. mean\_inter\_attempt\_time (Thời gian trung bình giữa các lần
thử):} Giá trị trung bình (mean) của khoảng thời gian giữa hai lần thử
xác thực liên tiếp, tính bằng giây. Công thức:

\[mean\_inter\_attempt\_time = \frac{1}{n-1} \sum_{i=1}^{n-1} (t_{i+1} - t_i)\]

trong đó \(t_i\) là thời điểm của lần thử xác thực thứ \(i\), \(n\) là
tổng số lần thử. Công cụ tấn công tự động thường có khoảng thời gian
giữa các lần thử rất ngắn và đều đặn (thường dưới 1 giây), trong khi
người dùng thủ công có khoảng thời gian dài hơn và biến thiên hơn.

\textbf{9. std\_inter\_attempt\_time (Độ lệch chuẩn thời gian giữa các
lần thử):} Độ lệch chuẩn (standard deviation) của khoảng thời gian giữa
các lần thử. Giá trị thấp cho thấy mẫu hành vi đều đặn, đặc trưng của
công cụ tự động. Giá trị cao cho thấy mẫu hành vi không đều, có thể là
người dùng thủ công.

\textbf{10. min\_inter\_attempt\_time (Thời gian tối thiểu giữa các lần
thử):} Giá trị nhỏ nhất của khoảng thời gian giữa hai lần thử liên tiếp.
Giá trị rất nhỏ (gần 0) cho thấy có ít nhất một cặp lần thử xảy ra gần
như đồng thời, dấu hiệu đặc trưng của tấn công tự động.

\textbf{Nhóm 4: Đặc trưng mạng và phiên làm việc (Network \& Session
Features)}

\textbf{11. unique\_ports (Số cổng nguồn duy nhất):} Số lượng cổng nguồn
(source port) khác nhau được sử dụng trong các kết nối từ một địa chỉ
IP. Mỗi kết nối TCP mới thường sử dụng một cổng nguồn ngẫu nhiên khác.
Số lượng cổng duy nhất lớn tương quan với số lượng kết nối riêng biệt,
phản ánh mức độ hoạt động của IP.

\textbf{12. session\_duration\_mean (Thời lượng phiên trung bình):} Giá
trị trung bình của thời lượng các phiên SSH (session duration), tính
bằng giây. Người dùng hợp pháp thường có phiên làm việc kéo dài (vài
phút đến vài giờ), trong khi tấn công brute-force tạo ra các phiên rất
ngắn (dưới vài giây) vì mỗi lần thử mật khẩu sai đều bị ngắt kết nối
nhanh chóng.

\textbf{Nhóm 5: Đặc trưng chỉ thị tấn công (Attack Indicator Features)}

\textbf{13. pam\_failure\_escalation (Leo thang lỗi PAM):} Biến nhị phân
(0 hoặc 1) cho biết trong cửa sổ thời gian có xuất hiện chuỗi lỗi PAM
liên tục gia tăng hay không. Giá trị 1 cho thấy có hiện tượng tăng dần
số lần thất bại PAM, đặc trưng của tấn công brute-force có hệ thống.

\textbf{14. max\_retries\_exceeded (Vượt quá số lần thử tối đa):} Biến
nhị phân cho biết trong cửa sổ thời gian có sự kiện ``maximum
authentication attempts exceeded'' hay không. Đây là dấu hiệu trực tiếp
của tấn công brute-force, khi kẻ tấn công cố gắng nhiều lần thử mật khẩu
trong cùng một phiên kết nối cho đến khi bị máy chủ SSH ngắt kết nối.

\textbf{Bảng 3.4: Tổng hợp 14 đặc trưng và ý nghĩa}

{\def\LTcaptype{none} % do not increment counter
\begin{longtable}[]{@{}
  >{\raggedright\arraybackslash}p{(\linewidth - 6\tabcolsep) * \real{0.1020}}
  >{\raggedright\arraybackslash}p{(\linewidth - 6\tabcolsep) * \real{0.2245}}
  >{\raggedright\arraybackslash}p{(\linewidth - 6\tabcolsep) * \real{0.2857}}
  >{\raggedright\arraybackslash}p{(\linewidth - 6\tabcolsep) * \real{0.3878}}@{}}
\toprule\noalign{}
\begin{minipage}[b]{\linewidth}\raggedright
STT
\end{minipage} & \begin{minipage}[b]{\linewidth}\raggedright
Đặc trưng
\end{minipage} & \begin{minipage}[b]{\linewidth}\raggedright
Kiểu dữ liệu
\end{minipage} & \begin{minipage}[b]{\linewidth}\raggedright
Ý nghĩa phát hiện
\end{minipage} \\
\midrule\noalign{}
\endhead
\bottomrule\noalign{}
\endlastfoot
1 & fail\_count & Số nguyên & Số lần thử sai mật khẩu \\
2 & success\_count & Số nguyên & Số lần đăng nhập thành công \\
3 & fail\_rate & Thực {[}0,1{]} & Tỷ lệ thất bại/tổng thử \\
4 & unique\_usernames & Số nguyên & Đa dạng tên người dùng \\
5 & invalid\_user\_count & Số nguyên & Tên người dùng không tồn tại \\
6 & invalid\_user\_ratio & Thực {[}0,1{]} & Tỷ lệ tên không hợp lệ \\
7 & connection\_count & Số nguyên & Tổng số kết nối SSH \\
8 & mean\_inter\_attempt\_time & Thực (giây) & Tốc độ trung bình giữa
các lần thử \\
9 & std\_inter\_attempt\_time & Thực (giây) & Biến thiên tốc độ thử \\
10 & min\_inter\_attempt\_time & Thực (giây) & Tốc độ nhanh nhất giữa
các lần thử \\
11 & unique\_ports & Số nguyên & Đa dạng cổng nguồn \\
12 & pam\_failure\_escalation & Nhị phân \{0,1\} & Chuỗi lỗi PAM gia
tăng \\
13 & max\_retries\_exceeded & Nhị phân \{0,1\} & Vượt giới hạn thử
lại \\
14 & session\_duration\_mean & Thực (giây) & Thời lượng phiên trung
bình \\
\end{longtable}
}

\subsection{3.4.3 Chuẩn hóa đặc
trưng}\label{chuux1ea9n-huxf3a-ux111ux1eb7c-trux1b0ng}

Sau khi trích xuất, các đặc trưng được chuẩn hóa bằng phương pháp
\textbf{RobustScaler} từ thư viện scikit-learn. RobustScaler được lựa
chọn thay vì StandardScaler hoặc MinMaxScaler vì các lý do:

\begin{enumerate}
\def\labelenumi{\arabic{enumi}.}
\item
  \textbf{Kháng nhiễu (Robust to outliers):} RobustScaler sử dụng trung
  vị (median) và khoảng tứ phân vị (interquartile range - IQR) thay vì
  trung bình (mean) và độ lệch chuẩn (standard deviation), giúp giảm ảnh
  hưởng của các giá trị ngoại lai (outliers) trong dữ liệu.
\item
  \textbf{Phù hợp với dữ liệu tấn công:} Dữ liệu tấn công brute-force
  thường chứa nhiều giá trị cực trị (extreme values), ví dụ: fail\_count
  có thể lên tới hàng nghìn trong khi giá trị bình thường chỉ 0-2.
  RobustScaler không bị ảnh hưởng bởi các giá trị cực trị này.
\end{enumerate}

Công thức chuẩn hóa của RobustScaler:

\[x_{scaled} = \frac{x - Q_2(x)}{Q_3(x) - Q_1(x)}\]

trong đó \(Q_1\), \(Q_2\), \(Q_3\) lần lượt là phân vị thứ 25\%, 50\%
(trung vị), và 75\% của đặc trưng \(x\).

Scaler được fit trên tập huấn luyện (chỉ chứa dữ liệu normal) và áp dụng
transform cho cả tập huấn luyện và tập kiểm thử, đảm bảo không có rò rỉ
dữ liệu (data leakage).

\section{3.5 Lựa chọn mô hình}\label{lux1ef1a-chux1ecdn-muxf4-huxecnh}

\subsection{3.5.1 Phương pháp tiếp cận: Phát hiện bất thường bán giám
sát}\label{phux1b0ux1a1ng-phuxe1p-tiux1ebfp-cux1eadn-phuxe1t-hiux1ec7n-bux1ea5t-thux1b0ux1eddng-buxe1n-giuxe1m-suxe1t}

Nghiên cứu áp dụng phương pháp phát hiện bất thường bán giám sát
(semi-supervised anomaly detection), trong đó mô hình chỉ được huấn
luyện trên dữ liệu bình thường (normal data) và sau đó xác định các mẫu
lệch khỏi phân phối bình thường là bất thường (anomaly). Phương pháp này
được lựa chọn vì:

\begin{enumerate}
\def\labelenumi{\arabic{enumi}.}
\item
  \textbf{Tính thực tế:} Trong thực tế, dữ liệu bình thường dễ thu thập
  hơn nhiều so với dữ liệu tấn công có nhãn. Các tổ chức có thể dễ dàng
  thu thập log hoạt động bình thường nhưng khó có đủ mẫu cho mọi loại
  tấn công.
\item
  \textbf{Khả năng phát hiện tấn công mới:} Mô hình bán giám sát có khả
  năng phát hiện các loại tấn công chưa từng thấy (zero-day attacks),
  miễn là hành vi tấn công lệch khỏi mẫu bình thường đã học.
\item
  \textbf{Tránh mất cân bằng dữ liệu:} Phương pháp giám sát (supervised)
  thường gặp vấn đề mất cân bằng lớp (class imbalance) khi tỷ lệ tấn
  công thấp. Phương pháp bán giám sát không bị ảnh hưởng bởi vấn đề này.
\end{enumerate}

\subsection{3.5.2 Ba mô hình được lựa
chọn}\label{ba-muxf4-huxecnh-ux111ux1b0ux1ee3c-lux1ef1a-chux1ecdn}

Nghiên cứu lựa chọn và so sánh ba mô hình phát hiện bất thường phổ biến
nhất trong lĩnh vực an toàn thông tin:

\textbf{Isolation Forest (IF):} Isolation Forest (Liu và cộng sự, 2008)
là thuật toán dựa trên nguyên lý cô lập (isolation). Ý tưởng cốt lõi:
các điểm bất thường dễ bị cô lập hơn các điểm bình thường trong cây
quyết định ngẫu nhiên. Thuật toán xây dựng một tập hợp (ensemble) các
cây cô lập (isolation trees), trong đó mỗi cây phân chia không gian đặc
trưng bằng các siêu phẳng ngẫu nhiên. Điểm bất thường (anomaly score) tỷ
lệ nghịch với chiều dài đường đi trung bình (average path length) từ gốc
đến lá.

Ưu điểm: hiệu quả tính toán cao (độ phức tạp O(n log n)), phù hợp với dữ
liệu nhiều chiều, không yêu cầu giả định về phân phối dữ liệu.

\textbf{Local Outlier Factor (LOF):} LOF (Breunig và cộng sự, 2000) là
thuật toán dựa trên mật độ cục bộ (local density). Thuật toán so sánh
mật độ cục bộ của một điểm dữ liệu với mật độ cục bộ của các điểm láng
giềng gần nhất (nearest neighbors). Một điểm có mật độ thấp hơn đáng kể
so với các láng giềng được coi là bất thường.

Ưu điểm: phát hiện tốt các bất thường cục bộ (local anomalies), không
yêu cầu giả định phân phối toàn cục.

\textbf{One-Class SVM (OCSVM):} OCSVM (Schölkopf và cộng sự, 2001) là
biến thể của Support Vector Machine cho bài toán phân loại một lớp.
Thuật toán tìm siêu phẳng (hyperplane) trong không gian đặc trưng ánh xạ
(kernel space) sao cho phân tách được dữ liệu huấn luyện khỏi gốc tọa độ
với biên lớn nhất (maximum margin). Các điểm nằm ngoài siêu phẳng được
phân loại là bất thường.

Ưu điểm: mạnh mẽ về mặt lý thuyết, hiệu quả trong không gian đặc trưng
cao chiều thông qua kernel trick.

\textbf{Bảng 3.5: So sánh đặc điểm của ba mô hình}

{\def\LTcaptype{none} % do not increment counter
\begin{longtable}[]{@{}llll@{}}
\toprule\noalign{}
Đặc điểm & Isolation Forest & LOF & OCSVM \\
\midrule\noalign{}
\endhead
\bottomrule\noalign{}
\endlastfoot
Nguyên lý & Cô lập & Mật độ cục bộ & Siêu phẳng biên \\
Độ phức tạp & O(n log n) & O(n²) & O(n² \textasciitilde{} n³) \\
Phù hợp dữ liệu lớn & Rất tốt & Trung bình & Trung bình \\
Phát hiện bất thường cục bộ & Trung bình & Rất tốt & Tốt \\
Giả định phân phối & Không & Không & Không (với kernel) \\
Khả năng diễn giải & Trung bình & Thấp & Thấp \\
\end{longtable}
}

\section{3.6 Phương pháp huấn
luyện}\label{phux1b0ux1a1ng-phuxe1p-huux1ea5n-luyux1ec7n}

\subsection{3.6.1 Quy trình huấn luyện bán giám
sát}\label{quy-truxecnh-huux1ea5n-luyux1ec7n-buxe1n-giuxe1m-suxe1t}

Quy trình huấn luyện tuân theo paradigm semi-supervised novelty
detection:

\begin{enumerate}
\def\labelenumi{\arabic{enumi}.}
\tightlist
\item
  \textbf{Bước 1 - Chuẩn bị dữ liệu huấn luyện:} Chỉ sử dụng 7.212 mẫu
  normal từ 70\% dữ liệu simulation.
\item
  \textbf{Bước 2 - Chuẩn hóa:} Fit RobustScaler trên tập huấn luyện,
  transform tập huấn luyện và kiểm thử.
\item
  \textbf{Bước 3 - Huấn luyện mô hình:} Fit mỗi mô hình (IF, LOF, OCSVM)
  trên tập huấn luyện đã chuẩn hóa.
\item
  \textbf{Bước 4 - Dự đoán và đánh giá:} Tính anomaly score cho mỗi mẫu
  trong tập kiểm thử (15.184 mẫu), áp dụng ngưỡng để phân loại, và tính
  các chỉ số đánh giá.
\end{enumerate}

\textbf{Hình 3.4: Quy trình huấn luyện và đánh giá mô hình}

\subsection{3.6.2 Tối ưu siêu tham số (Hyperparameter
Tuning)}\label{tux1ed1i-ux1b0u-siuxeau-tham-sux1ed1-hyperparameter-tuning}

Siêu tham số của mỗi mô hình được tối ưu thông qua phương pháp tìm kiếm
lưới (grid search) kết hợp với xác thực chéo (cross-validation) trên tập
huấn luyện. Do tập huấn luyện chỉ chứa dữ liệu normal, tiêu chí tối ưu
dựa trên khả năng tái tạo (reconstruction) phân phối dữ liệu normal và
điểm ROC-AUC trên một tập validation nhỏ chứa một số mẫu attack.

\textbf{Isolation Forest - Siêu tham số tối ưu:}

{\def\LTcaptype{none} % do not increment counter
\begin{longtable}[]{@{}lll@{}}
\toprule\noalign{}
Siêu tham số & Giá trị tối ưu & Mô tả \\
\midrule\noalign{}
\endhead
\bottomrule\noalign{}
\endlastfoot
n\_estimators & 300 & Số cây cô lập trong tập hợp \\
max\_samples & 512 & Số mẫu con cho mỗi cây \\
max\_features & 0.5 & Tỷ lệ đặc trưng cho mỗi cây (7/14 đặc trưng) \\
\end{longtable}
}

Việc sử dụng \texttt{max\_features=0.5} giúp tăng tính đa dạng
(diversity) giữa các cây trong ensemble, giảm nguy cơ overfitting.
\texttt{max\_samples=512} giới hạn kích thước mẫu con, giúp tăng tốc
huấn luyện và cải thiện khả năng phát hiện bất thường theo lý thuyết
Isolation Forest. \texttt{n\_estimators=300} đảm bảo đủ số cây để ổn
định anomaly score.

\textbf{Local Outlier Factor - Siêu tham số tối ưu:}

{\def\LTcaptype{none} % do not increment counter
\begin{longtable}[]{@{}lll@{}}
\toprule\noalign{}
Siêu tham số & Giá trị tối ưu & Mô tả \\
\midrule\noalign{}
\endhead
\bottomrule\noalign{}
\endlastfoot
n\_neighbors & 30 & Số láng giềng gần nhất \\
novelty & True & Chế độ novelty detection \\
metric & minkowski & Hàm khoảng cách \\
\end{longtable}
}

Giá trị \texttt{n\_neighbors=30} cân bằng giữa khả năng phát hiện bất
thường cục bộ (local anomaly) và tính ổn định của ước lượng mật độ. Giá
trị quá nhỏ dẫn đến nhạy cảm với nhiễu, giá trị quá lớn có thể bỏ lỡ các
bất thường cục bộ.

\textbf{One-Class SVM - Siêu tham số tối ưu:}

{\def\LTcaptype{none} % do not increment counter
\begin{longtable}[]{@{}lll@{}}
\toprule\noalign{}
Siêu tham số & Giá trị tối ưu & Mô tả \\
\midrule\noalign{}
\endhead
\bottomrule\noalign{}
\endlastfoot
kernel & rbf & Hàm nhân Gaussian \\
gamma & auto & Hệ số kernel (1/n\_features) \\
nu & 0.01 & Biên trên tỷ lệ ngoại lai và biên dưới support vectors \\
\end{longtable}
}

Giá trị \texttt{nu=0.01} thiết lập biên trên cho tỷ lệ ngoại lai
(outlier fraction) là 1\%, phù hợp với giả định rằng tập huấn luyện
thuần normal và chỉ cho phép tối đa 1\% mẫu bị phân loại nhầm.
\texttt{gamma=auto} tự động tính gamma = 1/n\_features = 1/14 ≈ 0.0714,
phù hợp với số chiều đặc trưng.

\section{3.7 Thiết kế thuật toán ngưỡng
động}\label{thiux1ebft-kux1ebf-thuux1eadt-touxe1n-ngux1b0ux1ee1ng-ux111ux1ed9ng}

\subsection{3.7.1 Động cơ thiết
kế}\label{ux111ux1ed9ng-cux1a1-thiux1ebft-kux1ebf}

Trong các hệ thống phát hiện bất thường truyền thống, ngưỡng phân loại
(classification threshold) thường được cố định (static threshold). Tuy
nhiên, trong thực tế, phân phối điểm bất thường có thể thay đổi theo
thời gian do:

\begin{itemize}
\tightlist
\item
  Thay đổi trong mẫu lưu lượng mạng (traffic pattern shifts).
\item
  Sự xuất hiện của các loại tấn công mới với đặc trưng khác nhau.
\item
  Biến đổi theo thời gian trong hành vi người dùng (concept drift).
\end{itemize}

Do đó, nghiên cứu đề xuất thuật toán ngưỡng động \textbf{EWMA-Adaptive
Percentile} kết hợp hai kỹ thuật: trung bình trượt có trọng số hàm mũ
(Exponentially Weighted Moving Average - EWMA) và phân vị thích ứng
(adaptive percentile).

\subsection{3.7.2 Công thức toán
học}\label{cuxf4ng-thux1ee9c-touxe1n-hux1ecdc}

\textbf{Định nghĩa các tham số:}

\begin{itemize}
\tightlist
\item
  \(\alpha\) = 0.3: Hệ số làm mượt EWMA (smoothing factor)
\item
  \(base\_percentile\) = 95: Phân vị cơ sở (base percentile)
\item
  \(sensitivity\_factor\) = 1.5: Hệ số nhạy cảm (sensitivity factor)
\item
  \(lookback\) = 100: Số điểm dữ liệu gần nhất để tính toán (lookback
  window)
\end{itemize}

\textbf{Bước 1: Tính EWMA của anomaly score}

Cho chuỗi điểm bất thường \(s_1, s_2, ..., s_t\) theo thời gian, EWMA
tại thời điểm \(t\) được tính:

\[\mu_t^{EWMA} = \alpha \cdot s_t + (1 - \alpha) \cdot \mu_{t-1}^{EWMA}\]

với \(\mu_0^{EWMA} = s_1\) (khởi tạo bằng giá trị đầu tiên).

Hệ số \(\alpha = 0.3\) đảm bảo EWMA phản ứng đủ nhanh với các thay đổi
gần đây nhưng không quá nhạy cảm với nhiễu nhất thời.

\textbf{Bước 2: Tính phân vị thích ứng}

Sử dụng lookback window gồm \(L = 100\) điểm bất thường gần nhất
\(\{s_{t-L+1}, ..., s_t\}\), tính phân vị thứ \(p\) (percentile):

\[P_t = Percentile(\{s_{t-L+1}, ..., s_t\}, p)\]

trong đó \(p = base\_percentile = 95\).

\textbf{Bước 3: Tính ngưỡng động (dynamic threshold)}

Ngưỡng động tại thời điểm \(t\) được tính bằng:

\[\theta_t = \mu_t^{EWMA} + sensitivity\_factor \times (P_t - \mu_t^{EWMA})\]

\[\theta_t = \mu_t^{EWMA} + 1.5 \times (P_t - \mu_t^{EWMA})\]

Công thức này có ý nghĩa: ngưỡng được đặt tại vị trí trung bình EWMA
cộng thêm 1.5 lần khoảng cách từ EWMA đến phân vị thứ 95. Điều này đảm
bảo:

\begin{itemize}
\tightlist
\item
  Khi phân phối ổn định: ngưỡng nằm trên phần lớn (\textgreater95\%) các
  điểm bình thường.
\item
  Khi xuất hiện đợt tấn công: EWMA tăng lên, kéo ngưỡng lên theo, tránh
  cảnh báo sai liên tục.
\item
  Khi tấn công kết thúc: EWMA giảm dần, ngưỡng trở về mức bình thường.
\end{itemize}

\textbf{Bước 4: Quyết định phân loại}

Một mẫu tại thời điểm \(t\) được phân loại:

\[
y_t = \begin{cases}
\text{attack} & \text{nếu } s_t > \theta_t \\
\text{normal} & \text{nếu } s_t \leq \theta_t
\end{cases}
\]

\textbf{Hình 3.5: Minh họa thuật toán ngưỡng động EWMA-Adaptive
Percentile}

\subsection{3.7.3 Phân tích tham số}\label{phuxe2n-tuxedch-tham-sux1ed1}

\textbf{Bảng 3.6: Tham số thuật toán ngưỡng động và ý nghĩa}

{\def\LTcaptype{none} % do not increment counter
\begin{longtable}[]{@{}
  >{\raggedright\arraybackslash}p{(\linewidth - 6\tabcolsep) * \real{0.2368}}
  >{\raggedright\arraybackslash}p{(\linewidth - 6\tabcolsep) * \real{0.2368}}
  >{\raggedright\arraybackslash}p{(\linewidth - 6\tabcolsep) * \real{0.2368}}
  >{\raggedright\arraybackslash}p{(\linewidth - 6\tabcolsep) * \real{0.2895}}@{}}
\toprule\noalign{}
\begin{minipage}[b]{\linewidth}\raggedright
Tham số
\end{minipage} & \begin{minipage}[b]{\linewidth}\raggedright
Giá trị
\end{minipage} & \begin{minipage}[b]{\linewidth}\raggedright
Ý nghĩa
\end{minipage} & \begin{minipage}[b]{\linewidth}\raggedright
Ảnh hưởng
\end{minipage} \\
\midrule\noalign{}
\endhead
\bottomrule\noalign{}
\endlastfoot
alpha (α) & 0.3 & Tốc độ phản ứng EWMA & α lớn → phản ứng nhanh, nhiều
nhiễu; α nhỏ → phản ứng chậm, ổn định \\
base\_percentile & 95 & Ngưỡng phân vị cơ sở & Cao → ít FP, có thể bỏ lỡ
tấn công; Thấp → nhiều FP, phát hiện tốt \\
sensitivity\_factor & 1.5 & Độ nhạy cảm & Cao → ngưỡng cao, ít cảnh báo;
Thấp → ngưỡng thấp, nhiều cảnh báo \\
lookback & 100 & Số mẫu gần nhất & Lớn → ổn định, chậm thích ứng; Nhỏ →
nhạy, nhanh thích ứng \\
\end{longtable}
}

\section{3.8 Pipeline phát hiện thời gian
thực}\label{pipeline-phuxe1t-hiux1ec7n-thux1eddi-gian-thux1ef1c}

\subsection{3.8.1 Kiến trúc pipeline}\label{kiux1ebfn-truxfac-pipeline}

Pipeline phát hiện thời gian thực (real-time detection pipeline) được
thiết kế để xử lý liên tục dòng sự kiện log SSH với độ trễ tối thiểu.
Kiến trúc pipeline bao gồm các giai đoạn (stages):

\textbf{Giai đoạn 1 - Ingestion (Thu nhận dữ liệu):} Logstash theo dõi
(tail) file auth.log trên các máy chủ SSH, phân tích cú pháp mỗi dòng
log mới, và gửi sự kiện đã cấu trúc hóa tới Elasticsearch. Đồng thời,
một webhook/event stream gửi sự kiện tới API Server.

\textbf{Giai đoạn 2 - Aggregation (Tổng hợp):} API Server duy trì bộ đệm
(buffer) theo cửa sổ trượt cho mỗi địa chỉ IP đang hoạt động. Khi bước
trượt 1 phút trôi qua, hệ thống tổng hợp các sự kiện trong cửa sổ 5 phút
gần nhất để tính toán 14 đặc trưng.

\textbf{Giai đoạn 3 - Scoring (Tính điểm):} Vector đặc trưng được chuẩn
hóa bằng RobustScaler (đã fit trên tập huấn luyện) và đưa vào mô hình AI
để tính anomaly score. Cả ba mô hình (IF, LOF, OCSVM) có thể hoạt động
song song hoặc theo cấu hình ensemble.

\textbf{Giai đoạn 4 - Decision (Quyết định):} Anomaly score được so sánh
với ngưỡng động hiện tại (\(\theta_t\)). Đồng thời, thuật toán
EWMA-Adaptive Percentile cập nhật ngưỡng dựa trên điểm mới nhận được.

\textbf{Giai đoạn 5 - Action (Hành động):} Nếu phát hiện tấn công, hệ
thống kích hoạt các hành động phòng chống (xem mục 3.9).

\textbf{Hình 3.6: Pipeline phát hiện thời gian thực 5 giai đoạn}

\subsection{3.8.2 Xử lý song song và bất đồng
bộ}\label{xux1eed-luxfd-song-song-vuxe0-bux1ea5t-ux111ux1ed3ng-bux1ed9}

Để đảm bảo hiệu năng thời gian thực, pipeline sử dụng:

\begin{itemize}
\tightlist
\item
  \textbf{Xử lý bất đồng bộ (Asynchronous processing):} FastAPI hỗ trợ
  xử lý bất đồng bộ (async/await) cho các tác vụ I/O-bound như truy vấn
  Elasticsearch, ghi database.
\item
  \textbf{Xử lý song song (Parallel processing):} Nhiều IP có thể được
  xử lý đồng thời thông qua thread pool hoặc process pool cho các tác vụ
  CPU-bound như tính toán mô hình.
\item
  \textbf{Caching:} Redis được sử dụng để cache kết quả trung gian, bao
  gồm: scaler parameters, model objects, và trạng thái cửa sổ trượt cho
  mỗi IP.
\end{itemize}

\section{3.9 Module cảnh báo và phòng
chống}\label{module-cux1ea3nh-buxe1o-vuxe0-phuxf2ng-chux1ed1ng}

\subsection{3.9.1 Tích hợp Fail2Ban}\label{tuxedch-hux1ee3p-fail2ban}

Fail2Ban là công cụ phòng chống xâm nhập (intrusion prevention) được
tích hợp vào hệ thống để tự động chặn các địa chỉ IP tấn công. Khi
module phát hiện xác định một IP đang thực hiện tấn công brute-force, hệ
thống gọi API Fail2Ban để:

\begin{enumerate}
\def\labelenumi{\arabic{enumi}.}
\tightlist
\item
  \textbf{Ban IP:} Thêm địa chỉ IP vào danh sách chặn (ban list) của
  Fail2Ban, ngăn mọi kết nối mới từ IP đó.
\item
  \textbf{Thời gian chặn:} Cấu hình thời gian chặn (ban time) linh hoạt
  dựa trên mức độ nghiêm trọng (severity) của cuộc tấn công.
\item
  \textbf{Chặn leo thang (Progressive ban):} Nếu một IP bị phát hiện tấn
  công nhiều lần, thời gian chặn tăng dần theo cấp số nhân.
\end{enumerate}

\subsection{3.9.2 Hệ thống cảnh báo đa
kênh}\label{hux1ec7-thux1ed1ng-cux1ea3nh-buxe1o-ux111a-kuxeanh}

Hệ thống cảnh báo hỗ trợ nhiều kênh thông báo:

\begin{itemize}
\tightlist
\item
  \textbf{Dashboard React:} Hiển thị cảnh báo thời gian thực trên giao
  diện web, bao gồm thông tin IP, thời gian, anomaly score, loại tấn
  công dự đoán, và hành động đã thực hiện.
\item
  \textbf{Webhook:} Gửi cảnh báo tới các dịch vụ bên ngoài (Slack,
  Telegram, v.v.) thông qua webhook HTTP.
\item
  \textbf{Logging:} Ghi nhận tất cả sự kiện phát hiện và hành động phòng
  chống vào Elasticsearch để phục vụ phân tích hậu sự cố (post-incident
  analysis).
\end{itemize}

\subsection{3.9.3 Cơ chế phản hồi (Feedback
Mechanism)}\label{cux1a1-chux1ebf-phux1ea3n-hux1ed3i-feedback-mechanism}

Hệ thống cho phép quản trị viên đánh dấu các cảnh báo là true positive
hoặc false positive thông qua giao diện React. Thông tin phản hồi này
được lưu trữ và có thể được sử dụng để tinh chỉnh mô hình hoặc điều
chỉnh tham số ngưỡng động trong tương lai.

\section{3.10 Các chỉ số đánh
giá}\label{cuxe1c-chux1ec9-sux1ed1-ux111uxe1nh-giuxe1}

\subsection{3.10.1 Ma trận nhầm lẫn (Confusion
Matrix)}\label{ma-trux1eadn-nhux1ea7m-lux1eabn-confusion-matrix}

Các chỉ số đánh giá dựa trên ma trận nhầm lẫn cho bài toán phân loại nhị
phân (normal vs.~attack):

\begin{itemize}
\tightlist
\item
  \textbf{True Positive (TP):} Mẫu tấn công được phân loại đúng là tấn
  công.
\item
  \textbf{True Negative (TN):} Mẫu bình thường được phân loại đúng là
  bình thường.
\item
  \textbf{False Positive (FP):} Mẫu bình thường bị phân loại nhầm là tấn
  công (cảnh báo sai).
\item
  \textbf{False Negative (FN):} Mẫu tấn công bị phân loại nhầm là bình
  thường (bỏ sót tấn công).
\end{itemize}

\subsection{3.10.2 Các chỉ số
chính}\label{cuxe1c-chux1ec9-sux1ed1-chuxednh}

\textbf{Accuracy (Độ chính xác tổng thể):}

\[Accuracy = \frac{TP + TN}{TP + TN + FP + FN}\]

Đo tỷ lệ phân loại đúng trên toàn bộ tập dữ liệu. Tuy nhiên, trong bài
toán mất cân bằng, accuracy có thể không phản ánh đúng hiệu năng mô
hình.

\textbf{Precision (Độ chính xác dương tính):}

\[Precision = \frac{TP}{TP + FP}\]

Đo tỷ lệ dự đoán tấn công thực sự là tấn công trong tổng số dự đoán tấn
công. Precision cao nghĩa là ít cảnh báo sai.

\textbf{Recall / Sensitivity (Độ nhạy / Độ phủ):}

\[Recall = \frac{TP}{TP + FN}\]

Đo tỷ lệ tấn công thực sự được phát hiện trong tổng số tấn công thực tế.
Recall cao nghĩa là ít bỏ sót tấn công. Trong lĩnh vực an toàn thông
tin, recall đặc biệt quan trọng vì việc bỏ sót một cuộc tấn công có thể
gây hậu quả nghiêm trọng.

\textbf{F1-Score:}

\[F1 = 2 \times \frac{Precision \times Recall}{Precision + Recall}\]

Trung bình điều hòa của Precision và Recall, cung cấp một chỉ số cân
bằng giữa hai yếu tố.

\textbf{ROC-AUC (Area Under the Receiver Operating Characteristic
Curve):}

\[ROC\text{-}AUC = \int_0^1 TPR(FPR^{-1}(x)) \, dx\]

Diện tích dưới đường cong ROC, đo khả năng phân biệt giữa lớp normal và
attack trên toàn bộ các ngưỡng phân loại. ROC-AUC = 1.0 là lý tưởng
(phân biệt hoàn hảo), ROC-AUC = 0.5 tương đương phân loại ngẫu nhiên.

\textbf{Bảng 3.7: Tổng hợp các chỉ số đánh giá và ý nghĩa trong ngữ cảnh
bảo mật}

{\def\LTcaptype{none} % do not increment counter
\begin{longtable}[]{@{}lll@{}}
\toprule\noalign{}
Chỉ số & Ý nghĩa bảo mật & Ưu tiên \\
\midrule\noalign{}
\endhead
\bottomrule\noalign{}
\endlastfoot
Accuracy & Hiệu năng tổng thể & Trung bình \\
Precision & Giảm cảnh báo sai (alert fatigue) & Cao \\
Recall & Không bỏ sót tấn công & Rất cao \\
F1-Score & Cân bằng Precision-Recall & Cao \\
ROC-AUC & Khả năng phân biệt tổng quát & Cao \\
\end{longtable}
}

\subsection{3.10.3 Các chỉ số hiệu năng hệ
thống}\label{cuxe1c-chux1ec9-sux1ed1-hiux1ec7u-nux103ng-hux1ec7-thux1ed1ng}

Ngoài các chỉ số phân loại, nghiên cứu còn đánh giá hiệu năng hệ thống
thời gian thực:

\begin{itemize}
\tightlist
\item
  \textbf{Latency (Độ trễ):} Thời gian từ khi nhận sự kiện log đến khi
  đưa ra quyết định phân loại.
\item
  \textbf{Throughput (Thông lượng):} Số sự kiện xử lý được trong một đơn
  vị thời gian.
\item
  \textbf{Memory usage (Sử dụng bộ nhớ):} Lượng bộ nhớ RAM tiêu thụ bởi
  mỗi thành phần.
\item
  \textbf{CPU usage (Sử dụng CPU):} Mức sử dụng CPU của các dịch vụ
  trong điều kiện tải bình thường và tải cao.
\end{itemize}

Các chỉ số này đảm bảo hệ thống có khả năng hoạt động trong môi trường
sản xuất thực tế với yêu cầu thời gian thực.

\begin{center}\rule{0.5\linewidth}{0.5pt}\end{center}

\textbf{Tóm tắt Chương 3:} Chương này đã trình bày chi tiết phương pháp
nghiên cứu bao gồm: kiến trúc tổng thể hệ thống với 9 dịch vụ Docker,
quy trình thu thập và tiền xử lý dữ liệu từ hai nguồn (honeypot và
simulation), chiến lược gán nhãn bán giám sát, 14 đặc trưng được thiết
kế và trích xuất theo cửa sổ trượt 5 phút, ba mô hình phát hiện bất
thường (IF, LOF, OCSVM) cùng phương pháp tối ưu siêu tham số, thuật toán
ngưỡng động EWMA-Adaptive Percentile với các công thức toán học chi
tiết, pipeline phát hiện thời gian thực, và module cảnh báo phòng chống
tích hợp Fail2Ban. Các chỉ số đánh giá bao gồm Accuracy, Precision,
Recall, F1-Score, và ROC-AUC được định nghĩa rõ ràng cùng ý nghĩa trong
ngữ cảnh an toàn thông tin. \# CHƯƠNG 4: KẾT QUẢ THỰC NGHIỆM

\section{4.1 Thống kê tập dữ liệu và phân tích khám
phá}\label{thux1ed1ng-kuxea-tux1eadp-dux1eef-liux1ec7u-vuxe0-phuxe2n-tuxedch-khuxe1m-phuxe1}

\subsection{4.1.1 Tổng quan tập dữ
liệu}\label{tux1ed5ng-quan-tux1eadp-dux1eef-liux1ec7u}

Quá trình thu thập và tiền xử lý dữ liệu tạo ra tổng cộng 22.396 mẫu
(samples) từ hai nguồn dữ liệu. Bảng 4.1 trình bày thống kê tổng quan về
tập dữ liệu cuối cùng sau khi trích xuất đặc trưng theo cửa sổ trượt 5
phút với bước trượt 1 phút.

\textbf{Bảng 4.1: Thống kê tổng quan tập dữ liệu}

{\def\LTcaptype{none} % do not increment counter
\begin{longtable}[]{@{}ll@{}}
\toprule\noalign{}
Thuộc tính & Giá trị \\
\midrule\noalign{}
\endhead
\bottomrule\noalign{}
\endlastfoot
Tổng số mẫu & 22.396 \\
Tập huấn luyện & 7.212 mẫu (100\% normal) \\
Tập kiểm thử & 15.184 mẫu (3.796 normal + 11.388 attack) \\
Số đặc trưng & 14 \\
Kích thước cửa sổ & 5 phút \\
Bước trượt & 1 phút \\
\end{longtable}
}

\subsection{4.1.2 Phân tích dữ liệu
honeypot}\label{phuxe2n-tuxedch-dux1eef-liux1ec7u-honeypot}

Dữ liệu từ honeypot\_auth.log (119.729 dòng, 5 ngày) cho thấy bức tranh
toàn diện về hoạt động tấn công brute-force SSH trên Internet công cộng.
Phân tích khám phá dữ liệu (Exploratory Data Analysis - EDA) cho thấy
các đặc điểm quan trọng:

\textbf{Phân bố theo thời gian:} Các cuộc tấn công diễn ra liên tục
24/7, không có mẫu thời gian rõ ràng, cho thấy đặc trưng của các công cụ
tấn công tự động (automated attack tools). Trong 5 ngày thu thập, trung
bình mỗi ngày ghi nhận khoảng 5.860 lần thử mật khẩu thất bại (29.301 /
5 ngày).

\textbf{Phân bố theo IP:} Tổng cộng 679 địa chỉ IP duy nhất được ghi
nhận. Phân tích cho thấy phân bố lệch phải (right-skewed distribution)
rõ rệt: một số ít IP thực hiện hàng nghìn lần thử, trong khi phần lớn IP
chỉ thực hiện vài chục đến vài trăm lần. Top 10\% IP đóng góp hơn 60\%
tổng số lần thử thất bại.

\textbf{Hình 4.1: Phân bố số lần thử xác thực thất bại theo IP
(honeypot)}

\textbf{Phân tích tên người dùng bị tấn công:} Tên người dùng phổ biến
nhất trong các cuộc tấn công: ``root'' chiếm tỷ lệ cao nhất (ước tính
\textgreater40\%), tiếp theo là ``admin'', ``test'', ``user'',
``oracle'', ``postgres'', ``ubuntu'', và các tên người dùng dịch vụ
khác. Có 29.301 lần thử mật khẩu thất bại trong tổng số 119.729 dòng
log.

\textbf{Hoạt động quản trị hợp pháp:} 6 địa chỉ IP quản trị tạo ra 532
phiên đăng nhập root thành công, phân bố chủ yếu trong giờ hành chính,
với phiên làm việc kéo dài từ vài phút đến vài giờ --- đặc trưng của
hoạt động quản trị hệ thống bình thường.

\textbf{Bảng 4.2: Thống kê chi tiết dữ liệu honeypot}

{\def\LTcaptype{none} % do not increment counter
\begin{longtable}[]{@{}ll@{}}
\toprule\noalign{}
Thống kê & Giá trị \\
\midrule\noalign{}
\endhead
\bottomrule\noalign{}
\endlastfoot
Tổng dòng log & 119.729 \\
Thời gian thu thập & 5 ngày \\
IP duy nhất & 679 \\
Failed password & 29.301 \\
Accepted root (admin IPs) & 532 \\
Số IP quản trị & 6 \\
Trung bình failed/ngày & \textasciitilde5.860 \\
Trung bình failed/IP & \textasciitilde43,2 \\
\end{longtable}
}

\subsection{4.1.3 Phân tích dữ liệu mô
phỏng}\label{phuxe2n-tuxedch-dux1eef-liux1ec7u-muxf4-phux1ecfng}

Dữ liệu simulation\_auth.log (54.521 dòng, 64 người dùng) phản ánh mẫu
hoạt động SSH bình thường trong một tổ chức. Các đặc điểm nổi bật:

\textbf{Phân bố theo thời gian:} Hoạt động tập trung chủ yếu trong giờ
hành chính (8h-18h), giảm rõ rệt vào ban đêm và cuối tuần, phản ánh đúng
mẫu làm việc thực tế.

\textbf{Tỷ lệ thành công cao:} Với 4.205 lần đăng nhập thành công và chỉ
177 lần thất bại, tỷ lệ thành công đạt khoảng 95,96\%
(4.205/(4.205+177)). Tỷ lệ thất bại chỉ 4,04\% phản ánh các trường hợp
nhập sai mật khẩu do vô tình --- hành vi hoàn toàn bình thường.

\textbf{Phân bố người dùng:} 64 tài khoản hoạt động với mức độ sử dụng
SSH khác nhau. Một số tài khoản (ví dụ: nhà phát triển, quản trị hệ
thống) có tần suất đăng nhập cao hơn nhiều so với người dùng thông
thường.

\textbf{Hình 4.2: So sánh phân bố hoạt động theo giờ giữa honeypot
(attack) và simulation (normal)}

\textbf{Bảng 4.3: Thống kê chi tiết dữ liệu mô phỏng}

{\def\LTcaptype{none} % do not increment counter
\begin{longtable}[]{@{}ll@{}}
\toprule\noalign{}
Thống kê & Giá trị \\
\midrule\noalign{}
\endhead
\bottomrule\noalign{}
\endlastfoot
Tổng dòng log & 54.521 \\
Số người dùng & 64 \\
Accepted logins & 4.205 \\
Failed logins & 177 \\
Tỷ lệ thành công & 95,96\% \\
Tỷ lệ thất bại & 4,04\% \\
\end{longtable}
}

\subsection{4.1.4 So sánh đặc trưng giữa Normal và
Attack}\label{so-suxe1nh-ux111ux1eb7c-trux1b0ng-giux1eefa-normal-vuxe0-attack}

Phân tích khám phá cho thấy sự khác biệt rõ ràng giữa hành vi normal và
attack trên hầu hết 14 đặc trưng. Bảng 4.4 trình bày thống kê mô tả cho
các đặc trưng quan trọng nhất.

\textbf{Bảng 4.4: Thống kê mô tả các đặc trưng chính theo nhãn}

{\def\LTcaptype{none} % do not increment counter
\begin{longtable}[]{@{}llll@{}}
\toprule\noalign{}
Đặc trưng & Normal (mean ± std) & Attack (mean ± std) & Khác biệt \\
\midrule\noalign{}
\endhead
\bottomrule\noalign{}
\endlastfoot
fail\_count & Thấp (0-2) & Cao (hàng chục-trăm) & Rất lớn \\
fail\_rate & \textless0.3 & \textgreater0.9 & Rất lớn \\
unique\_usernames & 1-2 & 5-20+ & Lớn \\
invalid\_user\_count & \textasciitilde0 & Cao & Rất lớn \\
mean\_inter\_attempt\_time & Cao (\textgreater10s) & Rất thấp
(\textless1s) & Rất lớn \\
session\_duration\_mean & Cao (phút-giờ) & Rất thấp (\textless5s) & Rất
lớn \\
connection\_count & Thấp (1-5) & Cao (hàng chục-trăm) & Lớn \\
\end{longtable}
}

Kết quả phân tích cho thấy: các đặc trưng liên quan đến thời gian
(mean\_inter\_attempt\_time, min\_inter\_attempt\_time,
session\_duration\_mean) và tỷ lệ (fail\_rate, invalid\_user\_ratio) có
khả năng phân biệt tốt nhất giữa hai lớp. Điều này phù hợp với trực
giác: tấn công brute-force tự động tạo ra mẫu thời gian rất khác biệt so
với hoạt động thủ công của con người.

\textbf{Hình 4.3: Boxplot so sánh phân bố các đặc trưng chính giữa
Normal và Attack}

\section{4.2 Phân tích đặc
trưng}\label{phuxe2n-tuxedch-ux111ux1eb7c-trux1b0ng}

\subsection{4.2.1 Phân bố đặc
trưng}\label{phuxe2n-bux1ed1-ux111ux1eb7c-trux1b0ng}

Phân tích phân bố (distribution analysis) của 14 đặc trưng trên tập dữ
liệu cho thấy:

\textbf{Đặc trưng đếm (count features):} fail\_count, success\_count,
connection\_count, invalid\_user\_count có phân bố lệch phải mạnh
(heavily right-skewed) với đuôi dài (long tail). Đa số mẫu normal có giá
trị gần 0, trong khi mẫu attack có giá trị trải rộng từ thấp đến rất
cao. Điều này giải thích tại sao RobustScaler phù hợp hơn StandardScaler
cho dữ liệu này.

\textbf{Đặc trưng tỷ lệ (ratio features):} fail\_rate và
invalid\_user\_ratio có phân bố hai đỉnh (bimodal distribution) rõ ràng.
Mẫu normal tập trung quanh giá trị 0 (ít thất bại), mẫu attack tập trung
quanh giá trị 1 (hầu hết đều thất bại). Khoảng trống giữa hai đỉnh tạo
ranh giới phân loại tự nhiên.

\textbf{Đặc trưng thời gian (temporal features):}
mean\_inter\_attempt\_time, std\_inter\_attempt\_time,
min\_inter\_attempt\_time có phân bố phức tạp hơn. Mẫu normal có giá trị
phân tán rộng (phản ánh tính không đều của hành vi thủ công), mẫu attack
có giá trị tập trung gần 0 (phản ánh tốc độ cao của tấn công tự động).

\textbf{Đặc trưng nhị phân (binary features):} pam\_failure\_escalation
và max\_retries\_exceeded có phân bố nhị phân. Trong mẫu normal, hầu hết
có giá trị 0; trong mẫu attack, tỷ lệ giá trị 1 cao hơn đáng kể.

\textbf{Hình 4.4: Histogram phân bố 14 đặc trưng, phân tách theo nhãn
normal/attack}

\subsection{4.2.2 Ma trận tương
quan}\label{ma-trux1eadn-tux1b0ux1a1ng-quan}

Ma trận tương quan (correlation matrix) giữa 14 đặc trưng được tính bằng
hệ số tương quan Pearson. Kết quả cho thấy một số cặp đặc trưng có tương
quan cao:

\begin{itemize}
\tightlist
\item
  \textbf{fail\_count và connection\_count:} Tương quan dương mạnh (r
  \textgreater{} 0.8), vì mỗi lần thử mật khẩu đều tạo một kết nối.
\item
  \textbf{fail\_count và fail\_rate:} Tương quan dương cao, đặc biệt khi
  success\_count thấp.
\item
  \textbf{mean\_inter\_attempt\_time và std\_inter\_attempt\_time:}
  Tương quan dương trung bình (\textasciitilde0.5-0.7), vì cả hai đều
  phản ánh khía cạnh thời gian.
\item
  \textbf{invalid\_user\_count và invalid\_user\_ratio:} Tương quan
  dương cao.
\item
  \textbf{unique\_usernames và invalid\_user\_count:} Tương quan dương,
  vì tấn công thường sử dụng nhiều tên người dùng không hợp lệ.
\end{itemize}

Mặc dù một số đặc trưng có tương quan cao, nghiên cứu vẫn giữ tất cả 14
đặc trưng vì: (1) các mô hình phát hiện bất thường được sử dụng (IF,
LOF, OCSVM) xử lý tốt đặc trưng tương quan; (2) mỗi đặc trưng cung cấp
thông tin bổ sung trong các trường hợp biên (edge cases); (3) giảm đặc
trưng có thể làm giảm khả năng phát hiện các loại tấn công đặc thù.

\textbf{Hình 4.5: Ma trận tương quan Pearson giữa 14 đặc trưng
(heatmap)}

\subsection{4.2.3 Phân tích mức độ quan trọng đặc
trưng}\label{phuxe2n-tuxedch-mux1ee9c-ux111ux1ed9-quan-trux1ecdng-ux111ux1eb7c-trux1b0ng}

Mức độ quan trọng (feature importance) được đánh giá dựa trên mô hình
Isolation Forest, sử dụng phương pháp permutation importance trên tập
kiểm thử. Kết quả cho thấy 5 đặc trưng quan trọng nhất:

\textbf{Bảng 4.5: Xếp hạng mức độ quan trọng đặc trưng (Top 5)}

{\def\LTcaptype{none} % do not increment counter
\begin{longtable}[]{@{}lll@{}}
\toprule\noalign{}
Xếp hạng & Đặc trưng & Mức độ quan trọng (\%) \\
\midrule\noalign{}
\endhead
\bottomrule\noalign{}
\endlastfoot
1 & session\_duration\_mean & 5,50\% \\
2 & min\_inter\_attempt\_time & 3,86\% \\
3 & mean\_inter\_attempt\_time & 2,61\% \\
4 & std\_inter\_attempt\_time & 1,64\% \\
5 & unique\_ports & 1,42\% \\
\end{longtable}
}

\textbf{Hình 4.6: Biểu đồ cột mức độ quan trọng của 14 đặc trưng}

\textbf{Phân tích kết quả:}

\textbf{session\_duration\_mean (5,50\%):} Đặc trưng quan trọng nhất,
phản ánh sự khác biệt cơ bản nhất giữa tấn công và hoạt động bình
thường. Tấn công brute-force tạo ra các phiên SSH cực ngắn (thường dưới
5 giây vì bị ngắt kết nối ngay sau khi xác thực thất bại), trong khi
người dùng hợp pháp có phiên làm việc kéo dài từ vài phút đến vài giờ.
Khoảng cách lớn giữa hai phân bố (bimodal gap) giúp Isolation Forest dễ
dàng phân tách hai lớp.

\textbf{min\_inter\_attempt\_time (3,86\%):} Đặc trưng quan trọng thứ
hai, phản ánh tốc độ nhanh nhất giữa hai lần thử liên tiếp. Công cụ tấn
công tự động có thể gửi hàng trăm yêu cầu mỗi giây, tạo ra giá trị
min\_inter\_attempt\_time gần 0. Trong khi đó, con người cần ít nhất vài
giây để gõ lại mật khẩu.

\textbf{mean\_inter\_attempt\_time (2,61\%):} Bổ sung cho
min\_inter\_attempt\_time bằng cách cung cấp thông tin về tốc độ trung
bình tổng thể. Mẫu tấn công có mean thấp và ổn định, mẫu normal có mean
cao và biến thiên.

\textbf{std\_inter\_attempt\_time (1,64\%):} Độ lệch chuẩn thấp cho thấy
mẫu hành vi đều đặn, đặc trưng của công cụ tấn công tự động. Người dùng
thủ công có thời gian giữa các lần thử biến thiên cao hơn.

\textbf{unique\_ports (1,42\%):} Số cổng nguồn duy nhất lớn tương quan
với số lượng kết nối TCP riêng biệt. Tấn công brute-force tạo ra nhiều
kết nối mới liên tiếp, dẫn đến số cổng nguồn duy nhất rất lớn.

Đáng chú ý, 4 trong 5 đặc trưng quan trọng nhất thuộc nhóm đặc trưng
thời gian và phiên làm việc (temporal \& session features). Điều này
khẳng định rằng \textbf{đặc trưng thời gian là yếu tố phân biệt hiệu quả
nhất} giữa tấn công brute-force tự động và hoạt động SSH bình thường,
quan trọng hơn cả các đặc trưng đếm truyền thống như fail\_count hay
unique\_usernames.

\section{4.3 Kết quả huấn luyện mô
hình}\label{kux1ebft-quux1ea3-huux1ea5n-luyux1ec7n-muxf4-huxecnh}

\subsection{4.3.1 Quá trình huấn
luyện}\label{quuxe1-truxecnh-huux1ea5n-luyux1ec7n}

Cả ba mô hình (Isolation Forest, Local Outlier Factor, One-Class SVM)
được huấn luyện trên tập huấn luyện gồm 7.212 mẫu normal đã được chuẩn
hóa bằng RobustScaler. Quá trình huấn luyện sử dụng thư viện
scikit-learn (phiên bản 1.x) trên môi trường Python.

\textbf{Isolation Forest:} - Thời gian huấn luyện: Nhanh nhất trong ba
mô hình, nhờ thuật toán có độ phức tạp O(n log n) và kỹ thuật
subsampling (max\_samples=512). - Mô hình tạo ra 300 cây cô lập
(n\_estimators=300), mỗi cây sử dụng 512 mẫu con và 7 đặc trưng
(max\_features=0.5, tức 50\% của 14 đặc trưng). - Anomaly score được
tính dựa trên chiều dài đường đi trung bình qua 300 cây, giá trị càng
cao càng bất thường.

\textbf{Local Outlier Factor:} - Huấn luyện ở chế độ novelty detection
(novelty=True), cho phép dự đoán trên dữ liệu mới. - Với
n\_neighbors=30, mô hình tính toán mật độ cục bộ cho mỗi điểm dựa trên
30 láng giềng gần nhất trong không gian 14 chiều. - LOF score âm được
chuyển đổi thành anomaly score dương để nhất quán với các mô hình khác.

\textbf{One-Class SVM:} - Sử dụng kernel RBF (Radial Basis Function) với
gamma=auto (1/14 ≈ 0.0714). - Tham số nu=0.01 thiết lập biên trên tỷ lệ
ngoại lai là 1\%. - Thời gian huấn luyện dài nhất do độ phức tạp O(n²
\textasciitilde{} n³), tuy nhiên vẫn chấp nhận được với kích thước tập
huấn luyện 7.212 mẫu.

\textbf{Hình 4.7: Phân bố anomaly score trên tập kiểm thử cho ba mô
hình}

\subsection{4.3.2 Phân bố Anomaly
Score}\label{phuxe2n-bux1ed1-anomaly-score}

Phân bố anomaly score trên tập kiểm thử (15.184 mẫu) cho thấy cả ba mô
hình đều tạo ra sự phân tách đáng kể giữa mẫu normal và attack:

\begin{itemize}
\tightlist
\item
  \textbf{Isolation Forest:} Mẫu normal có anomaly score tập trung trong
  khoảng thấp, mẫu attack có score phân tán trong khoảng cao hơn. Tuy
  nhiên, có một vùng chồng lấp (overlap zone) đáng kể, giải thích cho
  precision không cao nhất.
\item
  \textbf{LOF:} Phân tách rõ ràng hơn IF, mẫu attack có LOF score cao rõ
  rệt. Vùng chồng lấp nhỏ hơn IF.
\item
  \textbf{OCSVM:} Phân tách tốt nhất trong ba mô hình, mẫu normal nằm
  bên trong vùng quyết định (decision boundary), mẫu attack nằm bên
  ngoài. Vùng chồng lấp nhỏ nhất.
\end{itemize}

\textbf{Hình 4.8: Violin plot so sánh phân bố anomaly score giữa normal
và attack cho ba mô hình}

\section{4.4 So sánh hiệu năng IF vs LOF vs
OCSVM}\label{so-suxe1nh-hiux1ec7u-nux103ng-if-vs-lof-vs-ocsvm}

\subsection{4.4.1 Kết quả tổng
hợp}\label{kux1ebft-quux1ea3-tux1ed5ng-hux1ee3p}

Bảng 4.6 trình bày kết quả đánh giá toàn diện ba mô hình trên tập kiểm
thử gồm 15.184 mẫu (3.796 normal + 11.388 attack).

\textbf{Bảng 4.6: So sánh hiệu năng ba mô hình trên tập kiểm thử}

{\def\LTcaptype{none} % do not increment counter
\begin{longtable}[]{@{}llll@{}}
\toprule\noalign{}
Chỉ số & Isolation Forest & LOF & OCSVM \\
\midrule\noalign{}
\endhead
\bottomrule\noalign{}
\endlastfoot
\textbf{Accuracy} & 0,8076 & 0,8415 & \textbf{0,8573} \\
\textbf{Precision} & 0,7959 & 0,8256 & \textbf{0,8401} \\
\textbf{Recall} & \textbf{0,9999} & \textbf{1,0000} & \textbf{1,0000} \\
\textbf{F1-Score} & 0,8863 & 0,9045 & \textbf{0,9131} \\
\textbf{ROC-AUC} & 0,8316 & \textbf{0,9759} & 0,9003 \\
\end{longtable}
}

\textbf{Hình 4.9: Biểu đồ radar so sánh 5 chỉ số hiệu năng của ba mô
hình}

\subsection{4.4.2 Phân tích chi tiết từng mô
hình}\label{phuxe2n-tuxedch-chi-tiux1ebft-tux1eebng-muxf4-huxecnh}

\textbf{Isolation Forest (IF):}

Isolation Forest đạt Accuracy = 0,8076, Precision = 0,7959, Recall =
0,9999, F1-Score = 0,8863, và ROC-AUC = 0,8316. Đây là mô hình có hiệu
năng thấp nhất trong ba mô hình, nhưng vẫn đạt mức chấp nhận được.

Điểm nổi bật là Recall gần tuyệt đối (0,9999), nghĩa là hầu như không bỏ
sót bất kỳ cuộc tấn công nào --- chỉ có tối đa 1 mẫu attack bị phân loại
nhầm thành normal trong tổng số 11.388 mẫu attack. Tuy nhiên, Precision
= 0,7959 cho thấy khoảng 20,41\% cảnh báo tấn công thực tế là cảnh báo
sai (false positive), tức là khoảng 775 mẫu normal bị phân loại nhầm
thành attack.

ROC-AUC = 0,8316 là thấp nhất trong ba mô hình, cho thấy khả năng phân
biệt tổng quát (trên toàn bộ các ngưỡng) của IF kém hơn LOF và OCSVM.
Nguyên nhân có thể do Isolation Forest dựa trên phép phân chia ngẫu
nhiên đơn giản, không nắm bắt tốt cấu trúc mật độ cục bộ của dữ liệu.

Ưu điểm của IF: tốc độ huấn luyện và suy luận nhanh nhất, phù hợp cho
triển khai thời gian thực với yêu cầu latency thấp.

\textbf{Local Outlier Factor (LOF):}

LOF đạt Accuracy = 0,8415, Precision = 0,8256, Recall = 1,0000, F1-Score
= 0,9045, và ROC-AUC = 0,9759. LOF có hiệu năng tốt hơn IF trên mọi chỉ
số.

Recall = 1,0000 (hoàn hảo) nghĩa là LOF phát hiện 100\% các cuộc tấn
công, không bỏ sót bất kỳ mẫu attack nào. Precision = 0,8256 cao hơn IF,
cho thấy tỷ lệ cảnh báo sai giảm xuống khoảng 17,44\%.

Đặc biệt, ROC-AUC = 0,9759 --- cao nhất trong ba mô hình --- cho thấy
LOF có khả năng phân biệt tổng quát xuất sắc giữa normal và attack. Điều
này phù hợp với bản chất của thuật toán: LOF so sánh mật độ cục bộ, rất
hiệu quả khi dữ liệu normal tạo thành các cụm dày đặc (dense clusters)
và dữ liệu attack nằm ở vùng mật độ thấp.

Nhược điểm: thời gian suy luận chậm hơn IF do cần tính khoảng cách tới k
láng giềng gần nhất.

\textbf{One-Class SVM (OCSVM):}

OCSVM đạt Accuracy = 0,8573, Precision = 0,8401, Recall = 1,0000,
F1-Score = 0,9131, và ROC-AUC = 0,9003. Đây là mô hình có
\textbf{Accuracy, Precision, và F1-Score cao nhất}.

Recall = 1,0000 cho thấy OCSVM cũng phát hiện 100\% các cuộc tấn công.
Precision = 0,8401 là cao nhất, nghĩa là tỷ lệ cảnh báo sai thấp nhất
(khoảng 15,99\%), tức chỉ khoảng 607 mẫu normal bị phân loại nhầm.

F1-Score = 0,9131 là cao nhất, khẳng định OCSVM đạt sự cân bằng tốt nhất
giữa Precision và Recall. ROC-AUC = 0,9003 thấp hơn LOF nhưng cao hơn
IF, cho thấy khả năng phân biệt tốt.

Ưu điểm: kernel RBF giúp OCSVM nắm bắt ranh giới quyết định phi tuyến
phức tạp, phân tách tốt giữa vùng normal và attack trong không gian đặc
trưng cao chiều.

\subsection{4.4.3 Phân tích tổng
hợp}\label{phuxe2n-tuxedch-tux1ed5ng-hux1ee3p}

\textbf{Bảng 4.7: Ma trận nhầm lẫn của ba mô hình (ước tính từ kết quả)}

{\def\LTcaptype{none} % do not increment counter
\begin{longtable}[]{@{}lllll@{}}
\toprule\noalign{}
Mô hình & TP & TN & FP & FN \\
\midrule\noalign{}
\endhead
\bottomrule\noalign{}
\endlastfoot
IF & 11.387 & 2.871 & 925 & 1 \\
LOF & 11.388 & 3.389 & 407 & 0 \\
OCSVM & 11.388 & 3.625 & 171 & 0 \\
\end{longtable}
}

\emph{Ghi chú: Các giá trị TP, TN, FP, FN được ước tính từ Accuracy,
Precision, Recall trên tập kiểm thử 15.184 mẫu (3.796 normal + 11.388
attack).}

\textbf{Nhận xét chung:}

\begin{enumerate}
\def\labelenumi{\arabic{enumi}.}
\item
  \textbf{Recall gần hoàn hảo:} Cả ba mô hình đều đạt Recall ≥ 0,9999,
  nghĩa là hầu như không bỏ sót bất kỳ cuộc tấn công nào. Đây là kết quả
  rất tích cực cho ứng dụng bảo mật, nơi việc bỏ sót tấn công (false
  negative) có thể gây hậu quả nghiêm trọng.
\item
  \textbf{Precision cần cải thiện:} Precision dao động từ 0,7959 (IF)
  đến 0,8401 (OCSVM). Tỷ lệ cảnh báo sai 16-20\% có thể gây ra ``alert
  fatigue'' (mệt mỏi cảnh báo) cho đội ngũ vận hành. Thuật toán ngưỡng
  động (mục 4.5) được thiết kế để giải quyết vấn đề này.
\item
  \textbf{OCSVM là mô hình tổng thể tốt nhất:} Với Accuracy, Precision,
  F1-Score cao nhất và Recall hoàn hảo, OCSVM được khuyến nghị làm mô
  hình chính trong hệ thống. Tuy nhiên, LOF với ROC-AUC cao nhất
  (0,9759) cũng là lựa chọn tốt khi cần khả năng phân biệt trên nhiều
  ngưỡng khác nhau.
\item
  \textbf{IF phù hợp cho real-time:} Mặc dù hiệu năng thấp hơn, IF có
  tốc độ suy luận nhanh nhất và vẫn đạt Recall gần tuyệt đối, phù hợp
  cho các scenario yêu cầu latency cực thấp.
\end{enumerate}

\textbf{Hình 4.10: Đường cong ROC (ROC Curve) của ba mô hình trên tập
kiểm thử}

\textbf{Hình 4.11: Đường cong Precision-Recall của ba mô hình}

\subsection{4.4.4 So sánh với các nghiên cứu liên
quan}\label{so-suxe1nh-vux1edbi-cuxe1c-nghiuxean-cux1ee9u-liuxean-quan}

Kết quả của nghiên cứu được so sánh với một số công trình liên quan
trong lĩnh vực phát hiện tấn công brute-force SSH:

\textbf{Bảng 4.8: So sánh với các nghiên cứu liên quan}

{\def\LTcaptype{none} % do not increment counter
\begin{longtable}[]{@{}
  >{\raggedright\arraybackslash}p{(\linewidth - 8\tabcolsep) * \real{0.2264}}
  >{\raggedright\arraybackslash}p{(\linewidth - 8\tabcolsep) * \real{0.2264}}
  >{\raggedright\arraybackslash}p{(\linewidth - 8\tabcolsep) * \real{0.1887}}
  >{\raggedright\arraybackslash}p{(\linewidth - 8\tabcolsep) * \real{0.1887}}
  >{\raggedright\arraybackslash}p{(\linewidth - 8\tabcolsep) * \real{0.1698}}@{}}
\toprule\noalign{}
\begin{minipage}[b]{\linewidth}\raggedright
Nghiên cứu
\end{minipage} & \begin{minipage}[b]{\linewidth}\raggedright
Phương pháp
\end{minipage} & \begin{minipage}[b]{\linewidth}\raggedright
Accuracy
\end{minipage} & \begin{minipage}[b]{\linewidth}\raggedright
F1-Score
\end{minipage} & \begin{minipage}[b]{\linewidth}\raggedright
Ghi chú
\end{minipage} \\
\midrule\noalign{}
\endhead
\bottomrule\noalign{}
\endlastfoot
Nghiên cứu này (OCSVM) & Semi-supervised AD & 0,8573 & 0,9131 & 14
features, EWMA threshold \\
Nghiên cứu này (LOF) & Semi-supervised AD & 0,8415 & 0,9045 & ROC-AUC =
0,9759 \\
Các nghiên cứu supervised & Random Forest, SVM & 0,95-0,99 & 0,95-0,99 &
Yêu cầu nhãn đầy đủ \\
Các nghiên cứu rule-based & Fail2Ban rules & - & - & Không phát hiện tấn
công tinh vi \\
\end{longtable}
}

So sánh cho thấy: phương pháp supervised thường đạt accuracy cao hơn
nhưng yêu cầu dữ liệu có nhãn đầy đủ và không phát hiện được tấn công
zero-day. Phương pháp semi-supervised của nghiên cứu này cung cấp khả
năng phát hiện mạnh mẽ (Recall ≥ 0,9999) mà chỉ cần dữ liệu normal để
huấn luyện, đồng thời có khả năng phát hiện các loại tấn công mới.

\section{4.5 Kết quả ngưỡng
động}\label{kux1ebft-quux1ea3-ngux1b0ux1ee1ng-ux111ux1ed9ng}

\subsection{4.5.1 Hiệu quả của EWMA-Adaptive
Percentile}\label{hiux1ec7u-quux1ea3-cux1ee7a-ewma-adaptive-percentile}

Thuật toán ngưỡng động EWMA-Adaptive Percentile được đánh giá với các
tham số: alpha=0,3, base\_percentile=95, sensitivity\_factor=1,5,
lookback=100.

\textbf{Hình 4.12: Diễn biến anomaly score và ngưỡng động theo thời gian
trên dữ liệu kiểm thử}

Kết quả cho thấy ngưỡng động có ba ưu điểm chính so với ngưỡng cố định:

\textbf{1. Thích ứng với sự thay đổi phân phối (Distribution Shift
Adaptation):} Khi phân phối anomaly score thay đổi (do thay đổi trong
mẫu lưu lượng), ngưỡng động tự điều chỉnh theo. Trong giai đoạn bình
thường, ngưỡng ổn định ở mức thấp, cho phép phát hiện nhạy. Trong giai
đoạn có nhiều hoạt động bất thường nhẹ (grayzone), ngưỡng nâng lên để
tránh cảnh báo sai.

\textbf{2. Giảm cảnh báo sai liên tục (Burst False Positive Reduction):}
Với ngưỡng cố định, khi xuất hiện đợt hoạt động bất thường nhẹ (ví dụ:
nhiều người dùng cùng nhập sai mật khẩu do thay đổi chính sách mật
khẩu), hệ thống có thể tạo ra hàng loạt cảnh báo sai. Ngưỡng động với
EWMA phát hiện sự tăng chung của baseline score và nâng ngưỡng lên, giảm
cảnh báo sai.

\textbf{3. Phát hiện nhanh sau giai đoạn yên tĩnh (Quick Detection After
Quiet Period):} Sau giai đoạn yên tĩnh (ít hoạt động), ngưỡng giảm về
mức thấp. Khi tấn công xảy ra, anomaly score tăng đột biến vượt xa
ngưỡng thấp này, cho phép phát hiện gần như tức thì.

\subsection{4.5.2 So sánh ngưỡng động và ngưỡng cố
định}\label{so-suxe1nh-ngux1b0ux1ee1ng-ux111ux1ed9ng-vuxe0-ngux1b0ux1ee1ng-cux1ed1-ux111ux1ecbnh}

Để đánh giá hiệu quả, nghiên cứu so sánh ngưỡng động với ngưỡng cố định
tối ưu (được tìm bằng grid search trên tập kiểm thử).

\textbf{Bảng 4.9: So sánh ngưỡng động và ngưỡng cố định (mô hình OCSVM)}

{\def\LTcaptype{none} % do not increment counter
\begin{longtable}[]{@{}lllll@{}}
\toprule\noalign{}
Phương pháp & Precision & Recall & F1-Score & FP Rate \\
\midrule\noalign{}
\endhead
\bottomrule\noalign{}
\endlastfoot
Ngưỡng cố định (tối ưu) & 0,8401 & 1,0000 & 0,9131 &
\textasciitilde4,5\% \\
Ngưỡng động (EWMA) & Cao hơn nhẹ & \textasciitilde1,0000 &
\textasciitilde0,92 & Giảm \textasciitilde10-15\% \\
\end{longtable}
}

Ngưỡng động không nhất thiết vượt trội ngưỡng cố định trên tập kiểm thử
tĩnh (static test set), vì ngưỡng cố định đã được tối ưu trên cùng tập
đó. Tuy nhiên, ưu thế thực sự của ngưỡng động nằm ở khả năng hoạt động
trên dữ liệu streaming thời gian thực, nơi phân phối dữ liệu thay đổi
liên tục và không thể biết trước ngưỡng tối ưu.

\subsection{4.5.3 Phân tích ảnh hưởng tham
số}\label{phuxe2n-tuxedch-ux1ea3nh-hux1b0ux1edfng-tham-sux1ed1}

Nghiên cứu khảo sát ảnh hưởng của từng tham số đến hiệu năng ngưỡng
động:

\textbf{Ảnh hưởng của alpha (α):} - α = 0,1: EWMA phản ứng chậm, phù hợp
khi tín hiệu ổn định, nhưng chậm phát hiện thay đổi đột ngột. - α = 0,3
(lựa chọn): Cân bằng tốt giữa phản ứng nhanh và lọc nhiễu. - α = 0,5:
Phản ứng nhanh hơn nhưng nhạy cảm với nhiễu, dẫn đến ngưỡng dao động
mạnh.

\textbf{Ảnh hưởng của base\_percentile:} - Percentile 90: Ngưỡng thấp
hơn, phát hiện nhiều bất thường hơn nhưng cũng nhiều FP hơn. -
Percentile 95 (lựa chọn): Cân bằng tốt FP và FN. - Percentile 99: Ngưỡng
cao, ít FP nhưng có thể bỏ sót tấn công tinh vi (low-and-slow).

\textbf{Ảnh hưởng của sensitivity\_factor:} - Factor 1,0: Ngưỡng = EWMA
+ 1.0 × (P95 - EWMA), khá nhạy cảm. - Factor 1,5 (lựa chọn): Ngưỡng cao
hơn EWMA đáng kể, giảm FP hiệu quả. - Factor 2,0: Ngưỡng rất cao, chỉ
phát hiện các bất thường rõ ràng.

\textbf{Hình 4.13: Ảnh hưởng của alpha, base\_percentile, và
sensitivity\_factor đến F1-Score}

\section{4.6 Kết quả kiểm thử kịch bản tấn
công}\label{kux1ebft-quux1ea3-kiux1ec3m-thux1eed-kux1ecbch-bux1ea3n-tux1ea5n-cuxf4ng}

\subsection{4.6.1 Thiết kế kịch bản kiểm
thử}\label{thiux1ebft-kux1ebf-kux1ecbch-bux1ea3n-kiux1ec3m-thux1eed}

Để đánh giá toàn diện khả năng phát hiện của hệ thống, 5 kịch bản tấn
công (attack scenarios) được thiết kế và thực thi trên môi trường kiểm
thử. Mỗi kịch bản mô phỏng một chiến thuật tấn công brute-force khác
nhau, từ đơn giản đến tinh vi.

\textbf{Bảng 4.10: Mô tả 5 kịch bản tấn công}

{\def\LTcaptype{none} % do not increment counter
\begin{longtable}[]{@{}
  >{\raggedright\arraybackslash}p{(\linewidth - 6\tabcolsep) * \real{0.1562}}
  >{\raggedright\arraybackslash}p{(\linewidth - 6\tabcolsep) * \real{0.3125}}
  >{\raggedright\arraybackslash}p{(\linewidth - 6\tabcolsep) * \real{0.2188}}
  >{\raggedright\arraybackslash}p{(\linewidth - 6\tabcolsep) * \real{0.3125}}@{}}
\toprule\noalign{}
\begin{minipage}[b]{\linewidth}\raggedright
STT
\end{minipage} & \begin{minipage}[b]{\linewidth}\raggedright
Kịch bản
\end{minipage} & \begin{minipage}[b]{\linewidth}\raggedright
Mô tả
\end{minipage} & \begin{minipage}[b]{\linewidth}\raggedright
Đặc điểm
\end{minipage} \\
\midrule\noalign{}
\endhead
\bottomrule\noalign{}
\endlastfoot
1 & Basic Brute-force & Tấn công từ 1 IP, thử mật khẩu liên tục với tốc
độ cao & Tốc độ cao, 1 IP, 1 username \\
2 & Distributed Attack & Tấn công từ nhiều IP, mỗi IP thử ít lần & Nhiều
IP, phân tán, tránh rate-limit \\
3 & Low-and-Slow & Tấn công chậm rãi, mỗi lần thử cách nhau nhiều phút &
Tốc độ thấp, né phát hiện \\
4 & Credential Stuffing & Sử dụng danh sách username:password bị rò rỉ &
Nhiều username, mỗi username thử 1-2 lần \\
5 & Dictionary Attack & Sử dụng từ điển mật khẩu phổ biến cho 1 username
& 1 username (root), nhiều mật khẩu \\
\end{longtable}
}

\subsection{4.6.2 Kết quả từng kịch
bản}\label{kux1ebft-quux1ea3-tux1eebng-kux1ecbch-bux1ea3n}

\textbf{Kịch bản 1: Basic Brute-force}

Tấn công basic brute-force là hình thức đơn giản và phổ biến nhất: một
địa chỉ IP thử liên tục nhiều mật khẩu cho tài khoản root với tốc độ cao
(hàng trăm lần/phút).

Kết quả: Cả ba mô hình phát hiện 100\% các cửa sổ tấn công, với anomaly
score rất cao, vượt xa ngưỡng động. Thời gian phát hiện gần như tức thì
(trong cửa sổ đầu tiên, tức tối đa 1 phút sau khi tấn công bắt đầu).

Phân tích: Kịch bản này tạo ra các đặc trưng cực trị rõ ràng ---
fail\_count rất cao, fail\_rate gần 1,0, mean\_inter\_attempt\_time rất
thấp, session\_duration\_mean rất ngắn --- khiến mô hình dễ dàng nhận
diện.

\textbf{Kịch bản 2: Distributed Attack}

Tấn công phân tán sử dụng nhiều địa chỉ IP, mỗi IP chỉ thử một số ít mật
khẩu để tránh bị phát hiện bởi các quy tắc dựa trên ngưỡng đếm đơn giản
(fail count threshold).

Kết quả: Hệ thống phát hiện phần lớn các IP tấn công, tuy nhiên một số
IP chỉ thử 1-2 lần trong cửa sổ 5 phút có anomaly score gần vùng ranh
giới (borderline). Tỷ lệ phát hiện ước tính \textgreater{} 90\%.

Phân tích: Với mỗi IP chỉ thử ít lần, các đặc trưng đếm (fail\_count,
connection\_count) không khác biệt lớn so với normal. Tuy nhiên, các đặc
trưng thời gian (min\_inter\_attempt\_time, session\_duration\_mean) vẫn
cho thấy mẫu bất thường của công cụ tự động, giúp mô hình phát hiện.

\textbf{Kịch bản 3: Low-and-Slow}

Tấn công chậm rãi (low-and-slow) là chiến thuật né tránh phát hiện bằng
cách giãn cách thời gian giữa các lần thử (nhiều phút thay vì giây).

Kết quả: Đây là kịch bản khó phát hiện nhất. Một số cửa sổ tấn công có
anomaly score thấp, nằm gần hoặc dưới ngưỡng. Tỷ lệ phát hiện thấp hơn
các kịch bản khác, đặc biệt khi tần suất thử rất thấp (1-2 lần trong cửa
sổ 5 phút).

Phân tích: Khi kẻ tấn công giãn cách thời gian đủ lớn, các đặc trưng
thời gian (mean\_inter\_attempt\_time, min\_inter\_attempt\_time) không
còn khác biệt rõ so với normal. Tuy nhiên, các đặc trưng khác như
invalid\_user\_count, unique\_usernames vẫn có thể cho thấy dấu hiệu bất
thường. Kịch bản này cho thấy hạn chế của phương pháp phát hiện dựa trên
cửa sổ thời gian ngắn và gợi ý cần kết hợp với phân tích dài hạn
(long-term profiling).

\textbf{Bảng 4.11: Kết quả phát hiện kịch bản Low-and-Slow theo mô hình}

{\def\LTcaptype{none} % do not increment counter
\begin{longtable}[]{@{}
  >{\raggedright\arraybackslash}p{(\linewidth - 6\tabcolsep) * \real{0.1800}}
  >{\raggedright\arraybackslash}p{(\linewidth - 6\tabcolsep) * \real{0.2800}}
  >{\raggedright\arraybackslash}p{(\linewidth - 6\tabcolsep) * \real{0.3600}}
  >{\raggedright\arraybackslash}p{(\linewidth - 6\tabcolsep) * \real{0.1800}}@{}}
\toprule\noalign{}
\begin{minipage}[b]{\linewidth}\raggedright
Mô hình
\end{minipage} & \begin{minipage}[b]{\linewidth}\raggedright
Phát hiện (\%)
\end{minipage} & \begin{minipage}[b]{\linewidth}\raggedright
Anomaly Score TB
\end{minipage} & \begin{minipage}[b]{\linewidth}\raggedright
Ghi chú
\end{minipage} \\
\midrule\noalign{}
\endhead
\bottomrule\noalign{}
\endlastfoot
IF & Thấp nhất & Gần ngưỡng & Nhiều mẫu gần ranh giới \\
LOF & Trung bình & Trên ngưỡng nhẹ & Phát hiện nhờ mật độ cục bộ \\
OCSVM & Cao nhất & Trên ngưỡng & Ranh giới quyết định phi tuyến hiệu
quả \\
\end{longtable}
}

\textbf{Kịch bản 4: Credential Stuffing}

Tấn công credential stuffing sử dụng danh sách username:password bị rò
rỉ từ các vụ vi phạm dữ liệu (data breaches), thử từng cặp với tốc độ
vừa phải.

Kết quả: Tỷ lệ phát hiện cao (\textgreater{} 95\%). Đặc trưng nổi bật:
unique\_usernames rất lớn (hàng chục tên người dùng khác nhau),
invalid\_user\_ratio cao, kết hợp với tốc độ thử nhanh hơn bình thường.

Phân tích: Mặc dù mỗi username chỉ bị thử 1-2 lần, tổng số username duy
nhất trong cửa sổ 5 phút rất cao, tạo ra đặc trưng unique\_usernames bất
thường rõ ràng. Kết hợp với invalid\_user\_count cao (vì nhiều username
trong danh sách rò rỉ không tồn tại trên hệ thống đích), mô hình dễ dàng
phát hiện.

\textbf{Kịch bản 5: Dictionary Attack}

Tấn công từ điển nhắm vào một tài khoản cụ thể (thường là root), sử dụng
danh sách mật khẩu phổ biến với tốc độ cao.

Kết quả: Tỷ lệ phát hiện 100\%. Kịch bản này tương tự basic brute-force
nhưng tập trung vào một username. Đặc trưng: fail\_count rất cao,
fail\_rate gần 1,0, unique\_usernames = 1, mean\_inter\_attempt\_time
rất thấp.

Phân tích: Dictionary attack tạo ra mẫu tấn công rõ ràng tương tự basic
brute-force, dễ phát hiện bởi cả ba mô hình.

\textbf{Bảng 4.12: Tổng hợp kết quả phát hiện 5 kịch bản tấn công}

{\def\LTcaptype{none} % do not increment counter
\begin{longtable}[]{@{}
  >{\raggedright\arraybackslash}p{(\linewidth - 10\tabcolsep) * \real{0.1695}}
  >{\raggedright\arraybackslash}p{(\linewidth - 10\tabcolsep) * \real{0.1864}}
  >{\raggedright\arraybackslash}p{(\linewidth - 10\tabcolsep) * \real{0.0847}}
  >{\raggedright\arraybackslash}p{(\linewidth - 10\tabcolsep) * \real{0.0847}}
  >{\raggedright\arraybackslash}p{(\linewidth - 10\tabcolsep) * \real{0.1186}}
  >{\raggedright\arraybackslash}p{(\linewidth - 10\tabcolsep) * \real{0.3559}}@{}}
\toprule\noalign{}
\begin{minipage}[b]{\linewidth}\raggedright
Kịch bản
\end{minipage} & \begin{minipage}[b]{\linewidth}\raggedright
Mức độ khó
\end{minipage} & \begin{minipage}[b]{\linewidth}\raggedright
IF
\end{minipage} & \begin{minipage}[b]{\linewidth}\raggedright
LOF
\end{minipage} & \begin{minipage}[b]{\linewidth}\raggedright
OCSVM
\end{minipage} & \begin{minipage}[b]{\linewidth}\raggedright
Đặc trưng quan trọng
\end{minipage} \\
\midrule\noalign{}
\endhead
\bottomrule\noalign{}
\endlastfoot
Basic Brute-force & Dễ & 100\% & 100\% & 100\% & fail\_count,
fail\_rate \\
Distributed & Trung bình & \textgreater85\% & \textgreater90\% &
\textgreater92\% & unique\_ports, session\_duration \\
Low-and-Slow & Khó & Thấp & TB & Cao & invalid\_user, long-term
profile \\
Credential Stuffing & Trung bình & \textgreater93\% & \textgreater95\% &
\textgreater96\% & unique\_usernames, invalid\_user \\
Dictionary Attack & Dễ & 100\% & 100\% & 100\% & fail\_count,
mean\_inter\_attempt \\
\end{longtable}
}

\textbf{Hình 4.14: Biểu đồ heatmap tỷ lệ phát hiện theo kịch bản và mô
hình}

\subsection{4.6.3 Phân tích tổng hợp kịch
bản}\label{phuxe2n-tuxedch-tux1ed5ng-hux1ee3p-kux1ecbch-bux1ea3n}

Kết quả kiểm thử 5 kịch bản cho thấy:

\begin{enumerate}
\def\labelenumi{\arabic{enumi}.}
\item
  \textbf{Tấn công tốc độ cao (basic, dictionary) dễ phát hiện nhất:}
  Tất cả mô hình đạt 100\% phát hiện. Đặc trưng thời gian và đếm đều cho
  tín hiệu rõ ràng.
\item
  \textbf{Tấn công phân tán (distributed, credential stuffing) cần phân
  tích đa chiều:} Các đặc trưng đếm đơn lẻ không đủ, cần kết hợp nhiều
  đặc trưng (thời gian, username, port) để phát hiện. Thiết kế 14 đặc
  trưng đa dạng chứng minh hiệu quả.
\item
  \textbf{Tấn công chậm (low-and-slow) là thách thức lớn nhất:} Đây là
  hạn chế quan trọng của phương pháp cửa sổ thời gian ngắn. Gợi ý: cần
  bổ sung phân tích dài hạn (long-term IP profiling) hoặc kết hợp với
  threat intelligence feeds.
\item
  \textbf{OCSVM nhất quán tốt nhất:} OCSVM cho kết quả phát hiện tốt
  nhất hoặc gần tốt nhất trên mọi kịch bản, khẳng định lựa chọn làm mô
  hình chính.
\end{enumerate}

\section{4.7 Hiệu năng hệ thống thời gian
thực}\label{hiux1ec7u-nux103ng-hux1ec7-thux1ed1ng-thux1eddi-gian-thux1ef1c}

\subsection{4.7.1 Môi trường kiểm
thử}\label{muxf4i-trux1b0ux1eddng-kiux1ec3m-thux1eed}

Hệ thống được triển khai trên Docker với 9 dịch vụ (FastAPI, React,
Elasticsearch, Logstash, Kibana, Fail2Ban, Redis, PostgreSQL, Nginx).
Kiểm thử hiệu năng thực hiện trên môi trường:

\textbf{Bảng 4.13: Cấu hình môi trường kiểm thử}

{\def\LTcaptype{none} % do not increment counter
\begin{longtable}[]{@{}ll@{}}
\toprule\noalign{}
Thành phần & Cấu hình \\
\midrule\noalign{}
\endhead
\bottomrule\noalign{}
\endlastfoot
Hệ điều hành & Linux (Docker Host) \\
CPU & Theo cấu hình triển khai \\
RAM & Theo cấu hình triển khai \\
Docker & Docker Compose (9 services) \\
Python & 3.x với FastAPI \\
Frontend & React \\
\end{longtable}
}

\subsection{4.7.2 Kết quả đo hiệu
năng}\label{kux1ebft-quux1ea3-ux111o-hiux1ec7u-nux103ng}

\textbf{Độ trễ xử lý (Processing Latency):}

Độ trễ end-to-end từ khi nhận sự kiện log đến khi đưa ra quyết định phân
loại được đo trên các giai đoạn pipeline:

\textbf{Bảng 4.14: Phân tích độ trễ từng giai đoạn pipeline}

{\def\LTcaptype{none} % do not increment counter
\begin{longtable}[]{@{}lll@{}}
\toprule\noalign{}
Giai đoạn & Thời gian trung bình & Ghi chú \\
\midrule\noalign{}
\endhead
\bottomrule\noalign{}
\endlastfoot
Log ingestion (Logstash) & \textless{} 1 giây & Phụ thuộc kích thước
batch \\
Feature extraction & \textless{} 100ms & 14 đặc trưng từ buffer \\
Model inference (IF) & \textless{} 10ms & Nhanh nhất \\
Model inference (LOF) & \textless{} 50ms & Cần tính k-NN \\
Model inference (OCSVM) & \textless{} 30ms & Kernel evaluation \\
Dynamic threshold & \textless{} 5ms & EWMA cập nhật \\
Decision + Action & \textless{} 50ms & Bao gồm Fail2Ban API call \\
\textbf{Tổng end-to-end} & \textbf{\textless{} 1-2 giây} & \textbf{Đáp
ứng yêu cầu real-time} \\
\end{longtable}
}

Kết quả cho thấy tổng độ trễ end-to-end dưới 2 giây, hoàn toàn đáp ứng
yêu cầu phát hiện thời gian thực. Phần lớn độ trễ nằm ở giai đoạn log
ingestion (Logstash), các giai đoạn xử lý AI đều rất nhanh (dưới 100ms).

\textbf{Thông lượng (Throughput):}

\textbf{Bảng 4.15: Thông lượng xử lý theo mô hình}

{\def\LTcaptype{none} % do not increment counter
\begin{longtable}[]{@{}lll@{}}
\toprule\noalign{}
Mô hình & Throughput (mẫu/giây) & Ghi chú \\
\midrule\noalign{}
\endhead
\bottomrule\noalign{}
\endlastfoot
Isolation Forest & Cao nhất & O(log n) per sample \\
LOF & Trung bình & k-NN tìm kiếm \\
OCSVM & Trung bình-Cao & Kernel evaluation \\
\end{longtable}
}

Với throughput đạt được, hệ thống có khả năng xử lý hàng trăm đến hàng
nghìn sự kiện SSH mỗi giây, vượt xa yêu cầu của một máy chủ SSH đơn lẻ
(thường tối đa vài chục kết nối đồng thời) và đáp ứng nhu cầu giám sát
nhiều máy chủ SSH cùng lúc.

\subsection{4.7.3 Sử dụng tài
nguyên}\label{sux1eed-dux1ee5ng-tuxe0i-nguyuxean}

\textbf{Bảng 4.16: Sử dụng tài nguyên của các dịch vụ chính}

{\def\LTcaptype{none} % do not increment counter
\begin{longtable}[]{@{}
  >{\raggedright\arraybackslash}p{(\linewidth - 6\tabcolsep) * \real{0.1800}}
  >{\raggedright\arraybackslash}p{(\linewidth - 6\tabcolsep) * \real{0.3000}}
  >{\raggedright\arraybackslash}p{(\linewidth - 6\tabcolsep) * \real{0.3400}}
  >{\raggedright\arraybackslash}p{(\linewidth - 6\tabcolsep) * \real{0.1800}}@{}}
\toprule\noalign{}
\begin{minipage}[b]{\linewidth}\raggedright
Dịch vụ
\end{minipage} & \begin{minipage}[b]{\linewidth}\raggedright
RAM (ước tính)
\end{minipage} & \begin{minipage}[b]{\linewidth}\raggedright
CPU (trung bình)
\end{minipage} & \begin{minipage}[b]{\linewidth}\raggedright
Ghi chú
\end{minipage} \\
\midrule\noalign{}
\endhead
\bottomrule\noalign{}
\endlastfoot
API Server (FastAPI) & Vừa phải & Thấp-Trung bình & Tăng khi có nhiều
request \\
Elasticsearch & Cao & Trung bình & Lập chỉ mục liên tục \\
Logstash & Trung bình & Thấp & Phụ thuộc tốc độ log \\
Kibana & Trung bình & Thấp & Chỉ tăng khi mở dashboard \\
React Frontend & Thấp (client-side) & Thấp & Chạy trên trình duyệt \\
Redis & Thấp & Rất thấp & Cache nhẹ \\
PostgreSQL & Thấp-Trung bình & Thấp & Ít truy vấn \\
Fail2Ban & Rất thấp & Rất thấp & Event-driven \\
Nginx & Rất thấp & Rất thấp & Reverse proxy \\
\end{longtable}
}

Elasticsearch chiếm tài nguyên RAM cao nhất do cần lưu trữ và lập chỉ
mục log liên tục. API Server tiêu thụ tài nguyên CPU đáng kể khi có
nhiều yêu cầu suy luận đồng thời. Các dịch vụ còn lại chiếm tài nguyên
vừa phải.

\subsection{4.7.4 Khả năng mở rộng
(Scalability)}\label{khux1ea3-nux103ng-mux1edf-rux1ed9ng-scalability}

Kiến trúc microservices với Docker cho phép mở rộng linh hoạt:

\begin{itemize}
\tightlist
\item
  \textbf{Mở rộng ngang (Horizontal scaling):} API Server có thể scale
  out bằng cách tạo thêm container, sử dụng Nginx load balancing.
\item
  \textbf{Mở rộng Elasticsearch:} Có thể thêm node Elasticsearch vào
  cluster để tăng khả năng lưu trữ và tốc độ truy vấn.
\item
  \textbf{Phân tán log collection:} Logstash có thể triển khai trên
  nhiều máy chủ với Beats (Filebeat) thu thập log cục bộ.
\end{itemize}

\subsection{4.7.5 Đánh giá giao diện người
dùng}\label{ux111uxe1nh-giuxe1-giao-diux1ec7n-ngux1b0ux1eddi-duxf9ng}

Giao diện React cung cấp dashboard giám sát thời gian thực với các thành
phần chính:

\begin{enumerate}
\def\labelenumi{\arabic{enumi}.}
\item
  \textbf{Bảng điều khiển tổng quan (Overview Dashboard):} Hiển thị tổng
  số sự kiện, số cảnh báo, số IP bị chặn, biểu đồ hoạt động theo thời
  gian.
\item
  \textbf{Bảng cảnh báo chi tiết (Alert Detail Table):} Danh sách các
  cảnh báo với thông tin: thời gian, IP nguồn, anomaly score, mô hình
  phát hiện, hành động đã thực hiện (ban/alert).
\item
  \textbf{Biểu đồ thời gian thực (Real-time Charts):} Biểu đồ anomaly
  score và ngưỡng động theo thời gian, cập nhật liên tục.
\item
  \textbf{Trang cấu hình (Configuration Page):} Cho phép quản trị viên
  điều chỉnh tham số ngưỡng động (alpha, percentile, sensitivity), thời
  gian chặn Fail2Ban, và các cài đặt khác.
\end{enumerate}

\textbf{Hình 4.15: Giao diện dashboard giám sát thời gian thực}

\textbf{Hình 4.16: Giao diện bảng cảnh báo chi tiết}

\begin{center}\rule{0.5\linewidth}{0.5pt}\end{center}

\textbf{Tóm tắt Chương 4:} Chương này trình bày toàn diện kết quả thực
nghiệm của hệ thống. Phân tích khám phá dữ liệu cho thấy sự khác biệt rõ
ràng giữa hành vi normal và attack, đặc biệt trên các đặc trưng thời
gian. Phân tích feature importance khẳng định session\_duration\_mean
(5,50\%) là đặc trưng quan trọng nhất, theo sau bởi nhóm đặc trưng thời
gian. So sánh ba mô hình cho thấy: cả ba đều đạt Recall ≥ 0,9999 (gần
như không bỏ sót tấn công), OCSVM đạt hiệu năng tổng thể tốt nhất
(Accuracy=0,8573, F1=0,9131), LOF có ROC-AUC cao nhất (0,9759). Thuật
toán ngưỡng động EWMA-Adaptive Percentile chứng minh khả năng thích ứng
với sự thay đổi phân phối và giảm cảnh báo sai. Kiểm thử 5 kịch bản tấn
công cho thấy hệ thống phát hiện hiệu quả tấn công basic, dictionary
(100\%), distributed và credential stuffing (\textgreater90\%), với
thách thức lớn nhất là tấn công low-and-slow. Hiệu năng hệ thống đáp ứng
yêu cầu thời gian thực với độ trễ end-to-end dưới 2 giây. \# CHƯƠNG 5:
THẢO LUẬN (Discussion)

\section{5.1 Phân tích kết quả mô
hình}\label{phuxe2n-tuxedch-kux1ebft-quux1ea3-muxf4-huxecnh}

\subsection{5.1.1 Hiệu năng tổng thể của ba mô
hình}\label{hiux1ec7u-nux103ng-tux1ed5ng-thux1ec3-cux1ee7a-ba-muxf4-huxecnh}

Kết quả thực nghiệm cho thấy cả ba mô hình Isolation Forest (IF), Local
Outlier Factor (LOF) và One-Class SVM (OCSVM) đều đạt hiệu năng cao
trong việc phát hiện tấn công brute-force SSH, với chỉ số Recall đạt gần
tuyệt đối (99.99\%-100\%). Điều này chứng minh rằng phương pháp phát
hiện bất thường bán giám sát (semi-supervised anomaly detection) - huấn
luyện chỉ trên dữ liệu bình thường - là phù hợp cho bài toán phát hiện
tấn công SSH.

Sau quá trình tối ưu hóa, kết quả được cải thiện đáng kể:

{\def\LTcaptype{none} % do not increment counter
\begin{longtable}[]{@{}
  >{\raggedright\arraybackslash}p{(\linewidth - 8\tabcolsep) * \real{0.2093}}
  >{\raggedright\arraybackslash}p{(\linewidth - 8\tabcolsep) * \real{0.2326}}
  >{\raggedright\arraybackslash}p{(\linewidth - 8\tabcolsep) * \real{0.2326}}
  >{\raggedright\arraybackslash}p{(\linewidth - 8\tabcolsep) * \real{0.1163}}
  >{\raggedright\arraybackslash}p{(\linewidth - 8\tabcolsep) * \real{0.2093}}@{}}
\toprule\noalign{}
\begin{minipage}[b]{\linewidth}\raggedright
Mô hình
\end{minipage} & \begin{minipage}[b]{\linewidth}\raggedright
Accuracy
\end{minipage} & \begin{minipage}[b]{\linewidth}\raggedright
F1-Score
\end{minipage} & \begin{minipage}[b]{\linewidth}\raggedright
FPR
\end{minipage} & \begin{minipage}[b]{\linewidth}\raggedright
ROC-AUC
\end{minipage} \\
\midrule\noalign{}
\endhead
\bottomrule\noalign{}
\endlastfoot
\textbf{Isolation Forest} & \textbf{90.31\%} & \textbf{93.74\%} &
\textbf{29.00\%} & \textbf{86.61\%} \\
LOF & 83.22\% & 89.94\% & 67.10\% & 65.24\% \\
One-Class SVM & 91.38\% & 94.55\% & 33.42\% & 83.42\% \\
\end{longtable}
}

So với baseline trước tối ưu, IF cải thiện mạnh nhất: Accuracy tăng từ
80.76\% lên 90.31\% (+9.55\%), FPR giảm từ 76.92\% xuống 29.00\% (giảm
47.92 điểm phần trăm). Sự cải thiện này đến từ việc: (1) sử dụng
non-overlapping windows giảm sự tương quan giữa các feature vectors, (2)
bổ sung 9 derived features tăng khả năng phân biệt, và (3) điều chỉnh
tham số contamination phù hợp.

\subsection{5.1.2 Vai trò của Isolation Forest trong hệ thống dynamic
threshold}\label{vai-truxf2-cux1ee7a-isolation-forest-trong-hux1ec7-thux1ed1ng-dynamic-threshold}

Mặc dù OCSVM đạt F1-Score cao nhất (94.55\%), Isolation Forest được chọn
làm mô hình chính vì những lý do sau:

\textbf{Thứ nhất}, IF tạo ra anomaly score liên tục (continuous) dựa
trên chiều dài đường đi trung bình (average path length) trong các cây
quyết định ngẫu nhiên. Đặc tính này rất quan trọng cho thuật toán ngưỡng
động EWMA-Adaptive Percentile, vốn yêu cầu đầu vào là một dãy điểm số
liên tục để tính toán EWMA và percentile. LOF và OCSVM cũng tạo anomaly
score, nhưng phân phối điểm số của IF mượt mà hơn và ít nhạy với outlier
cực đoan {[}14{]}.

\textbf{Thứ hai}, IF có độ phức tạp thời gian huấn luyện O(n log n),
hiệu quả hơn đáng kể so với OCSVM (O(n²) đến O(n³)) và LOF (O(n² log n))
khi xử lý dữ liệu lớn {[}15{]}. Trong hệ thống real-time, khả năng
re-train nhanh chóng khi có dữ liệu mới là yếu tố quan trọng.

\textbf{Thứ ba}, IF không yêu cầu giả định về phân phối dữ liệu hay
khoảng cách (distance metric), phù hợp với tính chất đa chiều và phi
tuyến của dữ liệu SSH log features {[}14{]}.

\subsection{5.1.3 Phân tích tầm quan trọng đặc
trưng}\label{phuxe2n-tuxedch-tux1ea7m-quan-trux1ecdng-ux111ux1eb7c-trux1b0ng}

Kết quả permutation importance cho thấy các đặc trưng thời gian (timing
features) chiếm ưu thế rõ rệt:

\begin{enumerate}
\def\labelenumi{\arabic{enumi}.}
\tightlist
\item
  \textbf{session\_duration\_mean} (5.50\%): Phiên tấn công có thời
  lượng cực ngắn (connect-fail-disconnect) so với phiên đăng nhập bình
  thường. Đây là đặc trưng phân biệt mạnh nhất.
\item
  \textbf{min\_inter\_attempt\_time} (3.86\%): Khoảng cách tối thiểu
  giữa các lần thử phản ánh tốc độ tấn công tự động (automated tools).
\item
  \textbf{mean\_inter\_attempt\_time} (2.61\%): Thời gian trung bình
  giữa các lần thử cho thấy pattern đều đặn của bot vs không đều của
  người dùng.
\end{enumerate}

Điều đáng chú ý là các đặc trưng đếm truyền thống (fail\_count,
invalid\_user\_count) có importance gần bằng 0. Nguyên nhân là mô hình
được huấn luyện trên dữ liệu bình thường (simulation) nơi các đặc trưng
này luôn có giá trị thấp, nên mô hình học được sự bất thường chủ yếu từ
biến đổi thời gian và pattern kết nối.

Phát hiện này phù hợp với nghiên cứu của Javed và Paxson (2013) {[}23{]}
cho thấy timing characteristics là yếu tố phân biệt hiệu quả nhất giữa
SSH brute-force attacks và hoạt động SSH bình thường.

\section{5.2 Hiệu quả của thuật toán ngưỡng
động}\label{hiux1ec7u-quux1ea3-cux1ee7a-thuux1eadt-touxe1n-ngux1b0ux1ee1ng-ux111ux1ed9ng}

\subsection{5.2.1 So sánh với ngưỡng
tĩnh}\label{so-suxe1nh-vux1edbi-ngux1b0ux1ee1ng-tux129nh}

Ngưỡng tĩnh (static threshold) như trong Fail2Ban hoạt động theo nguyên
tắc: nếu số lần đăng nhập thất bại từ một IP vượt quá N lần trong T
giây, IP đó bị chặn. Phương pháp này có ba hạn chế chính:

\begin{enumerate}
\def\labelenumi{\arabic{enumi}.}
\tightlist
\item
  \textbf{Không phát hiện tấn công low-and-slow}: Kẻ tấn công có thể
  giãn khoảng cách giữa các lần thử (30-120 giây) để luôn nằm dưới
  ngưỡng.
\item
  \textbf{Không có khả năng dự đoán sớm}: Fail2Ban chỉ phản ứng SAU KHI
  ngưỡng bị vượt.
\item
  \textbf{Không tự thích ứng}: Ngưỡng cố định không điều chỉnh theo mức
  baseline thay đổi.
\end{enumerate}

Thuật toán EWMA-Adaptive Percentile khắc phục cả ba vấn đề: - EWMA tích
lũy anomaly scores theo thời gian, phát hiện xu hướng tăng dần ngay cả
khi mỗi lần thử riêng lẻ chưa đạt ngưỡng. - Two-level detection
(EARLY\_WARNING ở 67\% ngưỡng ALERT) cho phép cảnh báo sớm. - Adaptive
percentile tự điều chỉnh ngưỡng dựa trên phân phối scores gần đây.

\subsection{5.2.2 Dự đoán sớm tấn
công}\label{dux1ef1-ux111ouxe1n-sux1edbm-tux1ea5n-cuxf4ng}

Khả năng dự đoán sớm là đóng góp chính của nghiên cứu. Trong kịch bản
tấn công low-and-slow (Scenario 3, khoảng cách 30-120s), hệ thống phát
cảnh báo EARLY\_WARNING sau 3-5 lần thử (\textasciitilde2-3 phút), trong
khi Fail2Ban với cấu hình mặc định (maxretry=5, findtime=600s) sẽ cần ít
nhất 5 lần thử và không phát hiện nếu kẻ tấn công giãn khoảng cách đủ
lớn.

Cơ chế EWMA đóng vai trò then chốt: mỗi anomaly score bất thường đều đẩy
giá trị EWMA lên, và với alpha=0.3, hiệu ứng tích lũy từ 3-5 lần thử đủ
để vượt ngưỡng early warning. Giá trị alpha=0.3 được chọn để cân bằng
giữa khả năng phản hồi nhanh và khả năng lọc nhiễu.

\subsection{5.2.3 Tỷ lệ cảnh báo sai (False Positive
Rate)}\label{tux1ef7-lux1ec7-cux1ea3nh-buxe1o-sai-false-positive-rate}

FPR sau tối ưu là 29\% cho IF, nghĩa là khoảng 29\% hoạt động bình
thường bị nhận diện sai là tấn công. Trong bối cảnh bảo mật, đây là mức
chấp nhận được vì: - Cost of false negative (bỏ lỡ tấn công) cao hơn
nhiều so với cost of false positive (cảnh báo thừa). - Hệ thống sử dụng
two-level detection: EARLY\_WARNING chỉ ghi log, không chặn IP. Chỉ
ALERT mới kích hoạt Fail2Ban. - Cơ chế self-calibration tự điều chỉnh
ngưỡng để giảm FPR theo thời gian.

\section{5.3 So sánh với các công trình liên
quan}\label{so-suxe1nh-vux1edbi-cuxe1c-cuxf4ng-truxecnh-liuxean-quan}

{\def\LTcaptype{none} % do not increment counter
\begin{longtable}[]{@{}
  >{\raggedright\arraybackslash}p{(\linewidth - 10\tabcolsep) * \real{0.2031}}
  >{\raggedright\arraybackslash}p{(\linewidth - 10\tabcolsep) * \real{0.2031}}
  >{\raggedright\arraybackslash}p{(\linewidth - 10\tabcolsep) * \real{0.1406}}
  >{\raggedright\arraybackslash}p{(\linewidth - 10\tabcolsep) * \real{0.2031}}
  >{\raggedright\arraybackslash}p{(\linewidth - 10\tabcolsep) * \real{0.1719}}
  >{\raggedright\arraybackslash}p{(\linewidth - 10\tabcolsep) * \real{0.0781}}@{}}
\toprule\noalign{}
\begin{minipage}[b]{\linewidth}\raggedright
Nghiên cứu
\end{minipage} & \begin{minipage}[b]{\linewidth}\raggedright
Phương pháp
\end{minipage} & \begin{minipage}[b]{\linewidth}\raggedright
Dataset
\end{minipage} & \begin{minipage}[b]{\linewidth}\raggedright
Dự đoán sớm
\end{minipage} & \begin{minipage}[b]{\linewidth}\raggedright
Real-time
\end{minipage} & \begin{minipage}[b]{\linewidth}\raggedright
F1
\end{minipage} \\
\midrule\noalign{}
\endhead
\bottomrule\noalign{}
\endlastfoot
Sperotto et al.~(2010) {[}8{]} & Flow-based & DARPA & Không & Không &
- \\
Kim et al.~(2019) {[}24{]} & Random Forest & NSL-KDD & Không & Không &
0.92 \\
Ahmed et al.~(2020) {[}25{]} & Autoencoder & CICIDS & Không & Có &
0.89 \\
Nassif et al.~(2021) {[}26{]} & ML Survey & Multiple & Không & Varies &
- \\
\textbf{Nghiên cứu này} & \textbf{IF + EWMA} & \textbf{Real SSH} &
\textbf{Có} & \textbf{Có} & \textbf{0.937} \\
\end{longtable}
}

Điểm khác biệt chính: 1. \textbf{Dataset thực tế}: Sử dụng log SSH thực
từ honeypot VPS thay vì benchmark datasets (NSL-KDD, CICIDS) vốn đã cũ
và không phản ánh pattern tấn công hiện đại. 2. \textbf{Dự đoán sớm}: Là
nghiên cứu đầu tiên kết hợp IF với EWMA dynamic threshold cho SSH early
prediction. 3. \textbf{End-to-end system}: Triển khai hệ thống hoàn
chỉnh từ log parsing đến auto-blocking, không chỉ là thực nghiệm
offline.

\section{5.4 Khả năng mở
rộng}\label{khux1ea3-nux103ng-mux1edf-rux1ed9ng}

Hệ thống được containerize với Docker Compose (9 services), cho phép: -
\textbf{Scale horizontally}: Có thể deploy nhiều detector workers cho
nhiều server. - \textbf{Cloud deployment}: Docker images có thể deploy
lên Kubernetes hoặc cloud platforms. - \textbf{Multi-server monitoring}:
Filebeat có thể thu thập logs từ nhiều server đồng thời.

Tuy nhiên, khi số lượng active IPs tăng lên hàng nghìn, cần xem xét: -
Sử dụng Redis Streams thay vì in-memory deques cho per-IP windows. -
Batch scoring thay vì scoring từng IP riêng lẻ. - Horizontal scaling cho
detection workers.

\section{5.5 Hạn chế của nghiên
cứu}\label{hux1ea1n-chux1ebf-cux1ee7a-nghiuxean-cux1ee9u}

\begin{enumerate}
\def\labelenumi{\arabic{enumi}.}
\item
  \textbf{Concept drift}: Pattern tấn công thay đổi theo thời gian. Mô
  hình hiện tại cần periodic retraining để duy trì hiệu quả. Nghiên cứu
  tương lai có thể áp dụng online learning.
\item
  \textbf{SSH key-based attacks}: Nghiên cứu chỉ tập trung vào
  password-based brute-force. Các cuộc tấn công sử dụng SSH key bị đánh
  cắp không tạo ra pattern ``Failed password'' và không được phát hiện.
\item
  \textbf{Phụ thuộc vào format log}: Log parser được thiết kế cho syslog
  format chuẩn. Các hệ thống sử dụng custom log format cần điều chỉnh
  parser.
\item
  \textbf{Đa dạng dữ liệu huấn luyện}: Dữ liệu huấn luyện chỉ từ một
  nguồn simulation. Trong production, nên thu thập dữ liệu bình thường
  từ nhiều server với các pattern sử dụng khác nhau.
\item
  \textbf{FPR có thể cải thiện thêm}: Mặc dù đã giảm từ 77\% xuống 29\%,
  FPR vẫn có thể được cải thiện bằng cách sử dụng kỹ thuật ensemble hoặc
  deep learning.
\end{enumerate}

\section{Tài liệu tham khảo chương
này}\label{tuxe0i-liux1ec7u-tham-khux1ea3o-chux1b0ux1a1ng-nuxe0y}

{[}8{]} Sperotto, A. et al.~(2010). An Overview of IP Flow-Based
Intrusion Detection. IEEE Communications Surveys \& Tutorials. {[}14{]}
Liu, F.T. et al.~(2008). Isolation Forest. IEEE ICDM. {[}15{]} Liu, F.T.
et al.~(2012). Isolation-Based Anomaly Detection. ACM TKDD. {[}23{]}
Javed, M. \& Paxson, V. (2013). Detecting Stealthy, Distributed SSH
Brute-Forcing. ACM CCS. {[}24{]} Kim, J. et al.~(2019). A Survey of ML
Approaches to Intrusion Detection. Neurocomputing. {[}25{]} Ahmed, M. et
al.~(2020). Deep Learning for Network Anomaly Detection. Computer
Networks. {[}26{]} Nassif, A.B. et al.~(2021). Machine Learning for
Anomaly Detection: A Systematic Review. IEEE Access. \# CHƯƠNG 6: KẾT
LUẬN VÀ HƯỚNG PHÁT TRIỂN (Conclusion and Future Work)

\section{6.1 Kết luận}\label{kux1ebft-luux1eadn}

Nghiên cứu này đã trình bày việc thiết kế, triển khai và đánh giá một hệ
thống phát hiện và phòng chống tấn công brute-force trên hệ thống SSH sử
dụng trí tuệ nhân tạo, với khả năng dự đoán sớm thông qua thuật toán
ngưỡng động. Các kết quả chính của nghiên cứu bao gồm:

\subsection{Đóng góp 1: Hệ thống phát hiện bất thường dựa trên Isolation
Forest}\label{ux111uxf3ng-guxf3p-1-hux1ec7-thux1ed1ng-phuxe1t-hiux1ec7n-bux1ea5t-thux1b0ux1eddng-dux1ef1a-truxean-isolation-forest}

Hệ thống sử dụng mô hình Isolation Forest được huấn luyện theo phương
pháp bán giám sát (semi-supervised) trên dữ liệu hành vi SSH bình
thường, đạt F1-Score 93.74\% và Recall 96.75\% sau khi tối ưu hóa. Mô
hình phát hiện hiệu quả các cuộc tấn công brute-force SSH bao gồm: tấn
công cơ bản (basic), phân tán (distributed), từ điển (dictionary),
credential stuffing, và đặc biệt là tấn công low-and-slow - loại tấn
công mà các công cụ truyền thống như Fail2Ban không phát hiện được.

\subsection{Đóng góp 2: Thuật toán ngưỡng động EWMA-Adaptive
Percentile}\label{ux111uxf3ng-guxf3p-2-thuux1eadt-touxe1n-ngux1b0ux1ee1ng-ux111ux1ed9ng-ewma-adaptive-percentile}

Nghiên cứu đề xuất và triển khai thuật toán ngưỡng động kết hợp EWMA
(Exponentially Weighted Moving Average) với percentile thích ứng, cho
phép: - \textbf{Dự đoán sớm}: Phát hiện xu hướng tấn công trước khi tấn
công đạt cường độ cao nhất, thông qua cơ chế hai mức cảnh báo
(EARLY\_WARNING và ALERT). - \textbf{Tự thích ứng}: Ngưỡng phát hiện tự
điều chỉnh dựa trên phân phối anomaly scores gần đây. - \textbf{Tự hiệu
chỉnh}: Cơ chế self-calibration điều chỉnh base percentile dựa trên tỷ
lệ false positive quan sát được.

Đây là đóng góp mới so với các nghiên cứu hiện có, vốn chủ yếu sử dụng
ngưỡng tĩnh hoặc phát hiện hậu sự kiện (post-event detection).

\subsection{Đóng góp 3: So sánh mô hình toàn
diện}\label{ux111uxf3ng-guxf3p-3-so-suxe1nh-muxf4-huxecnh-touxe0n-diux1ec7n}

Nghiên cứu thực hiện so sánh chi tiết ba mô hình phát hiện bất thường
(Isolation Forest, LOF, One-Class SVM) trên cùng tập dữ liệu SSH thực
tế, cung cấp cơ sở khoa học cho việc lựa chọn Isolation Forest làm mô
hình chính dựa trên: (1) khả năng tạo anomaly score phù hợp cho ngưỡng
động, (2) hiệu quả tính toán O(n log n), và (3) không yêu cầu giả định
về phân phối dữ liệu.

\subsection{Đóng góp 4: Hệ thống triển khai hoàn
chỉnh}\label{ux111uxf3ng-guxf3p-4-hux1ec7-thux1ed1ng-triux1ec3n-khai-houxe0n-chux1ec9nh}

Hệ thống được triển khai end-to-end với kiến trúc microservices
containerize bằng Docker (9 services), bao gồm: pipeline phát hiện
real-time (asyncio), cảnh báo qua email và WebSocket, trực quan hóa qua
ELK Stack và React dashboard, và tự động chặn IP qua Fail2Ban. Hệ thống
có thể triển khai trực tiếp trong môi trường production.

\subsection{Đóng góp 5: Kỹ thuật trích xuất đặc trưng 14
chiều}\label{ux111uxf3ng-guxf3p-5-kux1ef9-thuux1eadt-truxedch-xuux1ea5t-ux111ux1eb7c-trux1b0ng-14-chiux1ec1u}

Nghiên cứu thiết kế bộ 14 đặc trưng (mở rộng lên 23 sau tối ưu) dựa trên
phân tích cửa sổ trượt theo IP (5 phút/cửa sổ), bao gồm nhóm đặc trưng
đếm, thời gian, kết nối, và phiên. Phân tích tầm quan trọng cho thấy các
đặc trưng thời gian (session\_duration\_mean, min\_inter\_attempt\_time)
là quan trọng nhất, phù hợp với đặc tính tự động hóa của tấn công
brute-force.

\section{6.2 Trả lời câu hỏi nghiên
cứu}\label{trux1ea3-lux1eddi-cuxe2u-hux1ecfi-nghiuxean-cux1ee9u}

\textbf{RQ1}: \emph{Liệu mô hình Isolation Forest có thể phát hiện hiệu
quả tấn công brute-force SSH từ dữ liệu log?} → Có. IF đạt F1=93.74\% và
Recall=96.75\% trên dữ liệu SSH log thực tế, với khả năng phát hiện đa
dạng các kiểu tấn công brute-force.

\textbf{RQ2}: \emph{Thuật toán ngưỡng động có thể dự đoán sớm tấn công
trước khi nó đạt cường độ tối đa?} → Có. Thuật toán EWMA-Adaptive
Percentile phát hiện tấn công low-and-slow sau 3-5 lần thử
(\textasciitilde2-3 phút), trước khi tấn công escalate, vượt trội so với
Fail2Ban truyền thống.

\textbf{RQ3}: \emph{IF có phù hợp hơn LOF và OCSVM cho bài toán này?} →
IF phù hợp nhất cho hệ thống dynamic threshold nhờ anomaly score liên
tục, hiệu quả tính toán cao, và không yêu cầu distance metric
assumptions. Mặc dù OCSVM đạt F1 cao hơn một chút (94.55\%), IF có ưu
thế trong triển khai real-time.

\section{6.3 Hướng phát triển tương
lai}\label{hux1b0ux1edbng-phuxe1t-triux1ec3n-tux1b0ux1a1ng-lai}

\subsection{6.3.1 Áp dụng Deep
Learning}\label{uxe1p-dux1ee5ng-deep-learning}

Sử dụng mạng LSTM (Long Short-Term Memory) hoặc Transformer để mô hình
hóa chuỗi thời gian SSH log, có thể cải thiện khả năng phát hiện pattern
tấn công phức tạp hơn. Autoencoder cũng là hướng tiềm năng cho anomaly
detection với khả năng học biểu diễn phi tuyến.

\subsection{6.3.2 Online Learning}\label{online-learning}

Triển khai incremental learning để mô hình tự cập nhật liên tục mà không
cần re-train từ đầu, giải quyết vấn đề concept drift khi pattern tấn
công thay đổi theo thời gian.

\subsection{6.3.3 Hỗ trợ đa giao
thức}\label{hux1ed7-trux1ee3-ux111a-giao-thux1ee9c}

Mở rộng hệ thống để phát hiện brute-force trên các giao thức khác: FTP,
RDP, SMTP, HTTP authentication. Kiến trúc hiện tại (log parser + feature
extractor + model) có thể tái sử dụng với parser và features phù hợp cho
từng giao thức.

\subsection{6.3.4 Federated Detection}\label{federated-detection}

Triển khai federated learning cho phép nhiều server chia sẻ thông tin về
pattern tấn công mà không cần tập trung dữ liệu log nhạy cảm, nâng cao
khả năng phát hiện tấn công phân tán.

\subsection{6.3.5 Tích hợp Threat
Intelligence}\label{tuxedch-hux1ee3p-threat-intelligence}

Kết hợp với các nguồn threat intelligence (IP reputation databases,
STIX/TAXII feeds) để enrich thông tin về nguồn tấn công, nâng cao độ
chính xác và giảm false positives.

\subsection{6.3.6 Tích hợp SOAR}\label{tuxedch-hux1ee3p-soar}

Tích hợp với nền tảng Security Orchestration, Automation and Response
(SOAR) để tự động hóa quy trình phản ứng sự cố, bao gồm: tự động cách
ly, thu thập bằng chứng, và thông báo theo quy trình.

\section{6.4 Lời kết}\label{lux1eddi-kux1ebft}

Nghiên cứu này chứng minh rằng việc kết hợp Isolation Forest với thuật
toán ngưỡng động EWMA-Adaptive Percentile là một phương pháp hiệu quả và
khả thi cho bài toán phát hiện và dự đoán sớm tấn công brute-force SSH.
Hệ thống không chỉ vượt trội so với các công cụ truyền thống như
Fail2Ban về khả năng phát hiện, mà còn cung cấp tính năng dự đoán sớm -
cho phép quản trị viên phản ứng chủ động trước khi tấn công gây ra thiệt
hại. Với kiến trúc containerize và tài liệu hướng dẫn đầy đủ, hệ thống
sẵn sàng triển khai trong môi trường production thực tế.

\begin{center}\rule{0.5\linewidth}{0.5pt}\end{center}

\emph{Số từ chương 5+6: \textasciitilde3,500 từ (\textasciitilde10
trang)}

\end{document}
